\documentclass[12pt]{book}
\usepackage{amsthm}
\usepackage{libertine}
\usepackage[utf8]{inputenc}
\usepackage[margin=1in]{geometry}
\usepackage{amsmath,amssymb}
\usepackage{multicol}
\usepackage[shortlabels]{enumitem}
\usepackage{siunitx}
\usepackage{cancel}
\usepackage{graphicx}
\usepackage{float}
\usepackage[dvipsnames]{xcolor}
\usepackage{pgfplots}
\usepackage{listings}
\usepackage{tikz}
\usepackage{hyperref}
\usepackage{cleveref}
\usepackage{fancyhdr}
\usepackage{setspace}
\usepackage{refcheck}
\usepackage{refcount}
\usepackage{soul}
\usepackage{subfiles}

\allowdisplaybreaks

\newcommand{\R}{\mathbb{R}}
\newcommand{\Rd}{\mathbb{R}^d}
\newcommand{\Q}{\mathbb{Q}}
\newcommand{\N}{\mathbb{N}}
\newtheorem{thm}{Theorem}[section]
\newtheorem*{ax}{Axiom}
\newtheorem{cor}{Corollary}[section]
\newtheorem{prop}{Proposition}[section]
\newtheorem{lemma}{Lemma}[section]
\theoremstyle{remark}
\newtheorem{rem}{Remark}[section]
\theoremstyle{definition}
\newtheorem{eg}{Example}[section]
\newtheorem{ex}{Exercise}[section]
\newtheorem{defn}{Definition}[section]
\newtheorem{q}{Question}[section]

\newenvironment{solution}{\begin{proof}[Solution]}{\end{proof}}
\renewcommand{\theequation}{\thesection.\arabic{equation}}

\title{Math 2180 -- Real Analysis 1 notes}
\author{Jeremy Wu}
\date{\today}
\doublespacing

\begin{document}
	\maketitle
	
	\textcolor{blue}{Introduction to metric spaces including connectedness, compactness and continuity; topics in infinite series of numbers, and sequences and series of functions. - U of M}
	To do:
	\begin{itemize}
		\item Write down product of complete / compact is complete / compact.
		\item Function spaces -- completeness of $C(X)$.
	\end{itemize}
	
	\tableofcontents
	
	\pagestyle{fancy}
	\lhead{Math 2180 Lecture notes}
	\section{Introduction and motivation}
	\label{sec:intro}
	These lecture notes cover the Math 2180 -- Real Analysis 1 course at the University of Manitoba. In particular, they will hopefully introduce some of the basic concepts of Analysis. The first part of this course covering abstract metric spaces and topology will follow selected sections from Wilson A. Sutherland's `Introduction to Metric and Topological Spaces.'
	
	My personal bias leans towards differential equations, thus these notes emphasise completeness (\Cref{sec:Cauchy}) and compactness (\Cref{sec:cpct_met_spaces}) \textcolor{blue}{and function spaces}. Connectedness (\Cref{sec:connected}) is relatively underappreciated here (my apologies to algebraic topology).
	
	To give a brief motivation of what these notes cover and why, I will begin by assuming that you, the reader, have already taken a course equivalent to Math 2080 -- Introduction to Analysis. There you learned (among other things) the rigorous meaning of the following statements:
	\begin{enumerate}
		\item A sequence of \textit{real numbers} $(a_n)_{n\in\N}\subset \R$ converges to $a\in\R$ as $n\to \infty$.
		\item A series of, again, \textit{real numbers} $\sum_{j=1}^na_j$ converges to $\sum_{j=1}^\infty a_j\in\R$.
	\end{enumerate}
	In the huge universe of the mathematical sciences, we don't always work with just real numbers. We deal with vectors, matrices, functions, \dots This course is meant to generalise the above items to more than just real numbers. For example, from your probability / statistics courses, you may have encountered (a version of) the \textbf{law of large numbers}. It says the following: If $X_1,X_2,\dots$ are independently and identically distributed random variables with $\mathbb{E}[X_1] = \mathbb{E}[X_2]=\dots = \mu$, then the sample average given by
	\[
	\bar{X}_n = \frac{1}{n}\sum_{j=1}^nX_j
	\]
	`converges' as $n\to \infty$ to the mean: $\bar{X}_n \to \mu$ as $n\to \infty$. Here, I have written `converges' with quotation marks because this notion of convergence is \textit{not} the same as the notion of convergence from Math 2080. Moreover, depending on which assumptions are made, the meaning of convergence in the above statement is different:
	\begin{itemize}
		\item For the \href{https://en.wikipedia.org/wiki/Law_of_large_numbers#Weak_law}{\textit{weak} law of large numbers}, the sample mean converges \textit{in probability}.
		\item For the \href{https://en.wikipedia.org/wiki/Law_of_large_numbers#Strong_law}{\textit{strong} law of large numbers}, the sample mean converges \textit{almost surely}.
	\end{itemize}
	The point that I want to emphasise here is that the above example with the law of large numbers uses different notions of convergence for sequences of \textit{random variables}. The \href{https://en.wikipedia.org/wiki/Law_of_large_numbers#Differences_between_the_weak_law_and_the_strong_law}{Wikipedia page} for the law of large numbers explains the difference between the weak and strong versions.
	
	Concerning the convergence of series, you have definitely studied \textbf{Taylor series} and you may have learned something about \textbf{Fourier series}. Both of these ideas have the same framework: You have a function $f(x)$ and the goal is to represent this function by a sequence of functions $f_n(x)$ in the following way.
	\begin{equation}
		\label{eq:intro_series_function}
	f(x) = \sum_{n=0}^\infty f_n(x).
	\end{equation}
	For example, if $f:[-\pi,\pi]\to \R$, then
	\begin{itemize}
		\item the Maclaurin series (Taylor series around the origin) representation takes the form
		\[
		f_n(x) = \frac{f^{(n)}(0)}{n!}x^n.
		\]
		\item The Fourier series representation takes the form
		\[
		f_n(x) = a_n \cos (nx) + b_n \sin (nx)
		\]
		where
		\[
		a_n = \frac{1}{\pi}\int_{-\pi}^\pi f(x) \cos nx \,\mathrm{d}x\text{ and }b_n = \frac{1}{\pi}\int_{-\pi}^\pi f(x) \sin nx\,\mathrm{d}x.
		\]
	\end{itemize}
	Returning to~\eqref{eq:intro_series_function}, a fundamental question is to understand what is meant by
	\[
	f(x) = \lim_{n\to \infty}\sum_{j=0}^nf_j(x)?
	\]
	One of the convenient (but not completely obvious nor rigorous!) things about~\eqref{eq:intro_series_function} is that we may be tempted to perform a computation like
	\[
	\frac{df}{dx} = \frac{d}{dx}\left(\sum_{n=0}^\infty f_n(x)\right) \overset{!}{=}\sum_{n=0}^\infty \frac{df_n}{dx}.
	\]
	This is very handy but is not always true! The possible failure or legitimacy of the above operation is sensitive to the notion of \textit{convergence} of the series~\eqref{eq:intro_series_function}.
	
	As we will see in this course, it is important to understand the precise notion of convergence. To recap: in Math 2080, you learned what it meant for a sequence / series of \textit{real numbers} to converge. This course is meant to generalise, from those principles, what it means for a sequence / series of \textit{mathematical objects} (real numbers, vectors, matrices, functions, random variables, \dots) to converge. In analogy to the previous two points summarising Math 2080, we aim to answer the following questions in this course.
	\begin{enumerate}
		\item Let $(a_n)_{n\in\N}$ and $a$ be mathematical objects belonging to a set $X$ (of real numbers, or of functions, or \dots). What does it mean for $a_n \to a$ as $n\to \infty$ and what is the role of $X$ in this definition?
		\item Let $(f_n)_{n\in\N}$ and $f$ be real-valued functions (but their inputs may belong in some abstract set $X$). What does it mean for $\sum_{j=1}^nf_j\to f$ as $n\to \infty$?
	\end{enumerate}
	\begin{ex}[Failure of convergence of Taylor Series]
	Consider the function
	\[
	f(x) = \left\{
	\begin{array}{cl}
		0,	&\text{if }x\le 0 	\\
		e^{-\frac{1}{x}},	&\text{if }x>0
	\end{array}
	\right..
	\]
	You may use the following result (it is also a good exercise to prove it!):
	\begin{equation}
		\label{eq:intro_seq_exp}
		\lim_{x\to \infty} x^p e^{-x} = 0,\quad \forall p\ge 0.
	\end{equation}
	This exercise concerns the Taylor series of $f$ around the point $x=0$.
	\begin{enumerate}
		\item Sketch $f$.
		\item Prove that $f$ is continuous at $x=0$.
		\item Let $n\in\N$ and remember that we use the symbol $f^{(n)}(a)$ to denote the value of the $n$th derivative of $f$ evaluated at $a$. Calculate $f^{(n)}(0)$. \textit{Hint: }They are all the same value.
		\item Let $n\in\N$ and, using the previous part, write down the $n$th order Taylor series of $f$.
		\item Let $x>0$. Based on the previous part, is it true or false that $f(x)$ equals its infinite Taylor series expansion around 0? In other words, is the following formula valid?
		\[
		f(x) \overset{?}{=} \sum_{n=0}^\infty \frac{f^{(n)}(0)}{n!}x^n.
		\]
	\end{enumerate}
	\end{ex}
%	\section{Comments (to an instructor)}
%	As mentioned in~\Cref{sec:intro}, these notes largely follow Terence Tao's `Analysis II' but I should clarify one large distinction especially pertaining to chapters 1 and 2 of his book. In his book, all notions of topology are made in reference to some metric and this is how he introduced many definitions. Where possible, I have elected to introduce definitions with as little reliance on metric-exclusive ideas as possible. That being said, these notes do not cover topological spaces in general, only metric topology.
	\section{Some notations}
	I will use various abbreviations and symbols which you have hopefully seen from earlier foundational (pure) maths courses.
	\[
	\begin{array}{c|c}
	\text{Symbol} 	&\text{Meaning} 	\\\hline
	\text{s.t.} 	&\text{such that} 	\\
	\text{iff} 		&\text{if and only if, }\iff 	\\
	\text{TFAE} 	&\text{the following are equivalent} 	\\
	\text{wlog} 	&\text{without loss of generality} 	\\
	\forall 		&\text{for every, for all, for any, \dots} 	\\
	\exists 		&\text{there exists, there is, \dots} 	\\
	\subset 	&\text{(not necessarily strict )subset of\dots }
	\end{array}
	\]
	Most, if not all, results here about a metric (or topological) space $X$ implicitly assume $X\neq \emptyset$. I have explicitly and redundantly stated this non-emptiness as an assumption in a lot of results only to emphasise its importance in those results. For simplicity, the reader may just assume that all metric spaces are non-empty unless otherwise stated.
	\section{Reader's guide}
	\label{sec:guide}
	These notes are written in a different order than how I would, personally, deliver a lecture course. In particular, these notes cover (strictly) more than I would have time to lecture on in a 1-semester / term / quarter course.
	
	\Cref{ch:metric_spaces} introduces the concepts of metric spaces. As mentioned in~\Cref{sec:intro}, completeness, compactness, and connectedness  (\Cref{sec:Cauchy}, \Cref{sec:cpct_met_spaces}, and~\Cref{sec:connected} respectively) with an emphasis on the first two due to my own bias. The length of~\Cref{sec:cpct_met_spaces} reflects my attempt at proving results with \textit{minimal reliance on metrics}. This is meant to pre-emptively illustrate the flavour of \textit{topological spaces}, the generalization of metric spaces. Similarly, \Cref{sec:connected} relies entirely on topological arguments, rather than explicit metric structures. Finally, this chapter focuses on metric spaces in isolation \textbf{without relation to continuous functions yet.}
	
	\Cref{ch:cts_met} introduces and discusses continuous functions on metric spaces. This chapter also investigates the interplay between continuous functions and complete / compact / connected spaces. In a lecture course, I would (probably) discuss this interplay at the same time as I introduce the notions of complete / compact / connectedness to some extent. In other words, my lectures would likely \textbf{jump back and forth between the first two chapters.}
	\chapter{Metric Spaces}
	\label{ch:metric_spaces}
	This chapter is dedicated to an introduction to the theory of metric spaces. For all students, I strongly suggest reading Wilson A. Sutherland's `Introduction to Metric and Topological Spaces' which is the basis for this.
	\section{Definitions and examples}
\label{sec:def_met}
Recall from Math 2080 the following definition for the convergence of a sequence of real numbers.
\begin{defn}[Convergence - sequence of real numbers]
	\label{defn:cvg_real}
	Let $(a_n)_{n\in\N}\subset \R$ and $a\in \R$. We say that $(a_n)_{n\in\N}$ \textit{converges to }$a$, written $a_n\overset{n\to \infty}{\to}a$, if and only if the following statement is true:
	\[
	\forall \varepsilon>0,\exists N\in\N\text{ such that }\forall n\ge N, \, |a_n - a|<\varepsilon.
	\]
\end{defn}
Before you took Math 2080, you may have informally described convergence as something like `the points $a_n$ \textit{get closer} to $a$ as $n$ gets larger.' Within~\Cref{defn:cvg_real}, the rigorous notion of `closeness' is captured precisely by
\[
|a_n-a| = \text{`distance between }a_n\text{ and }a\text{'}
\]
We want to generalise~\Cref{defn:cvg_real} when $a_n$ and $a$ are not necessarily real numbers. In order to do so, we need to be able to estimate the \textit{distance} between $a_n$ and $a$, whatever kind of mathematical object they may be. This motivates our first definition.
\begin{defn}[Metric spaces]
	\label{defn:metric}
	Let $X$ be a non-empty set and $d:X\times X \to [0,+\infty)$ a function. We say that the pair $(X,d)$ is a \textit{metric space} provided the function $d$ satisfies the following:
	\begin{enumerate}[label = \textbf{M\arabic*)}]
		\item (Non-negativity) For any $x,y\in X$, we have $d(x,y)\ge 0$. Moreover, $d(x,y) = 0 \iff x=y$. \label{m:pos}
		\item (Symmetry) For any $x,y\in X$, we have $d(x,y) = d(y,x)$. \label{m:sym}
		\item (Triangle inequality) For any $x,y,z\in X$, we have
		\[
		d(x,z) \le d(x,y) + d(y,z).
		\]
		\label{m:tri}
	\end{enumerate}
	In this case, we refer to the function $d$ as the \textit{metric} or \textit{distance}. We will also call the elements of $X$ \textit{points}.
\end{defn}
Sometimes we refer to $X$ as a metric space if the metric is clear or given from context. A set $X$ may or may not have many distinct metrics associated to it. The motivation for studying metric spaces is that distance function allows us to quantify how far two elements are from each other. Before listing (many) examples of metric spaces, we will generalise~\Cref{defn:cvg_real} to metric spaces.
\begin{defn}[Convergence in a metric space]
	\label{defn:cvg_metric}
Let $(X,d)$ be a metric space be given together with a sequence $(a_n)_{n\in\N}\subset X$ and a point $a\in X$. We say that $(a_n)_{n\in\N}$ \textit{converges to }$a$ \textit{(with respect to the metric $d$)}, written $a_n\overset{n\to \infty}{\to} a$, if and only if the following statement is true:
\[
\forall \varepsilon>0,\exists N\in\N\text{ such that }\forall n\ge N, \, d(a_n,a)<\varepsilon.
\]
\end{defn}
Let us start with the most familiar example of a metric space.
\begin{eg}[The reals]
\label{eg:real_metric}
Let $d:\R\times\R\to [0,+\infty)$ be defined by $d(x,y) := |x-y|$. Then, the pair $(\mathbb{R}, d)$ is a metric space. Thus, \Cref{defn:cvg_metric} is equivalent to~\Cref{defn:cvg_real} in this case.
\end{eg}
\begin{eg}[Euclidean spaces]
Let $n\in\N$ be a natural number and denote $\R^n$ the vector space of $n$-tuples of real numbers:
\[
\R^n := \left\{
(x_1,x_2,\dots, x_n)\, : \, x_i\in\R \, \forall i=1,\dots, n
\right\}.
\]
We define the \textit{Euclidean metric} (also called the $\ell^2$ metric) $d_2:\R^n\times \R^n\to [0,+\infty)$ by
\begin{align*}
d_2((x_1,x_2,\dots, x_n),\,(y_1,y_2,\dots, y_n))&:= \sqrt{(x_1-y_1)^2 + (x_2-y^2)^2 + \dots + (x_n-y_n)^2} 	\\
&= \left(\sum_{i=1}^n (x_i-y_i)^2\right)^\frac{1}{2}.
\end{align*}
The pair $(\R^n,\, d_2)$ is a metric space. Usually, when we work in $\R^n$, we will refer to $d_2$ as the \textit{standard metric} (on $\R^n$).
\end{eg}
\begin{eg}[$\ell^1$ and $\ell^\infty$ metrics on $\R^n$]
Again, let $n\in\N$ be a natural number and consider $\R^n$ the vector space of $n$-tuples of real numbers. Let $x,y\in\R^n$ be vectors with components
\[
x=(x_1,x_2,\dots, x_n)\text{ and }y = (y_1,y_2,\dots, y_n).
\]
\begin{enumerate}
	\item The \textit{taxi-cab metric} (also called the $\ell^1$ metric) is defined by
	\begin{align*}
	d_1(x,y) 	&:= |x_1-y_1| + |x_2-y_2| + \dots + |x_n - y_n| = \sum_{i=1}^n|x_i-y_i|.
	\end{align*}
	For example, if $n=2$, then $d_1((1,6),(4,2)) = |1 - 4| + |6 - 2| = 3 + 4 = 7$. Intuitively, the $\ell^1$ metric models the distance it takes for a taxi-cab to get from one point to another by only moving up / right / down / left (North / East / South / West) since most North American cities have streets built on a grid system which prevents diagonal movement.
	\item The \textit{sup norm metric} (also called the $\ell^\infty$ metric) is defined by
	\begin{align*}
		d_\infty(x,y) 	&:= \sup_{i=1,\dots, n}|x_i-y_i|
	\end{align*}
	For example, if $n=2$, then $d_\infty((1,6),(4,2)) = \sup (3,4) = 4$.
\end{enumerate}
\end{eg}
Notice from the previous example in $\R^2$ that we had
\[
d_1((1,6),(4,2)) = 7,\quad d_\infty((1,6),(4,2)) = 4,\quad d_2((1,6),(4,2)) = 5.
\]
So all three metrics $d_1,d_\infty, d_2$ measure the distance between the points $(1,6)$ and $(4,2)$ on the plane in different ways. However, they are all \textit{equivalent} in the following sense (which we shall hopefully revisit).
\begin{prop}
[Equivalence of $d_1,d_2,d_\infty$ on $\R^n$]
\label{prop:equiv_d12inf}
Fix $n\in\N$ and let $(a_j)_{j\in\N}\subset \R^n$ and $a\in \R^n$. The following are equivalent:
\begin{itemize}
	\item $a_j\overset{j\to \infty}{\to}a$ with respect to $d_1$.
	\item $a_j\overset{j\to \infty}{\to}a$ with respect to $d_2$.
	\item $a_j\overset{j\to \infty}{\to}a$ with respect to $d_\infty$.
\end{itemize}
\end{prop}
In words, \Cref{prop:equiv_d12inf} is saying that if $a_j$ converges to $a$ in any one of the metrics $d_1 / d_2 / d_\infty$, then $a_j$ converges to $a$ in \textit{all} three metrics.
\begin{proof}
We will take the following inequalities for granted (see the first problem sheet for a proof of~\eqref{eq:hold_1} and~\eqref{eq:hold_2}).
\begin{equation}
	\label{eq:hold_1}
	d_2(x,y) \le d_1(x,y)\le \sqrt{n} d_2(x,y),\quad \forall x,y\in\R^n,
\end{equation}
and
\begin{equation}
	\label{eq:hold_2}
	\frac{1}{\sqrt{n}}d_2(x,y)\le d_\infty(x,y)\le d_2(x,y),\quad\forall x,y\in\R^n.
\end{equation}
We will just prove that $a_j \to a$ in $d_1$ if and only if $a_j\to a$ in $d_2$.

Let $a_j\overset{j\to \infty}{\to}a$ in $d_1$. Now, \textcolor{blue}{fix $\varepsilon>0$.} By assumption, we know that \textcolor{blue}{there exists $N\in\N$ such that $j\ge N$} implies
\[
d_1(a_j,a)<\varepsilon.
\]
But the left inequality of~\eqref{eq:hold_1} says $d_2(x,y)\le d_1(x,y)$ and so when $j\ge N$, we have
\[
{\color{blue}
d_2(a_j,a)<\varepsilon.
}
\]
The text in \textcolor{blue}{blue} shows that $a_j \overset{j\to \infty}{\to} a$ in $d_2$. A similar argument with the other inequality of~\eqref{eq:hold_1} shows that $d_2$ convergence implies $d_1$ convergence.

In order to show the equivalence of $d_2$ and $d_\infty$ convergence, one uses~\eqref{eq:hold_2}.
\end{proof}
\begin{eg}
	\label{eg_diag_vanish}
	Let us work in $\R^2$ with the standard metric $d_2$ and consider the sequence $(a_j)_{j\in\N}\subset \R^2$ defined by
	\[
	a_j := \left(\frac{1}{j},\frac{1}{j}\right).
	\]
	We claim and prove that $a_j\to (0,0)$ with respect to $d_2$ (hence, by~\Cref{prop:equiv_d12inf}, $a_j$ also converges to $(0,0)$ in $d_1$ and $d_\infty$). Before setting $\varepsilon>0$ and $\delta>0$, let's work backwards and calculate, informally, what we need to show.
	
	After fixing $\varepsilon>0$, we need to find some $N\in\N$ such that, whenever $j\ge N$, we have
	\[
	d_2(a_j,(0,0))<\varepsilon.
	\]
	We can explicitly calculate
	\[
	d_2(a_j,(0,0)) = \sqrt{\left(\frac{1}{j} - 0\right)^2 + \left(\frac{1}{j} - 0\right)^2} = \sqrt{\frac{2}{j^2}} = \frac{\sqrt{2}}{j}\overset{j\to \infty}{\to}0.
	\]
	In words, the ($\ell^2$) distance between $a_j$ and the origin is vanishing as $j\to \infty$.
	\begin{proof}
		\textcolor{blue}{Fix $\varepsilon>0$.} We know from Math 2080 / Introduction to Analysis that $\frac{\sqrt{2}}{j}\overset{j\to \infty}{\to} 0$. Hence, \textcolor{blue}{there exists $N\in\N$ such that $j\ge N$ implies}
		\[
		\left|\frac{\sqrt{2}}{j} - 0\right| = \frac{\sqrt{2}}{j} < \varepsilon.
		\]
		By the calculation done just before this proof, we therefore have, whenever $j\ge N$, that
		\[
		{\color{blue}d_2(a_j,(0,0))} = \frac{\sqrt{2}}{j} {\color{blue}<\varepsilon.}
		\]
		The text in \textcolor{blue}{blue} is precisely what it means for $a_j\to (0,0)$ according to~\Cref{defn:cvg_metric}.
	\end{proof}
\end{eg}
\begin{eg}[The discrete metric]
\label{eg:disc_metric}
Given any set $X$, we define the \textit{discrete metric} $d_\text{disc}:X\times X \to [0,\infty)$ by
\[
d_\text{disc}(x,y) := \left\{
\begin{array}{cl}
	0, 	&\text{if }x=y 	\\
	1, 	&\text{if }x\neq y
\end{array}
\right..
\]
This example shows that every set has at least one metric (the discrete metric) on it. However, in terms of measuring distance between points, the discrete metric is not great and is usually only considered for pedagogical purposes.

Let us return to the previous example of $\R^2$ with the sequence $a_j = \left(\frac{1}{j},\frac{1}{j}\right)_{j\in\N}$. This time, however, we will work the discrete metric $d_\text{disc}$. We see that
\[
d_\text{disc}(a_j,(0,0)) = 1,\quad \forall j\in\N.
\]
Hence, the discrete distance between $a_j$ and the origin can \textit{never} be less than $\varepsilon$ for any $0<\varepsilon<1$ nor any $j\in\N$. In this way, the sequence $a_j$ does \textit{not} converge to $(0,0)$ in the discrete metric.
\end{eg}
\Cref{eg:disc_metric} highlights two points:
\begin{enumerate}
	\item Every set has at least one metric on it -- the discrete metric.
	\item The convergence of a sequence \textbf{depends on the metric}.
\end{enumerate}
\begin{eg}[Metric subspaces]
	\label{eg:sub_met}
	Given a metric space $(X,d)$ and a subset $A\subset X$, we can define a metric on $A$ by restricting $d$ to $A\times A$. In detail, we define
	\[
	d_A:A\times A \to [0,+\infty),\quad d_A(x,y) = d(x,y),\quad x,y\in A.
	\]
	Of course, $(A,d_A)$ is a metric space according to~\Cref{defn:metric} since~\Cref{m:pos,m:sym,m:tri} hold for $d$ and so they are also true of $d_A$. We will refer to the metric space $(A,d_A)$ as a \textit{(metric) subspace of }$(X,d)$ and $d_A$ as the \textit{metric on $A$ induced by $d$}. When the metric $d$ is clear, we may just say that $A$ is a subspace of $X$. For subsets $A\subset \R^d$, we will refer to $A$ as a metric space with the assumption that the Euclidean metric $d_2$ is used.
	
	It is, however, often the case that we may assign a distance function $d_A$ to $A$ which is \textbf{not} the restriction of $d$ onto $A\times A$.
	
	We will revisit this notion of subspaces in some detail later on.
\end{eg}
\begin{eg}[Product spaces]
	Suppose we have two metric spaces $(X,d_X)$ and $(Y,d_Y)$. We can define several metrics on the product space $X\times Y$. Fix $x_1,x_2\in X$ and $y_1,y_2\in Y$. We may define, for example:
	\begin{itemize}
		\item $d_1((x_1,y_1),\,(x_2,y_2)):= d_X(x_1,x_2) + d_Y(y_1,y_2)$, or
		\item $d_2((x_1,y_1),\,(x_2,y_2)):= \left(d_X(x_1,x_2)^2 + d_Y(y_1,y_2)^2\right)^\frac{1}{2}$, or
		\item $d_\infty((x_1,y_1),\,(x_2,y_2)):= \max\left\{d_X(x_1,x_2),\,d_Y(y_1,y_2)\right\}$.
	\end{itemize}
	As an exercise, show that all such functions $d_1,d_2,d_\infty$ are indeed metrics on $X\times Y$.
\end{eg}
Finally, we leave as an exercise the following result.
\begin{ex}
\label{ex:cvg_metric}
Let $(X,d)$ be a metric space with a sequence $(a_n)_{n\in\N}\subset X$ and a point $a\in X$. TFAE:
\begin{enumerate}
	\item $a_n\overset{n\to \infty}{\to} a$.
	\item $d(a_n,a)\overset{n\to \infty}{\to}0$.
\end{enumerate}
\end{ex}
\begin{proof}
	This is an exercise from~\Cref{defn:cvg_metric}.
\end{proof}
Before concluding, we will introduce below function spaces as metric spaces. We will leave a detailed discussion for later. Perhaps the most important point to grasp, if this is the first exposure to function spaces, is that we will consider elements of $X$ to be functions (as opposed to merely numbers or vectors).
\begin{eg}[Sup norm for continuous functions]
	Let $a<b$ and we consider
	\[
	X:= \left\{f:[a,b]\to \R\, : \, f\text{ is continuous}\right\}.
	\]
	Given two elements $f,g\in X$, we define
	\[
	d(f,g):= \sup_{x\in[a,b]}|f(x) - g(x)|.
	\]
	According to the Extreme Value Theorem (see Math 2080), the right-hand side exists.
\end{eg}
	\section{Some point-set topology of metric spaces}
\label{sec:point_set}
\begin{defn}[Balls]
\label{defn:ball}
Let $(X,d)$ be a metric space with $x_0\in X$ and $r>0$. We define the \textit{(open) ball of radius $r$ centred around $x_0$} to be the following set
\[
B_r(x_0):= \left\{
x\in X \, : \, d(x,x_0)<r
\right\}.
\]
\end{defn}
\begin{rem}
	\begin{itemize}
		\item The notation for balls is sometimes written as $B_{(X,d)}(x_0,r) =B(x_0,r) = B_r(x_0)$. The first expression is meant to emphasise the role of the metric. See~\Cref{eg:ball_metric} where the metric can change the shape of the ball.
		\item On $\R$ (equipped with the metric $d(x,y) = |x-y|$), balls are open intervals (see below in~\Cref{eg:ball_R}):
		\[
		B_r(x_0) = (x_0-r,x_0+r).
		\]
		\item Conversely, for abstract metric spaces $(X,d)$, balls generalise the notion of (open) intervals.
	\end{itemize}
\end{rem}
\begin{eg}[Balls on $\R$]
	\label{eg:ball_R}
	As we saw in the remarks above, balls on $\R$ are just open intervals:
	\begin{proof}
	\begin{align*}
		x \in B_r(x_0) 	&\iff d(x,x_0)< r \iff |x-x_0| < r 	\\
		&\iff -r < x-x_0 < r \iff x_0-r<x<x_0+r 	\\
		&\iff x\in (x_0-r,x_0+r).
	\end{align*}
	Hence, we have proven both inclusions $B_r(x_0) \subset (x_0-r,x_0+r)$ and $B_r(x_0) \supset (x_0-r,x_0+r)$ which verifies $B_r(x_0) = (x_0-r,x_0+r)$.
	\end{proof}
\end{eg}
\begin{eg}[Balls on $\R^2$]
	\label{eg:ball_metric}
	Take $x_0$ to be the origin in the plane (denoted by $0 = (0,0)$) and $r=1$. Then, the ball $B_{(\R^2,d_2)}(0,1)$ (notice the $d_2$ metric) is the open \textbf{disc}
	\[
	B_{(\R^2,d_2)}(0,1) = \left\{
	x = (x_1,x_2) \in \R^2 \, : \, x_1^2 + x_2^2 < 1 
	\right\}.
	\]
	Indeed, the inequality $d_2(x,0) < 1$ can equivalently be expressed as
	\[
	d_2(x,0) < 1 \iff \sqrt{(x_1-0)^2 + (x_2-0)^2}< 1 \iff x_1^2 + x_2^2 < 1.
	\]
	On the other hand, the shape of the ball can change with a \textit{different} metric. If we use the $d_1$ metric, then the ball $B_{(\R^2,d_1)}(0,1)$ has a \textbf{diamond} shape
	\[
	B_{(\R^2,d_1)}(0,1) = \left\{
	x = (x_1,x_2)\in \R^2 \, : \, |x_1| + |x_2| < 1
	\right\}.
	\]
	Exercise: Draw $B_{(\R^2,d_1)}(0,1)$.
\end{eg}
\begin{ex}
	Verify in the previous example that
	\[
	B_{(\R^2,d_1)}(0,1) = \left\{
	x = (x_1,x_2)\in \R^2 \, : \, |x_1| + |x_2| < 1
	\right\}.
	\]
\end{ex}
\begin{eg}[Balls with the discrete metric]
	Let $X$ be a non-empty set and suppose $x\in X$. Let's endow $X$ with the discrete metric and set $d = d_\text{disc}$. Then, we have
	\[
	B_{(X,d)}(x,1) = \{ x \}.
	\]
	Now, if we increase the radius by choosing $r>1$, we obtain
	\[
	B_{(X,d)}(x,r) = X.
	\]
\end{eg}
\begin{ex}
	Verify the previous example.
\end{ex}
The next notion we will introduce is that of \textit{open} and \textit{closed} sets. The informal idea of open sets is that at every point, one can `fit balls inside' them.
\begin{defn}
	\label{defn:open_closed}
	Let $(X,d)$ be a metric space and let $A\subset X$. We say that $A$ is
	\begin{itemize}
		\item \textit{open} iff:
		\[
		\forall x \in A, \exists r>0\text{ such that }B_r(x) \subset A.
		\]
		\item \textit{closed} iff $X\setminus A$ is open. Here, $X\setminus A$ means the \textit{complement} of $A$ (with respect to $X$) i.e.
		\[
		X\setminus A = A^\complement = \{x \in X \, : \, x \notin A\}.
		\]
	\end{itemize}
\end{defn}
\begin{eg}
	\label{eg:R_open_closed}
	Consider $\R$ with the Euclidean metric. We usually refer to $(a,b)$ as an open interval and $[a,b]$ as a closed interval. The adjectives `open' and `closed' are consistent with~\Cref{defn:open_closed}.
\begin{itemize}
	\item The interval $(a,b)$ is open in the sense that it satisfies the definition in~\Cref{defn:open_closed}. 
	\begin{proof}
	\textcolor{blue}{Fix any $x \in (a,b)$}. Then, let us \textcolor{blue}{define}
	\[
	\textcolor{blue}{r:= \min (x-a,b-x)>0}.
	\]As we saw in~\Cref{eg:ball_R}, we have
	\[
	B_r(x) = (x-r,x+r).
	\]
	Since $r \le x-a$, we see that $a \le x-r$. Similarly, since $r\le b-x$, we have $x+r \le b$. These inequalities show that
	\[
	\textcolor{blue}{B_r(x)} = (x-r,x+r)\textcolor{blue}{\subset (a-b).} 
	\]
	The text in \textcolor{blue}{blue} yields precisely the criteria that $(a,b)$ is open.
	\end{proof}
	\item The interval $[a,b]$ is closed because its complement is $(-\infty,a)\cup (b,\infty)$ which is open according to~\Cref{defn:open_closed}. The verification of this is left as an exercise.
	\item The interval $(a,b]$ is \textit{neither open nor closed}. Exercise: Why?
	\item The empty set $\emptyset$ and $\R$ are \textit{both open and closed}. Exercise: Why?
	\item The previous points highlight that the notions of open and closed are \textbf{not} negations of each other. \textit{Sets can be exactly one of, neither of, or both of open and closed.}
\end{itemize}
\end{eg}
\begin{eg}
	\label{eg:balls_are_open}
	Let $(X,d)$ be a non-empty metric space. For any $x_0 \in X$ and $\delta>0$, the ball $B_\delta(x_0)$ is open according to~\Cref{defn:open_closed}.
	\begin{proof}
		\textcolor{blue}{Fix any $x\in B_\delta(x_0)$.} If $x=x_0$, then we are done because \textcolor{blue}{we can take $r = \delta$ and so $B_r(x) = B_\delta(x_0)$.} Assume then that $x\neq x_0$ which implies (c.f.~\ref{m:pos}) that $d(x,x_0)>0$. \textcolor{blue}{We set $r:= \delta - d(x,x_0)>0$ and fix any $y \in B_r(x)$.} By definition, this means that $d(x,y) < r$. According to~\ref{m:tri}, we have
		\[
		d(y,x_0) \le d(y,x) + d(x,x_0) < \delta - d(x,x_0) + d(x,x_0) = \delta. 
		\]
		This shows that \textcolor{blue}{$y\in B_\delta(x_0)$, hence $B_r(x) \subset B_\delta(x_0)$.}
	\end{proof}
\end{eg}
The notions of open and closed sets in~\Cref{defn:open_closed} might appear very abstract. These actually play a fundamental role in \textit{topological spaces} (generalisations of metric spaces) which we may study, time permitting. We will return to more intuitive characterisations of them. For now, let us discuss other terminology which will eventually be linked to open and closed sets.
\begin{defn}[Interior and limit points]
	\label{defn:int_lim_pts}
Let $(X,d)$ be a metric space and let $A\subset X$. We say that a point $x\in X$ is
\begin{itemize}
	\item an \textit{interior point of $A$} iff:
	\[
	\exists r>0\text{ such that }B_r(x) \subset A.
	\]
	\item a \textit{limit point of $A$} iff:
	\[
	\forall r>0, B_r(x) \cap A \neq \emptyset.
	\]
\end{itemize}
\end{defn}
During lecture, I will (hopefully remember to) illustrate these concepts. In words, $x\in X$ is an interior point of $A$ if one can always find a positive radius $r$ for which the ball $B_r(x)$ is a subset of $A$. Intuitively, one can `wiggle' $x$ by a distance $r>0$ and still remain in $A$. As for limit points, we will prove later on a result which more firmly justifies the word `limit' in this definition.
\begin{rem}
	\label{rem:int_lim_pts}
	Concerning interior points, here are some facts.
	\begin{itemize}
		\item Since $x \in B_r(x)$ for every $r>0$, it is clear that
		\[
		x\text{ is an interior point of }A\implies x\in A.
		\]
		\item $x\in A$ need \textbf{not} imply that $x$ is an interior point of $A$. A counter example, on $\R$, is $A = [a,b]$ and $x = a$ or $x=b$. Exercise: Verify this.
		\item Let's return to~\Cref{defn:open_closed}. Another way of describing open sets (exercise) is by saying
		\[
		A\text{ is open }\iff \text{every element in $A$ is an interior point of $A$}.
		\]
	\end{itemize}
	Concerning limit points, here are some facts.
	\begin{itemize}
		\item Since $x\in B_r(x)$ for every $r>0$, it is clear that
		\[
		x\in A \implies x\text{ is a limit point of $A$}.
		\]
		\item On the other hand, if $x$ is a limit point of $A$, then it is \textbf{not necessarily} an element of $A$. A counter example, on $\R$, is $A = (a,b)$ with $x = a$ or $x=b$. Exercise: Verify this.
		\item Analogous to the previous characterisation of open sets, one can show (tricky exercise)
		\[
		A\text{ is closed }\iff A\text{ contains all its limit points.}
		\]
	\end{itemize}
\end{rem}
We now prove a result which characterises limit points in a way that justifies attaching the word `limit' to the term.
\begin{lemma}[Limit points are limits of sequences]
\label{lem:lim_pts}
Let $(X,d)$ be a non-empty metric space with $A\subset X$ non-empty and $x\in X$. TFAE:
\begin{enumerate}
	\item \label{lim_pts_1}$x$ is a limit point of $A$.
	\item \label{lim_pts_2}There exists a sequence $(x_n)_{n\in\N}\subset A$ such that $x_n \to x$ as $n\to \infty$.
\end{enumerate}
\end{lemma}
\begin{proof}
	\ref{lim_pts_1} $\implies$ \ref{lim_pts_2}: We assume that $x$ is a limit point of $A$ and we wish to construct a sequence $(x_n)\subset A$ so that $x_n \overset{n\to \infty}{\to} x$. By~\Cref{defn:int_lim_pts}, we know the following:
	\begin{itemize}
		\item For $r=1>0$, we have that $B_1(x)\cap A\neq \emptyset$. This means that there is some $x_1 \in B_1(x)\cap A$. In particular, we have
		\[
		d(x_1,x) <1\quad\text{and}\quad x_1 \in A.
		\]
		\item For $r = \frac{1}{2}>0$, we have that $B_\frac{1}{2}(x)\cap A \neq \emptyset$. This means that there is some $x_2 \in B_\frac{1}{2}(x)\cap A$. In particular, we have
		\[
		d(x_2,x) < \frac{1}{2}\quad\text{and}\quad x_2\in A.
		\]
		\item \textcolor{blue}{For every $n\in\N$,} we can set $r = \frac{1}{n}>0$. We know that $B_r(x)\cap A \neq \emptyset$. Hence, this means that \textcolor{blue}{there exists some $x_n$} belonging to $B_\frac{1}{n}(x)\cap A$ yielding
		\[
		d(x_n,x)<\frac{1}{n}\quad\text{and}\quad\textcolor{blue}{x_n \in A.}
		\]
		By construction, $d(x_n,x)<\frac{1}{n}$ implies, by~\Cref{ex:cvg_metric}, that \textcolor{blue}{$x_n \overset{n\to \infty}{\to} x$.} The text in \textcolor{blue}{blue} is precisely what is demanded in this implication.
	\end{itemize}
	\ref{lim_pts_2} $\implies$ \ref{lim_pts_1}: Our starting assumption is that we have a sequence $(x_n)_{n\in\N}\subset A$ which converges to $x$ and we wish to prove that $x$ is a limit point of $A$. To this end, let us \textcolor{blue}{fix any $r>0$}. Since $x_n\to x$, for this $r>0$, we know that there exists some $N\in\N$ such that
	\[
	n\ge N \implies d(x_n,x)<r.
	\]
	In particular, we see that $d(x_N,x)<r$ and, since the sequence $(x_n)_{n\in\N}$ belongs in $A$, we know $x_N \in A$. Therefore, we obtain \textcolor{blue}{$B_r(x)\cap A \neq \emptyset$} since $x_N$ belongs to both sets.
\end{proof}
To intuitively summarise~\Cref{defn:int_lim_pts} and~\Cref{lem:lim_pts}, if $x$ is an interior point of $A$, then one can `wiggle / move' $x$ by a distance $r>0$ and still remain in $A$. If $x$ is a limit point of $A$, then one can `approximate' $x$ by constructing a sequence in $A$ which converges to $x$. The connection between interior points and open sets is quite clear from the definitions. What about limit points and closed sets?
\begin{lemma}[Closed sets consist of limit points]
\label{lem:closed_lim}
Let $(X,d)$ be a non-empty metric space and let $A\subset X$. TFAE:
\begin{enumerate}
	\item \label{closed_lim_1} $A$ is closed.
	\item \label{closed_lim_2} Suppose $(x_n)_{n\in\N}\subset A$ is a convergent sequence. Then, the limit $\lim_{n\to \infty}x_n$ must belong to $A$.
\end{enumerate}
\end{lemma}
\begin{proof}
We will show \ref{closed_lim_1} $\implies$ \ref{closed_lim_2} by proof by contradiction. The other direction can also use proof by contradiction and is left as an exercise.

Let us denote $x = \lim_{n\to \infty} x_n \in X$ and \textcolor{blue}{assume, for a contradiction, that $x \notin A$.} By assumption that $A$ is closed, we have that $X\setminus A$ is open and so there exists $r>0$ such that \textcolor{blue}{$B_r(x) \subset X\setminus A$}. However, since $(x_n)_{n\in\N}\subset A$ converges to $x$, for this $r>0$, we know that there exists some $N\in\N$ such that
\[
n\ge N \implies d(x_n,x)<r.
\]
In particular, we have \textcolor{blue}{$x_N \in B_r(x)$}. The text in \textcolor{blue}{blue} yields a contradiction because $x_N \in A$ but, simultaneously, $x_N \in B_r(x) \subset X\setminus A$.
\end{proof}
Given a set $A\subset X$, we may also be interested in the set of all interior points and / or the set of all limit points of $A$. These are defined now.
\begin{defn}[Interior and closure of a set]
	\label{defn:int_clos}
	Let $(X,d)$ be a metric space and let $A\subset X$. We define
	\begin{itemize}
		\item the \textit{interior of $A$} to be the set int$(A)$ consisting of all interior points of $A$ given by
		\[
		\text{int}(A) := \left\{
		x\in X \, : \, x\text{ is an interior point of }A
		\right\}.
		\]
		This is sometimes written as $\mathring{A} = $ int$(A)$.
		\item the \textit{closure of $A$} to be the set cl$(A)$ consisting of all limit points of $A$ given by
		\[
		\text{cl}(A):= \left\{
		x\in X \, : \, x\text{ is a limit point of }A
		\right\}.
		\]
		This is sometimes written as $\overline{A} = $ cl$(A)$.
	\end{itemize}
\end{defn}
Some remarks are in order.
\begin{rem}
	\label{rem:int_clos}
	Concerning the interior, we have the following facts.
	\begin{itemize}
		\item From~\Cref{rem:int_lim_pts}, we see that $\mathring{A} \subset A$.
		\item $\mathring{A}$ is an open set.
		\begin{proof}
			\textcolor{blue}{Fix any $x\in \mathring{A}$.} By definition, $x$ is an interior point of $A$, so \textcolor{blue}{there exists $r>0$} such that $B_r(x) \subset A$. Now, to conclude $\mathring{A}$ is open, we would like to prove that $B_r(x)\subset \mathring{A}$, which is the content of the rest of this proof. \textcolor{DarkOrchid}{Fix any $y \in B_r(x)$ and define $\rho:= r - d(x,y)$.} Notice that \textcolor{DarkOrchid}{$\rho>0$} since $y \in B_r(x)$ implies $d(x,y) < r$. We claim that $B_\rho(y) \subset B_r(x)$. According to~\ref{m:tri} and~\ref{m:sym}, for any $z \in B_\rho(y)$, we see that
			\[
			d(z,x) \le d(z,y) + d(y,x) < \rho + d(x,y) = r - d(x,y) + d(x,y) = r.
			\]
			Hence, every $z\in B_\rho(y)$ also satisfies $z\in B_r(x)$ and so $B_\rho(y) \subset B_r(x)$. Since we know that $B_r(x) \subset A$, we deduce \textcolor{DarkOrchid}{$B_\rho(y)\subset A$}. The text in \textcolor{DarkOrchid}{purple} tells us that every $y\in B_r(x)$ is an interior point of $A$. Hence, we deduce \textcolor{blue}{$B_r(x) \subset \mathring{A}$.} The text in \textcolor{blue}{blue} tells us that $\mathring{A}$ is open.
		\end{proof}
	\end{itemize}
	Concerning the closure, we have the following facts.
	\begin{itemize}
		\item From~\Cref{rem:int_lim_pts}, we see that $A \subset \overline{A}$.
		\item $\overline{A}$ is a closed set.
		\begin{proof}
			This is left as an exercise.
		\end{proof}
	\end{itemize}
\end{rem}
Other ways of thinking about the interior and closure of a set are the following:
\begin{lemma}[Interior / Closure as largest open subset / smallest closed supset]
	\label{lem:int_large}
	Let $(X,d)$ be a metric space with $A\subset X$. Then, $\mathring{A}$ is the largest open set contained in $A$ and $\overline{A}$ is the smallest closet set containing $A$. More precisely:
	\begin{enumerate}
		\item \label{int_large} Suppose $B\subset X$ is an open set such that $B \subset A$. Then, $B\subset \mathring{A}$.
		\item \label{clos_small} Suppose $B\subset X$ is a closed set such that $B \supset A$. Then, $B \supset \overline{A}$.
	\end{enumerate}
\end{lemma}
\begin{proof}
	We will only prove~\ref{int_large} and leave~\ref{clos_small} as an exercise. \textcolor{blue}{Suppose $B\subset X$ is an open set such that $B\subset A$. Fix any $x\in B$.} Since $B$ is open, we know that there exists $r>0$ such that $B_r(x) \subset B$. Since $x\in B\subset A$, the previous sentence tells us that $x$ is an interior point of $A$. Hence, we see that \textcolor{blue}{$x\in \mathring{A}$.} The text in \textcolor{blue}{blue} is exactly what it means for $B\subset \mathring{A}$.
\end{proof}
\begin{defn}[Boundary]
\label{defn:boundary}
Let $(X,d)$ be a metric space and $A\subset X$. The \textit{boundary of $A$} is the set $\partial A$ defined by
\[
\partial A = \overline{A} \setminus \mathring{A}.
\]
In words, $\partial A$ consists of elements $x\in X$ such that $x\in \overline{A}$ such that $x\notin \mathring{A}$.
\end{defn}
\begin{eg}
On $\R$, let's consider the interval $(a,b]$. Exercise: Verify the following statements.
\begin{enumerate}
	\item The interior is given by $\mathring{(a,b]} = (a,b)$.
	\item The closure is given by $\overline{(a,b]} = [a,b]$.
	\item The boundary is given by $\partial(a,b] = \{a,b\}$.
\end{enumerate}
In fact, there is nothing particular about the interval $(a,b]$: Suppose $I\subset \R$ is an interval (open / closed / half-and-half) with endpoints $-\infty < a < b < \infty$, then, we always have:
\[
\mathring{I} = (a,b),\quad \overline{I} = [a,b],\quad \partial I = \{a,b\}.
\]
So the notion of boundary, from~\Cref{defn:boundary} really does capture the intuition of `boundary' points for intervals on $\R$. Of course, \Cref{defn:boundary} applies to general metric spaces beyond $\R$.
\end{eg}
We end this section with some properties of the concepts introduced.
\begin{prop}
\label{prop:defn_top}
Let $(X,d)$ be a (non-empty) metric space with $A\subset X$ (non-empty) and $x_0\in X$. Then, we have the following statements.
\begin{enumerate}
	\item \label{prop_top_1} $A$ is open iff $A = \mathring{A}$.
	\item \label{prop_top_2} $A$ is closed iff $A = \overline{A}$.
	\item \label{prop_top_3} For any $r>0$, the ball $B_r(x_0)$ is open. For any $r>0$, the following set is closed
	\[
	\{x \in X\, : \, d(x,x_0)\le r\}.
	\]
	\item \label{prop_top_4} The singleton set $\{x_0\}$ is closed.
	\item \label{prop_top_5} Suppose $A_1,A_2,\dots, A_n\subset X$ are a finite number of open subsets of $X$. Then, their intersection $\bigcap_{j=1}^n A_j$ is also open.
	\item \label{prop_top_6} Suppose $A_1,A_2,\dots, A_n \subset X$ are a finite number of closed subsets of $X$. Then, their union $\bigcup_{j=1}^n A_j$ is also closed.
	\item \label{prop_top_7} Suppose $\{A_\alpha\}_{\alpha \in I}$ is a collection of open subsets of $X$. Here, the index set $I$ could be finite or infinite (countable or uncountable). Then, their union $\bigcup_{\alpha \in I}A_\alpha$ is also open.
	\item \label{prop_top_8} Suppose $\{A_\alpha\}_{\alpha \in I}$ is a collection of closed subsets of $X$. Here, the index set $I$ could be finite or infinite (countable or uncountable). Then, their intersection $\bigcap_{\alpha \in I}A_\alpha$ is also closed.
\end{enumerate}
\end{prop}
We will leave most of the proof as an exercise, but we will give some hints and references to the previous results which would be helpful.
\begin{rem}
\ref{prop_top_3} is deliberately phrased so that it does \textbf{not} suggest $\overline{B_r(x_0)} = \{x \in X \, : \, d(x,x_0)\le r\}$. It is always the case that $\overline{B_r(x_0)}$ and $\{x \in X \, : \, d(x,x_0)\le r\}$ are closed sets, but they are \textbf{not} necessarily the same set. Exercise: The classic counterexample is the discrete metric with $r = 1$, verify this.
\end{rem}
\begin{proof}[Proof ideas]
\begin{enumerate}
	\item Use~\Cref{lem:int_large} and~\Cref{rem:int_clos}.
	\item Use~\Cref{lem:int_large} and~\Cref{rem:int_clos}.
	\item Use~\Cref{defn:open_closed} and~\ref{m:tri}.
	\item Use~\Cref{lem:closed_lim}.
	\item Use~\Cref{defn:open_closed}.
	\item Use~\Cref{defn:open_closed}.
	\item Use~\Cref{defn:open_closed}.
	\item Use~\Cref{defn:open_closed}.
\end{enumerate}
\end{proof}
	\section{Relative topology}
\label{sec:rel_top}
\textcolor{ForestGreen}{Skip on first reading.}
When we introduced open sets in~\Cref{defn:open_closed}, the concept depended on the choice of metric (since the definition involves balls which, as we saw in~\Cref{eg:ball_metric}, changes depending on metric). For instance, on $\R$, the singleton $\{1\}$ is not open with respect to the Euclidean metric $d(x,y) = |x-y|$. However, $\{1\}$ is open with respect to the discrete metric (exercise).

Beyond just the choice of metric, the choice of \textit{ambient space} $X$ comes into play. Let us begin with a motivating example.
\begin{eg}[axis of $\R^2$ isomorphic to $\R$]
Let us consider the plane $\R^2$ with the Euclidean metric $d_2$. Within the plane, we define the $x$-axis as $X:= \{(x,0)\, :\, x\in \R\}\subset \R^2$. On the one hand, we can view $X$ as `essentially' $\R$ (it is just a 1-dimensional axis). On the other hand, we can also view it as a subspace of $\R^2$ and we will shortly return to this.

Let us consider the restriction of $d_2$ onto $X$. To recall notation, $\left.d_2\right|_{X\times X}$ means the restriction of $d_2$ to $X\times X$ i.e.
\[
\left. d_2\right|_{X\times X}: (a,b) \in X\times X \mapsto d_2(a,b)\in [0,\infty).
\]
If $a=(a_1,0) \in X$ and $b = (b_1,0) \in X$, then the restriction is nothing other than
\[
\left.d_2\right|_{X\times X}(a,b) = |a_1-b_1|.
\]
This restriction has given us two metric spaces: we have $(\R^2,d_2)$ and $(X,\left.d_2\right|_{X\times X})$ (Exercise: verify that this latter object is indeed a metric space).
Let us now consider another subset given by
\[
E = \{(x,0) \in \R^2\, : \, -1 < x < 1\}.
\]
Notice that $E \subset X$ and $E \subset \R^2$. When we view $E$ as a subspace of $\R^2$, it is \textbf{not} open. For instance, the origin $(0,0) \in E$, however any ball $B_{(\R^2,d_2)}((0,0),r)$ for $r>0$ will always contain the point $(0,\frac{r}{2})\notin E$. Although $(0,0)\in E$, it fails to be an interior point \textit{with respect to to $d_2$}.

On the other hand, when we view $E$ as a subspace of $X$, it is open \textit{with respect to $\left.d_2\right|_{X\times X}$}. For example, $(0,0)\in E$ is indeed an interior point with respect to $\left.d_2\right|_{X\times X}$. We can take, for example $r = \frac{1}{2}$ and so
\begin{align*}
B_{(X,\left.d_2\right|_{X\times X})}\left((0,0),\frac{1}{2}\right) &= \left\{
x=(x_1,0)\in X \, : \, \left.d_2\right|_{X\times X}(x,(0,0))<\frac{1}{2}
\right\} = \left\{
(x_1,0)\in X \, : \, |x_1| < \frac{1}{2}
\right\} 	\\
&= \left\{
(x_1,0) \in \R^2 \, : \, -\frac{1}{2}< x_1 < \frac{1}{2}
\right\} \subset E.
\end{align*}
\end{eg}
So changing the ambient space can change the open / closed property of a set.
\begin{eg}[Subspaces of $\R$]
Consider $\R$ with the Euclidean metric $d_2$ and define $X := (-1,1)$. By restricting $d_2$ to $X\times X$, we now have two metric spaces $(\R,d_2)$ and $(X, \left.d_2\right|_{X\times X})$. Let us consider the following set
\[
E:= [0,1).
\]
Similar to the example after~\Cref{defn:open_closed}, we see that $E$ is neither open nor closed \textit{with respect to $(\R,d_2)$}. In particular, 1 is a limit point of $E$ (exercise) but $1 \notin E$ and so (c.f.~\Cref{lem:closed_lim}) $E$ is not closed.

On the other hand, with respect to $(X,\left.d_2\right|_{X\times X})$, 1 does not even belong to $X$, so 1 is not a limit point of $E$. With respect to $(X,\left.d_2\right|_{X\times X})$, the set $E$ contains all its limit points.
\end{eg}
Let us formalise this phenomena in general.
\begin{defn}[Subspace / Relative topology]
\label{defn:rel_top}
Let $(X,d)$ be a metric space and consider subsets $E \subset Y \subset X$. By restricting $d$ to $Y\times Y$, we view $(Y,d|_{Y\times Y})$ as a metric space (distinct from $(X,d)$). We say that $E$ is \textit{relatively open / closed with respect to $Y$} iff $E$ is open / closed in the metric subspace $(Y,d|_{Y\times Y})$.
\end{defn}
To get a better understanding of relative open-ness, we start off with describing the relation between open balls.
\begin{lemma}[Relatively open balls]
\label{lem:rel_ball}
Let $(X,d)$ be a (non-empty) metric space with a (non-empty) subset $Y\subset X$ and let $y \in Y$ with $r>0$. We have
\[
B_{(Y,d|_{Y\times Y})}(y,r) = B_{(X,d)}(y,r)\cap Y.
\]
\end{lemma}
$B_{(Y,d|_{Y\times Y})}(y,r)$ is open in $Y$ and $B_{(X,d)}(y,r)$ is open in $X$. The content of~\Cref{lem:rel_ball} is that the former can (and should) be viewed as the latter intersected with $Y$.
\begin{proof}
\begin{align*}
x \in B_{(Y,d|_{Y\times Y})}(y,r) &\iff x\in \left\{
z \in Y \, : \, d|_{Y\times Y}(z,y)<r
\right\} \iff x\in \left\{
z \in Y \, : \, d(z,y)<r
\right\} 	\\
&\iff x \in Y \text{ and }d(x,y)<r 	\iff x \in Y\text{ and }x \in B_{(X,d)}(y,r) 	\\
&\iff x \in B_{(X,d)}(y,r) \cap Y.
\end{align*}
\end{proof}
The previous lemma will be used and generalised in the following result. Essentially, the role of the balls in~\Cref{lem:rel_ball} will be replaced by open sets while keeping track of what is meant by `open'.
\begin{prop}[Relatively open / closed sets]
\label{prop:rel_top}
Let $(X,d)$ be a metric space and consider subsets $E\subset Y \subset X$. Then, we have the following.
\begin{enumerate}
	\item \label{rel_top_open} $E$ is relatively open with respect to $Y$ iff there exists some $V \subset X$ which is open in $X$ and $E = V \cap Y$.
	\item \label{rel_top_closed} $E$ is relatively closed with respect to $Y$ iff there exists some $K\subset X$ which is closed in $X$ and $E = K \cap Y$.
\end{enumerate}
\begin{proof}
	We will only prove the equivalence in \ref{rel_top_open} and leave the proof of \ref{rel_top_closed} as an exercise.
	
	$\implies$: We assume that $E$ is relative open with respect to $Y$ and the goal is to construct some open $V \subset X$ such that $E = V \cap Y$. By definition, for every $x\in E$, we can find a radius $r_x>0$ (the radius depends on $x$ and we emphasise this dependence by writing $r_x$) such that $B_{(Y,d|_{Y\times Y})}(x,r_x) \subset E$. \textcolor{blue}{We construct the following set
	\[
	V := \bigcup_{x\in E} B_{(X,d)}(x,r_x).
	\]
	} Of course, we see that \textcolor{blue}{$V \subset X$}. Moreover, by~\ref{prop_top_7} in~\Cref{prop:defn_top}, as $V$ is the union of balls (which are open), we see that \textcolor{blue}{$V$ is open.} It remains to show that $E = V \cap Y$.
	
	Fix any $x\in E$. Since $E\subset Y$, we immediately obtain $x\in Y$. Moreover, since $x \in B_{(X,d)}(x,r_x)$, we also have $x\in V$. Hence, \textcolor{blue}{$E\subset V\cap Y$}.
	
	For the other containment, let us fix any $y \in V \cap Y$. Since $y \in V$, there must exist some $x\in E$ such that $y\in B_{(X,d)}(x,r_x)$. Moreover, since $y \in Y$ as well, we get that $y \in B_{(Y,d|_{Y\times Y})}(x,r_x)$. In the first paragraph, we saw that $B_{(Y,d|_{Y\times Y})}(x,r_x) \subset E$ and so $y \in E$. Hence, \textcolor{blue}{$V\cap Y \subset E$}. The text in \textcolor{blue}{blue} is precisely what was desired to prove.
	
	$\impliedby$: We assume that there is some open $V\subset X$ for which $E = V\cap Y$ and we wish to show that $E$ is relatively open with respect to $Y$. \textcolor{blue}{Fix any $x \in E$.} By assumption, we know that $x \in V$ (which is open in $X$) so \textcolor{blue}{we can find a radius $r>0$} such that $B_{(X,d)}(x,r)\subset V$. By intersecting both sides with $Y$, this containment implies
	\[
	B_{(X,d)}(x,r)\cap Y \subset V \cap Y = E.
	\]
	On the other hand, by~\Cref{lem:rel_ball} we have
	\[
	B_{(X,d)}(x,r)\cap Y = B_{(Y,d|_{Y\times Y})}(x,r).
	\]
	Putting the previous equations together gives \textcolor{blue}{$B_{(Y,d|_{Y\times Y})}(x,r) \subset E$.} According to~\Cref{defn:rel_top}, the text in \textcolor{blue}{blue} is precisely what it means for $E$ to be relatively open with respect to $Y$.
\end{proof}
\end{prop}
The pathology described in the first example of this section can be ruled out in the following case.
\begin{prop}
\label{prop:triple_rel}
Let $(X,d)$ be a metric space and consider subsets $E\subset Y \subset X$. We have the following.
\begin{enumerate}
\item \label{triple_rel_open} If $Y$ is open (with respect to $X$), then we have the equivalence:
\[
E\text{ is (relatively) open in }Y \iff E\text{ is open in }X.
\]
\item \label{triple_rel_closed} If $Y$ is closed (with respect to $X$), then we have the following equivalence:
\[
E\text{ is (relatively) closed in }Y \iff E\text{ is closed in }X.
\]
\end{enumerate}
\end{prop}
\begin{proof}
We only prove~\Cref{triple_rel_open} as the other case is very similar and also uses~\Cref{prop:rel_top} (replace all instances of the word `open' with `closed' below). 

$\implies$: If $E$ is (relatively) open in $Y$, then~\Cref{prop:rel_top} guarantees that $E = U\cap Y$ for some open $U\subset X$. Since both $Y,U$ are open in $X$, the finite intersection $E = U \cap Y$ is also open in $X$.

$\impliedby$: If $E$ is open in $X$, then since $E\subset Y$ we can write $E = E \cap Y$. In other words, we have expressed $E$ as the intersection between $Y$ and an open subset of $X$. Again by~\Cref{prop:rel_top}, we deduce that $E$ is (relatively) open in $Y$.
\end{proof}
	\section{Cauchy sequences and completeness}
\label{sec:Cauchy}
You learned from Math 2080 about Cauchy sequences on $\R$: $(a_n)_{n\in\N}\subset \R$ is said to be Cauchy iff
\[
\forall \varepsilon>0,\exists N\in\N\text{ s.t. }\forall n,m\ge N, \, |a_n-a_m|<\varepsilon.
\]
Perhaps you can already see how this easily generalises to metric spaces. Before writing down the formal definition in general metric spaces, let us remind ourselves of some of the intuition behind Cauchy sequences (of real numbers).
\begin{rem}
	\begin{itemize}
		\item Cauchy sequences describe the idea that `eventually, the elements of a Cauchy sequence get arbitrarily close to \textit{each other}.' This is in contrast with convergent sequences where the idea is that `eventually, the elements of a convergent sequence get arbitrarily close to \textit{one number (the limit)}.'
		\item You proved in Math 2080 that Cauchy sequences and convergent sequences on $\R$ are the same thing. In practice, if you wanted to show a sequence is convergent but did not have a good guess for what the limit is, it is much easier to try and prove that it is Cauchy.
		\item However, in general metric spaces it is \textbf{not} true that Cauchy sequences and convergent sequences are the same.
	\end{itemize}
\end{rem}
\begin{defn}[Cauchy]
\label{defn:Cauchy}
Let $(X,d)$ be a metric space and $(a_n)_{n\in\N}\subset X$ be a sequence. We say that $(a_n)_{n\in\N}$ is a \textit{Cauchy sequence} iff the following statement is true:
\[
\forall \varepsilon>0,\exists N\in\N\text{ such that }\forall n,m\ge N, \, d(a_n,a_m)<\varepsilon.
\]
\end{defn}
\begin{lemma}[Convergent $\implies$ Cauchy]
	\label{lem:cvg_imp_cauchy}
Let $(X,d)$ be a metric space and assume that $(a_n)_{n\in\N}\subset X$ is a convergent sequence with limit $a\in X$. Then, $(a_n)_{n\in\N}$ is Cauchy.
\end{lemma}
Compare the proof below to the proof of the statement `convergent sequences of real numbers are Cauchy.'
\begin{proof}
\textcolor{blue}{Fix any $\varepsilon>0$.} Since $(a_n)_{n\in\N}$ is convergent, \textcolor{blue}{we can find $N\in\N$} such that
\[
n\ge N \implies d(a_n,a)<\frac{\varepsilon}{2}.
\]
In particular, using~\ref{m:tri}, \textcolor{blue}{for any $n,m\ge N$, we have}
\[
\textcolor{blue}{d(a_n,a_m)} \le d(a_n,a) + d(a_m,a) \textcolor{blue}{<}\frac{\varepsilon}{2} + \frac{\varepsilon}{2} = \textcolor{blue}{\varepsilon}.
\]
The text in \textcolor{blue}{blue} is precisely what it means for $(a_n)_{n\in\N}$ to be Cauchy.
\end{proof}
As remarked earlier in this section, for general metric spaces, it is \textbf{not} true that Cauchy sequences and convergent sequences are the same. This means that the converse to~\Cref{lem:cvg_imp_cauchy} does not hold, in general.
\begin{eg}[The rationals]
Let us consider the rationals $\Q$ as a metric space with the usual Euclidean metric $d(x,y) = |x-y|$. Let us consider the following sequence
\[
3,\, 3.1, \, 3.14, \, 3.141, \, 3.1415, \, 3.14159,\, \dots
\]
i.e. finite decimal approximations of $\pi\notin \Q$. This is a Cauchy sequence, however it is not a convergent sequence.
\end{eg}
The metric spaces where Cauchy sequences are convergent sequences deserve special terminology.
\begin{defn}[Complete metric spaces]
\label{defn:complete}
We say that a metric space $(X,d)$ is \textit{complete} iff every Cauchy sequence of $(X,d)$ is, in fact, also convergent in $(X,d)$.
\end{defn}
We will revisit in detail the importance of complete metric spaces particularly when we deal with metric spaces of \textit{functions}. For now, let us reformulate the previous example with this new term.
\begin{eg}
From what you learned in Math 2080, you saw that $\R$ is a complete metric space. From the previous example, we see that $\Q$ is not complete.
\end{eg}
We saw in~\Cref{sec:rel_top} that the notions of open and closed sets can change depending on the ambient space. Complete metric spaces, however, will always remain closed no matter what ambient space they are viewed in. This is the content of the following result.
\begin{prop}[Completeness and closedness]
\label{prop:comp_clos}
Let $(X,d)$ be a metric space and let $Y\subset X$.
\begin{enumerate}
	\item \label{comp_clos_1} Assume that $(Y,d|_{Y\times Y})$ is complete. Then, $Y$ is closed in $X$.
	\item \label{comp_clos_2} Assume that $(X,d)$ is complete and that $Y$ is closed in $X$. Then, $(Y,d|_{Y\times Y})$ is complete.
\end{enumerate}
\end{prop}
\begin{proof}
\begin{enumerate}
	\item To prove this, we will use~\Cref{lem:closed_lim}. \textcolor{blue}{Fix $(x_n)_{n\in\N}\subset Y$ such that it converges to a limit $x\in X$.} For any $\varepsilon>0$, since $x_n\to x$, we know that there exists some $N\in\N$ such that
	\[
	n\ge N \implies d(x_n,x)<\frac{\varepsilon}{2}.
	\]
	In particular, using~\ref{m:tri} and~\ref{m:sym}, we have
	\[
	n,m\ge N \implies d|_{Y\times Y}(x_n,x_m) = d(x_n,x_m) \le d(x_n,x) + d(x,x_m) < \varepsilon.
	\]
	Thus, we have shown that $(x_n)_{n\in\N}\subset Y$ is a Cauchy sequence. Since we assume that $(Y,d|_{Y\times Y})$ is complete, the sequence must also be convergent to some limit $y\in Y\subset X$. Limits are unique and so \textcolor{blue}{$x= y\in Y$}. By~\Cref{lem:closed_lim}, the text in \textcolor{blue}{blue} asserts that $Y$ is closed in $X$.
	\item \textcolor{blue}{Fix $(y_n)_{n\in\N}\subset Y$ to be a Cauchy sequence}. Since $Y\subset X$, we can view $(y_n)_{n\in\N}$ as a Cauchy sequence in $X$. By assumption, $(X,d)$ is complete, so \textcolor{blue}{$(y_n)_{n\in\N}$ converges} to some limit $y\in X$. Again, by assumption, since $Y$ is closed in $X$, \Cref{lem:closed_lim} tells us that \textcolor{blue}{$y\in Y$}. The text in \textcolor{blue}{blue} give precisely that $(Y,d|_{Y\times Y})$ is complete.
\end{enumerate}
\end{proof}
\subsection{Subsequences}
\label{sec:subseq}
Another aspect of importance is that of subsequences which we review.
\begin{defn}[Subsequences]
Let $(X,d)$ be a metric space and suppose $(x_n)_{n\in \N}\subset X$ is a sequence. Suppose that $n: j\in \N \mapsto n_j \in \N$ is strictly increasing i.e.
\[
1 \le n_1 < n_2 < n_3 < \dots
\]
We define $(x_{n_j})_{j\in\N}$ to be a \textit{subsequence} of the original sequence $(x_n)_{n\in\N}$.
\end{defn}
So subsequences are built by taking an infinite subset of the original sequence and keeping the same order.
\begin{eg}
\label{eg:seq_alt}
On $\R$, consider $a_n = (-1)^n$ so that it reads
\[
1,-1,1,-1,1,-1,\dots
\]
By taking $n_j = 2j$, we have the subsequence
\[
a_{n_j} = 1,\quad \forall j\in\N.
\]
Alternatively, we could take $n_j = 2j+1$ to get the subsequence
\[
a_{n_j} = -1,\quad \forall j\in\N.
\]
\end{eg}
\begin{eg}
Consider the sequence defined by $a_n := \left(\frac{1}{n},\frac{1}{n^2}\right)\in \R^2$. Then, by taking $n_j = j^2$, one obtains a subsequence
\[
a_{n_j} = \left(\frac{1}{j^2},\, \frac{1}{j^4}\right),\quad j\in \N.
\]
\end{eg}
The following result says that all subsequences of a convergent sequence are also convergent (and to the same limit).
\begin{lemma}
\label{lem:cvg_imp_sub_cvg}
Let $(X,d)$ be a metric space and suppose $(x_n)_{n\in\N}\subset X$ is a convergent sequence with limit $x\in X$. Then, every subsequence of $(x_n)_{n\in\N}$ also converges to $x\in X$.
\end{lemma}
\begin{proof}
This is an exercise. It is an easy generalisation of what you learned in the case $X = \R$ from Math 2080.
\end{proof}
On the other hand, it is entirely possible for a sequence to not be convergent yet it might have multiple subsequences which are convergent (to different limits too!). We already saw an example of this in~\Cref{eg:seq_alt}. In order to deal with subsequences and possibly multiple limits therein, we define a familiar term.
\begin{defn}[(Subsequential) limit point]
\label{defn:subseq_lim}
Let $(X,d)$ be a metric space and suppose $(x_n)_{n\in\N}\subset X$ is a sequence with $x\in X$. We say that $x$ is a \textit{(subsequential) limit point of $(x_n)_{n\in\N}$} iff the following holds:
\[
\forall \varepsilon>0,\forall N\in\N, \exists n\ge N\text{ such that }d(x_n,x)<\varepsilon.
\]
\end{defn}
\begin{rem}
\begin{itemize}
	\item Notice the distinction between this definition and that of~\Cref{defn:cvg_metric} and~\Cref{defn:Cauchy}. The quantifier around `$N$' is `for all' instead of `there exists.'
	\item Returning to~\Cref{eg:seq_alt}, the sequence $a_n = (-1)^n$ has two (subsequential) limit points, namely 1 and -1.
	\item The notion of (subsequential) limit points in~\Cref{defn:subseq_lim} is deliberately written to remind us of limit points from~\Cref{defn:int_lim_pts}. Momentarily we will connect these two.
\end{itemize}
\end{rem}
\begin{prop}
\label{prop:subseq_lim}
Let $(X,d)$ be a metric space and suppose $(x_n)_{n\in\N}\subset X$ is a sequence with $x\in X$ such that $x\notin (x_n)_{n\in\N}$. TFAE:
\begin{enumerate}
	\item \label{subseq_lim_1} $x$ is a (subsequential) limit point of $(x_n)_{n\in\N}$ in the sense of~\Cref{defn:subseq_lim}.
	\item \label{subseq_lim_2} $x$ is a limit point of $\{x_n\}_{n\in\N}$ in the sense of~\Cref{defn:int_lim_pts}.
	\item \label{subseq_lim_3} There exists a subsequence $(x_{n_j})_{j\in\N}\subset (x_n)_{n\in\N}$ such that $x_{n_j}\overset{j\to \infty}{\to}x$.
\end{enumerate}
\end{prop}
Owing to the equivalence between~\ref{subseq_lim_1} and~\ref{subseq_lim_2}, we will just say `limit point' from now on.
\begin{proof}
Notice (convince yourself as an exercise) that~\Cref{lem:lim_pts} immediately gives us the equivalence~\ref{subseq_lim_2} $\iff$ \ref{subseq_lim_3}. We therefore only need to show~\ref{subseq_lim_1} $\iff$ \ref{subseq_lim_3}.

\ref{subseq_lim_1} $\implies$ \ref{subseq_lim_3}: By assumption that $x$ is a (subsequential) limit point for $(x_n)_{n\in\N}$, we can construct the following.
\begin{enumerate}
	\item Set $\varepsilon = 1$ and $N = 1\in \N$. There must exist some \textcolor{blue}{$n_1 \ge N = 1$} such that $d(x_{n_1},x)<1$.
	\item Set $\varepsilon = \frac{1}{2}$ and $N = n_1 + 1 \in \N$. There must exist some $n_2 \ge N = n_1+1 > n_1$ such that $d(x_{n_2},x)<\frac{1}{2}$.
	\item Set $\varepsilon = \frac{1}{3}$ and $N = n_2 + 1\in\N$. There must exist some $n_3 \ge N = n_2 + 1 > n_2$ such that $d(x_{n_3},x)<\frac{1}{3}$.
	\item Repeat this process so that, at step $j\in\N$, with $\varepsilon = \frac{1}{j}$, we can find \textcolor{blue}{$n_j > n_{j-1}$} such that $d(x_{n_j},x)<\frac{1}{j}$.
\end{enumerate}
Of course, we see that $d(x_{n_j},x)\overset{j\to \infty}{\to} 0$ which, according to~\Cref{ex:cvg_metric}, yields \textcolor{blue}{$x_{n_j}\overset{j\to \infty}{\to} x$.} The text in \textcolor{blue}{blue} outline precisely the construction of the desired subsequence.

\ref{subseq_lim_3} $\implies$ \ref{subseq_lim_1}: \textcolor{blue}{Fix $\varepsilon>0$ and $N\in\N$.} By assumption, we know that there exists some $J\in\N$ such that
\[
j\ge J \implies d(x_{n_j},x)<\varepsilon.
\]
In particular, let us define $K = \max(J,N)$ and consider $n:=n_K$. Certainly, since $K \ge J$, we have \textcolor{blue}{$d(x_n,x)<\varepsilon$.} Moreover, since $n:j\in \N \mapsto n_j\in\N$ is strictly increasing, we have $n_K \ge K \ge N$. Hence, \textcolor{blue}{$n\ge N$.} The text in \textcolor{blue}{blue} is precisely the criteria in~\Cref{defn:subseq_lim}.
\end{proof}
We have the following useful result which connects the ideas of Cauchy sequences and subsequences.
\begin{lemma}
\label{lem:cauchy_sub_cvg}
Let $(X,d)$ be a metric space and $(x_n)_{n\in\N}\subset X$ is a Cauchy sequence. Assume that there is some subsequence $(x_{n_j})_{j\in\N}\subset (x_n)_{n\in\N}$ which is convergent to a limit $x\in X$. Then, the original sequence $(x_n)_{n\in\N}$ converges to $x$.
\end{lemma}
\begin{proof}
This is an exercise. It is an easy generalisation of what you learned in the case $X = \R$ from Math 2080.
\end{proof}
\Cref{lem:cauchy_sub_cvg} says that if you have a Cauchy sequence and if you know it admits a convergent subsequence, then the entire original sequence is convergent. We end this section with a seemingly convoluted, but useful, criteria for convergence.
\begin{lemma}
\label{lem:cvg_sub_sub}
Let $(X,d)$ be a metric space, $(x_n)_{n\in\N}\subset X$ be a sequence, and fix $x\in X$. TFAE:
\begin{enumerate}
	\item \label{cvg_sub_sub_1} $x_n \overset{n\to \infty}{\to}x$.
	\item \label{cvg_sub_sub_2} For every subsequence $(x_{n_j})_{j\in \N}$, there exists a (further) subsequence $(x_{n_{j_k}})_{k\in\N}$ such that $x_{n_{j_k}}\overset{k\to \infty}{\to}x$.
\end{enumerate}
\end{lemma}
In~\ref{cvg_sub_sub_2}, $(x_{n_{j_k}})_{k\in\N}$ is a subsequence of $(x_{n_j})_{j\in\N}$ which is also a subsequence of $(x_n)_{n\in\N}$ -- two layers of subsequencing.
\begin{proof}
$\implies$: This is immediate from~\Cref{lem:cvg_imp_sub_cvg} as we can take the second subsequence $(x_{n_{j_k}})$ to just be $(x_{n_j})$.

$\impliedby$: This looks more complicated than the proof by contradiction below indicates. Suppose, for a contradiction, that $x_n$ does \textit{not} converge to $x$ as $n\to \infty$. According to~\Cref{defn:cvg_metric}, this means that there exists some $\varepsilon>0$ such that
\begin{equation}
	\label{eq:cvg_sub_sub}
	\forall N\in\N,\,\exists n\ge N\text{ s.t. }d(x_n,x)\ge \varepsilon.
\end{equation}
We will construct a subsequence $(x_{n_j})_{j\in\N}$ by repeatedly using~\eqref{eq:cvg_sub_sub}:
\begin{enumerate}
	\item Set $N = 1$. Then~\eqref{eq:cvg_sub_sub} says that there exists $n_1\ge N = 1$ such that $d(x_{n_1},x)\ge \varepsilon$.
	\item Set $N = n_1+1$. Then~\eqref{eq:cvg_sub_sub} says that there exists $n_2\ge N>n_1$ such that $d(x_{n_2},x)\ge \varepsilon$.
	\item Set $N = n_2+1$. Then~\eqref{eq:cvg_sub_sub} says that there exists $n_3\ge N > n_2$ such that $d(x_{n_3},x)\ge \varepsilon$.
	\item Iterate the previous steps so that we obtain a subseqeunce $(x_{n_j})_{j\in\N}$ satisfying
	\begin{equation}
		\label{eq:cvg_sub_sub_seq}
	d(x_{n_j},x)\ge \varepsilon,\quad \forall j\in \N.
	\end{equation}
\end{enumerate}
By assumption, we can extract a subsequence $(x_{n_{j_k}})_{k\in\N}\subset (x_{n_j})_{j\in\N}$ such that $x_{n_{j_k}}\overset{k\to \infty}{\to}x$. But this convergence implies that there exists some $K\in\N$ such that
\[
d(x_{n_{j_K}},x)< \frac{\varepsilon}{2}.
\]
However, this is in direct contradiction with~\eqref{eq:cvg_sub_sub_seq}.
\end{proof}
	\section{Compact metric spaces}
\label{sec:cpct_met_spaces}
We now come to one of the most important concepts in point-set topology which is central to many areas of research. Just to motivate this discussion from the previous concepts of subsequences, suppose you are trying to approximate a function $f$ (which many people do in machine learning, partial differential equations, optimisation, \dots). Suppose you have a (computer) algorithm which can generate a sequence of functions $(f_n)_{n\in\N}$. You believe, to some extent, that $f_n$ should converge to $f$ thereby validating your approximation. But how can you prove this convergence? In general, this is quite difficult. Perhaps you would be satisfied if only a \textit{subsequence} $(f_{n_j})_{j\in\N}$ converges. This is not as good as convergence of the full sequence, but we have more tools available that can yield positive answers to convergence along subsequences. This is expressed in the following idea of `sequential compactness' (although, at first glance, it might not appear so).
\begin{defn}[Sequential compactness]
\label{defn:seq_cpct}
Let $(X,d)$ be a metric spaces with $Y\subset X$. We say that
\begin{itemize}
	\item $(X,d)$ is \textit{sequentially compact} iff every sequence in $X$ admits at least one convergent subsequence.
	\item the subset $Y$ is \textit{sequentially compact in $X$} iff $(Y,d|_{Y\times Y})$ is sequentially compact.
\end{itemize}
\end{defn}
\begin{rem}
\begin{itemize}
	\item The notion of sequential compactness is \textit{intrinsic}. This property, unlike open and closedness, does not change depending on the ambient space. For this reason, we will usually just say $Y$ is sequentially compact when $Y$ is sequentially compact \textit{in $X$}.
	\item To clarify, in sequential compactness, the limit of the convergent subsequence \textbf{must also belong in that set}.
	\item Another way of looking at sequential compactness is the following: sequentially compact spaces are those for which every subset is guaranteed to have at least one limit point.
\end{itemize}
\end{rem}
In Math 2080, you learned the famous Bolzano-Weierstrass theorem which we remind ourselves now.
\begin{thm}[Bolzano-Weierstrass]
\label{thm:bw}
Let $(a_n)_{n\in\N}\subset \R$ be a bounded sequence. Then, it admits at least one convergent subsequence.
\end{thm}
\begin{cor}
	\label{cor:bw}
On $\R$, every closed and bounded set is sequentially compact.
\end{cor}
\begin{proof}
	This is an exercise. Use~\Cref{thm:bw} and~\Cref{lem:closed_lim}.
\end{proof}
The previous example suggest that sequential compactness and boundedness are both connected to each other. On $\R$, it is easy to describe what it means for a set to be bounded and we generalise this to metric spaces.
\begin{defn}[Bounded]
	\label{defn:bdd}
Let $(X,d)$ be a metric space with $Y\subset X$. We say that:
\begin{itemize}
	\item $Y$ is \textit{bounded} iff $\forall x\in X, \exists r>0$ such that $Y\subset B_r(x)$.
	\item $(X,d)$ is \textit{bounded} iff $X$ is bounded.
\end{itemize}
\end{defn}
\begin{ex}
Let $(X,d)$ be a non-empty metric space with non-empty $Y\subset X$. Show that TFAE:
\begin{enumerate}
	\item $Y$ is bounded.
	\item $\exists x\in X, r>0$ such that $Y\subset B_r(x)$.
\end{enumerate}
\textit{Hint: Notice the change in quantifier. Use~\ref{m:tri}.}
\end{ex}
\begin{prop}[Sequentially compact $\implies$ complete and bounded]
	\label{prop:seq_cpct_cmpl_bdd}
Let $(X,d)$ be a sequentially compact metric space. Then, $(X,d)$ is both complete and bounded.
\end{prop}
\begin{proof}
\ul{Complete:} \textcolor{blue}{Fix any Cauchy sequence $(x_n)_{n\in\N}\subset X$.} Since $X$ is sequentially compact, \Cref{defn:seq_cpct} tells us that there is a subsequence $(x_{n_j})_{j\in\N}\subset (x_n)_{n\in\N}$ which is convergent to some limit $x\in X$. According to~\Cref{lem:cauchy_sub_cvg}, \textcolor{blue}{the original sequence $(x_n)_{n\in\N}$ converges to $x$}.

\ul{Bounded:} \textcolor{blue}{Assume, for a contradiction, that $X$ is not bounded.} This means that there exists some $x\in X$ such that, for every $r>0$, we have $X \nsubseteq B_r(x)$. Alternatively, for every $r>0$, there exists some element $x_r \in X$ (this element depends on $r>0$) but $x_r \notin B_r(x)$. In particular,
\begin{enumerate}
	\item For $r=1$, there exists $x_1 \in X$ such that $x_1 \notin B_1(x)$ i.e. $d(x_1,x)\ge 1$.
	\item For $r=2$, there exists $x_2 \in X$ such that $x_2 \notin B_2(x)$ i.e. $d(x_2,x) \ge 2$.
	\item Repeating this process, for $r = n\in\N$, there exists $x_n\in X$ such that $d(x_n,x)\ge n$.
\end{enumerate}
We have therefore constructed a sequence $(x_n)_{n\in \N}\subset X$ such that $d(x_n,x)\ge n$ for every $n\in\N$. Since $(X,d)$ is sequentially compact, there must exist a subsequence $(x_{n_j})_{j\in\N}\subset (x_n)_{n\in\N}$ which is convergent to some limit $y\in X$. Let us denote $r:= d(x,y)\ge 0$. Since $x_{n_j}\overset{j\to \infty}{\to} y$ ($r$ does not depend on $n\in\N$), we can find $J_1\in \N$ such that
\[
j\ge J_1 \implies d(x_{n_j},y)<1.
\]
Now, let us take $J_2\in\N$ large such that $n_{J_2} - r \ge 1$ (why does this exist? Exercise). Let us define $J = \max (J_1,J_2)$. On the one hand, by construction of the sequence $(x_n)_{n\in\N}$ and~\ref{m:tri}, since $J\ge J_2$ we have
\[
\textcolor{blue}{1\le} n_{J} - r \le d(x_{n_J},x) - d(x,y) \le  \textcolor{blue}{d(x_{n_J},y)}.
\]
On the other hand, since $x_{n_j}\overset{j\to \infty}{\to}y$ and $J \ge J_1$, we have
\[
\textcolor{blue}{d(x_{n_J},y)<1.}
\]
The text in \textcolor{blue}{blue} reaches the desired contradiction.
\end{proof}
\begin{cor}[Sequentially compact $\implies$ closed and bounded]
\label{cor:seq_cpct_clos_bdd}
Let $(X,d)$ be a metric space and $Y\subset X$ is sequentially compact. Then, $Y$ is closed (in $X$) and bounded (in $X$).
\end{cor}
\begin{proof}
This is an exercise. Use~\Cref{prop:seq_cpct_cmpl_bdd} and~\Cref{prop:comp_clos}.
\end{proof}
The above~\Cref{cor:seq_cpct_clos_bdd} is one half of the famous Heine-Borel theorem. The other direction is only true in $\R^n$.
\begin{thm}[Heine-Borel]
	\label{thm:hb_real}
Consider $(\R^n,d)$ with the metric being either $d=d_1,d_2,$ or $d=d_\infty$. Let $E\subset \R^n$. The following are equivalent.
\begin{enumerate}
	\item \label{hb_real_1} $E$ is sequentially compact.
	\item \label{hb_real_2} $E$ is closed and bounded.
\end{enumerate}
\end{thm}
\begin{proof}
\ref{hb_real_1} $\implies$ \ref{hb_real_2} comes for free from~\Cref{cor:seq_cpct_clos_bdd}. We show the other direction only for $d=d_2$ and leave the other metrics as an exercise.

\ref{hb_real_2} $\implies$ \ref{hb_real_1}: The idea here is to apply~\Cref{thm:bw} to each coordinate. Since $E$ is bounded, there must exist $R>0$ such that
\begin{align*}
E \subset B_R(0) &= \left\{
x=(x_1,x_2,\dots, x_n)\in \R^n \, : \, d_2(x,0)<R
\right\} 	\\
&= \left\{
x=(x_1,x_2,\dots, x_n)\in \R^n \, : \, x_1^2 + x_2^2 + \dots+ x_n^2 < R^2
\right\}.
\end{align*}
In particular, if we consider the cube $Q = [-R,R]^n$, we see that $E\subset B_R(0) \subset Q$. Now, let us \textcolor{blue}{fix a sequence $(x^j)_{j\in \N}\subset E$} (the index is a superscript). Since $E\subset Q$, each component satisfies
\[
x_i^j \in [-R,R],\quad \forall i=1,2,\dots, n.
\]
\begin{enumerate}
	\item Let us focus on the $i=1$ component $(x_1^j)_{j\in\N}\subset [-R,R]\subset \R$. This sequence is bounded, so~\Cref{thm:bw} yields a convergent subsequence (which we do not relabel) such that $x_1^j \to x_1$ for some $x_1\in \R$ as $j\to \infty$. We do not write the relabeled subsequence to keep the notation as light as possible.
	\item Let us focus on the $i=2$ component $(x_2^j)\subset [-R,R]\subset \R$ \textbf{along the subsequence indexed by $j$ from the previous step}. Using~\Cref{thm:bw} again, we obtain \textit{yet another convergent subsequence (not relabeled)} such that $x_2^j \to x_2$ for some $x_2 \in \R$ as $j\to \infty$. Along this (further) subsequence, we still have $x_1^j \to x_1$ as $j\to \infty$.
	\item Let us focus on the $i=3$ component $(x_3^j)\subset [-R,R]\subset \R$ \textbf{along the subsequence indexed by $j$ from the previous step}. Using~\Cref{thm:bw} again, we obtain \textit{yet another convergent subsequence (not relabeled)} such that $x_3^j \to x_3$ for some $x_3 \in \R$ as $j\to \infty$. Along this (further) subsequence, we still have $x_1^j\to x_1$ and $x_2^j \to x_2$ as $j\to \infty$.
	\item We repeat this all the way down to the last $i=n$ component to find, for some $x = (x_1,x_2,\dots, x_n)\in \R^n$ that
	\[
	x_i^j \to x_i,\quad\text{as }j\to \infty,\quad \forall i=1,\dots, n.
	\]
	Again, this convergence is along some subsequence (we took a subsequence of a subsequence $n$ times) which we do not label for notational convenience.
\end{enumerate}
Each component of the original sequence admits a convergent subsequence
\[
x_i^j\overset{j\to \infty}{\to} x_i,\quad \forall i=1,\dots, n.
\]
\textcolor{blue}{Along the subsequence} above, this is equivalent (exercise) to the convergence
\[
\textcolor{blue}{x^j \overset{j\to \infty}{\to} x} = (x_1,x_2,\dots, x_n)\in \R^n.
\]
Finally, since $E$ is closed, we use~\Cref{lem:closed_lim} to deduce that \textcolor{blue}{$x\in E$.} The text in \textcolor{blue}{blue} shows that $E$ is sequentially compact.
\end{proof}
\begin{rem}
	\Cref{thm:hb_real} is \textbf{not} true for more general metric spaces. One needs to replace closed and boundedness with stronger notions (c.f.~\Cref{thm:hb_gen} and problem sheets).
\end{rem}
\begin{ex}
	\label{ex:r_closed_bdd}
	Let $A\subset \R$ be a non-empty, closed, and bounded set. Prove that $\min A$ and $\max A$ exist i.e. there exists $a_m,a_M\in A$ such that
	\[
	a_m \le a \le a_M,\quad \forall a\in A.
	\]
	\begin{proof}[Hint:]
		The non-empty and boundedness give the existence of $\inf A$ and $\sup A$ (why?). The closedness of $A$ gives $\inf A = \min A$ and $\max A = \sup A$ (why?).
	\end{proof}
\end{ex}
\begin{ex}
Consider the integers $\mathbb{Z}$ equipped with the discrete metric $d_\text{disc}$.
\begin{itemize}
	\item Show that $(\mathbb{Z},d_\text{disc})$ is a closed and bounded metric space.
	\item Show that $(\mathbb{Z},d_\text{disc})$ is not sequentially compact.
\end{itemize}
\end{ex}
As the above exercise demonstrates, closed and boundedness are not enough to guarantee sequential compactness. These must be replaced with stronger conditions to give the following result.
\begin{defn}[Totally bounded]
\label{defn:tot_bdd}
We say that a (non-empty) metric space $(X,d)$ is \textit{totally bounded} iff for every $\varepsilon>0$, there exists a finite number of points $x^{(1)},x^{(2)},\dots, x^{(N)}$ for some $N\in\N$ such that
\[
X = \bigcup_{i=1}^NB_\varepsilon(x^{(i)}).
\]
\end{defn}
\begin{thm}[Heine-Borel (general)]
\label{thm:hb_gen}
Let $(X,d)$ be a non-empty metric space. TFAE:
\begin{enumerate}
	\item \label{hb_1} $X$ is sequentially compact.
	\item \label{hb_2} $X$ is complete and totally bounded.
\end{enumerate}
\end{thm}
\begin{proof}
	\ul{$\implies$:} According to~\Cref{prop:seq_cpct_cmpl_bdd}, we deduce that $X$ is complete and bounded. We thus aim to show that sequential compactness implies total boundedness. Assume, for a contradiction, that $X$ is not totally bounded. This means that there exists $\varepsilon>0$ such that, for every $N\in\N$ and any choice of point $\{x^{(i)}\}_{i=1}^N\subset X$, we have
	\[
	\bigcup_{i=1}^NB_\varepsilon(x^{(i)}) \subsetneq X.
	\]
	We construct a sequence $(x_n)_{n\in\N}\subset X$ in the following way:
	\begin{enumerate}
	\item Since $X$ is non-empty, there exists an element which we label as $x_1\in X$.
	\item Since $X$ is not totally bounded, we know that
	\[
	B_\varepsilon(x_1)\subsetneq X.
	\]
	Hence, we can find an element $x_2 \in X \setminus B_\varepsilon(x_1)$. In other words, we have
	\[
	d(x_2,x_1)\ge \varepsilon.
	\]
	\item Since $X$ is not totally bounded, we know that $\bigcup_{i=1}^2B_\varepsilon(x_i) \subsetneq X$ and so we can find an element $x_3\in X\setminus (\bigcup_{i=1}^2B_\varepsilon(x_i))$. In other words, we have
	\[
	d(x_3,x_i) \ge \varepsilon,\quad i=1,2.
	\]
	\item Since $X$ is not totally bounded, we know that $\bigcup_{i=1}^3B_\varepsilon(x_i) \subsetneq X$ and so we can find an element $x_4 \in X\setminus\left(\bigcup_{i=1}^3B_\varepsilon(x_i)\right)$. In other words, we have
	\[
	d(x_4,x_i) \ge \varepsilon,\quad i=1,2,3.
	\]
	\item Continuing in this way, we can construct a sequence $(x_n)_{n\in\N}\subset X$ such that
	\begin{equation}
	\label{eq:tot_bdd_seq}
	d(x_m,x_n) \ge \varepsilon,\quad \forall m > n.
	\end{equation}
	\end{enumerate}
	Since $X$ is sequentially compact, we can extract a convergent subsequence $(x_{n_j})_{j\in\N}\subset (x_n)_{n\in\N}$. We deduce that $(x_{n_j})_{j\in\N}$ must also be Cauchy, meaning that we can find $J\in\N$ such that
	\[
	k,l\ge J \implies d(x_{n_k},x_{n_l}) <\varepsilon.
	\]
	However, this contradicts~\eqref{eq:tot_bdd_seq} since we can take $k = J+1, l = J$ which will give $n_k > n_l$ and so
	\[
	\varepsilon \le d(x_{n_k},x_{n_l}) < \varepsilon.
	\]
	\ul{$\impliedby$:} \textcolor{blue}{Let $(x_n)_{n\in\N}\subset X$ be a sequence.} We now use the totally bounded property with $\varepsilon = \frac{1}{k}$ for $k\in \N$.
	\begin{enumerate}
		\item There exists a finite number of points $a^{1,1},a^{1,2},\dots, a^{1,N_1}$ for some $N_1\in\N$ such that
		\[
		X = \bigcup_{i=1}^{N_1}B_1(a^{1,i}).
		\]
		At least one of the balls, say $B_1(a^{1,i_1})$ for some $i_1=1,\dots, N_1$, must contain an infinite number of terms in the sequence $(x_n)_{n\in\N}$ (why?). Let us denote $B_1 := B_1(a^{1,i_1})$ and refer to $A_1\subset \N$ as the (infinite) index set for which $(x_n)_{n\in A_1} \subset B_1$.
		\item There exists a finite number of points $a^{2,1},a^{2,2},\dots, a^{2,N_2}$ for some $N_2\in\N$ such that
		\[
		X = \bigcup_{i=1}^{N_2}B_\frac{1}{2}(a^{2,i}).
		\]
		At least one of the balls, say $B_\frac{1}{2}(a^{2,i_2})$ for some $i_2=1,\dots, N_2$, must contain an infinite number of elements of $(x_n)_{n\in A_1}$. Let us denote $B_2 := B_\frac{1}{2}(a^{2,i_2})$ and refer to $A_2 \subset A_1$ as the (further) subsequence for which $(x_n)_{n\in A_2}\subset B_2$.
		\item Repeat above for $\varepsilon = \frac{1}{3}$ and so on.
	\end{enumerate}
	We have thus extracted a nested index set $\N \supset A_1 \supset A_2\supset\dots$ From this, we will define a subsequence of $(x_n)_{n\in\N}$ in the following way. Since each $A_j$ is infinite, we can choose $n_j \in A_j$ such that $n_j < n_{j+1}$ for every $j\in\N$. This gives $(x_{n_j})_{j\in\N}$ such that $x_{n_j}\in B_\frac{1}{j}$ for each $j\in\N$. We seek to prove that $(x_{n_j})_{j\in\N}$ is Cauchy. This will finish the proof because the assumption of completeness will give us that \textcolor{blue}{$(x_{n_j})_{j\in\N}$ is convergent.} 
	
	As described above, we aim to show that $(x_{n_j})_{j\in\N}$ is Cauchy. \textcolor{blue}{Fix any $\varepsilon>0$. Take $J\in \N$} large enough such that $\frac{1}{J}< \frac{\varepsilon}{2}$. \textcolor{blue}{For any $j,k\ge J$,} we have that $x_{n_j},x_{n_k}$ belong to some ball with radius at most $\frac{1}{J}$. The triangle inequality then gives
	\[
	\textcolor{blue}{d(x_{n_j},d_{n_k})}\le \frac{2}{J} \textcolor{blue}{< \varepsilon.}
	\]
\end{proof}

\subsection{Compactness (independent of sequences)}
There is a more `fundamental' notion of compactness which does not involve sequences and does not rely on the metric. This section is intended to lend itself to a study of \textit{topological spaces} (generalisations of metric spaces) but \textcolor{ForestGreen}{may be skipped on first reading.} We introduce it here but we first need some preliminary definitions.
\begin{defn}[Covers]
Let $Y\subset X$ be two sets. Let $\{V_\alpha\}_{\alpha \in A}$ be a collection of subsets $V_\alpha \subset X$ indexed by $\alpha \in A$.
\begin{itemize}
	\item We say that $\{V_\alpha\}_{\alpha\in A}$ is \textit{finite} iff the index set $A$ is a finite set.
	\item We say that $\{V_\alpha\}_{\alpha \in A}$ \textit{covers / is a cover for $Y$} iff $Y \subset \bigcup_{\alpha \in A}V_\alpha.$
	\item For $B\subset A$, we say that $\{V_\alpha\}_{\alpha \in B}$ is a \textit{subcover for $Y$} iff $\{V_\alpha\}_{\alpha \in B}$ is a cover for $Y$.
\end{itemize}
When $\{V_\alpha\}_{\alpha \in A}$ is a cover for $Y$ with $V_\alpha$ open (relative to $X$), we refer to $\{V_\alpha\}_{\alpha \in A}$ as an \textit{open cover for $Y$}.
\end{defn}
\begin{eg}
Let $X = \R$ and $Y = (-1,1)$. Let's define $A = (0,2)$ and $V_\alpha = (-\alpha,\alpha)$ for $\alpha \in A$. It is clear to see that $\{V_\alpha\}_{\alpha \in A}$ is a cover of $Y$ because
\[
(-1,1) = Y \subset \bigcup_{\alpha \in A} V_\alpha = \bigcup_{\alpha \in (0,2)}(-\alpha,\alpha)
\]
In this example, the index set $A$ is infinite. But we could just as easily consider $B = \{2\}$ so that
\[
\bigcup_{\alpha \in B}V_\alpha = (-2,2) \cap Y.
\]
So $\{V_\alpha\}_{\alpha \in B}$ is a finite cover of $Y$.
Exercise: What is the set $\bigcup_{\alpha \in (0,2)}(-\alpha,\alpha)$?
\end{eg}
\begin{eg}
Let $X = \R$ and $Y = (-\infty, 0)$. Let's define $A = \N$ and $V_\alpha = (-\alpha - 2, -\alpha+1)$ for $\alpha \in A$. Then (exercise), $\{V_\alpha\}_{\alpha \in A}$ is a cover of $Y$.
\end{eg}
\begin{defn}[Compact]
	\label{defn:cpct}
Let $(X,d)$ be a metric space with $Y\subset X$. We say that
\begin{itemize}
	\item $(X,d)$ is \textit{compact} iff the following holds. Let $\{V_\alpha\}_{\alpha \in A}$ be any collection of open subsets of $X$ which cover $X$ i.e.
	\[
	V_\alpha\text{ is open }\forall \alpha \in A\text{ and }\bigcup_{\alpha \in A}V_\alpha = X.
	\]
	Then, there exists a finite subset $F\subset A$ such that $\{V_\alpha\}_{\alpha \in F}$ also covers $X$.
	\item $Y$ is \textit{compact} iff $(Y,d|_{Y\times Y})$ is compact.
\end{itemize}
\end{defn}
\Cref{defn:cpct} is usually expressed in an abbreviated way by saying `$X$ is compact iff every open cover of $X$ admits a finite subcover.' As remarked in~\Cref{defn:seq_cpct}, the notion of compactness is \textit{intrinsic} -- this property does not change even when the ambient space changes.
\begin{lemma}
\label{lem:cpct_int}
Let $(X,d)$ be a metric space with $Y\subset X$. TFAE:
\begin{enumerate}
	\item \label{cpct_int_1} $Y$ is compact.
	\item \label{cpct_int_2} If $\{V_\alpha\}_{\alpha \in A}$ is a collection of open subsets of $X$ which cover $Y$, then there exists a finite subset $F\subset A$ such that $\{V_\alpha\}_{\alpha \in F}$ also covers $Y$.
\end{enumerate}
\end{lemma}
\begin{proof}
\ref{cpct_int_1} $\implies$ \ref{cpct_int_2}: Assume $Y$ is compact and let us \textcolor{blue}{fix any collection $\{V_\alpha\}_{\alpha \in A}$ of open subsets of $X$ which cover $Y$.} In other words,
\[
V_\alpha\text{ is open in $X$ for every $\alpha \in A$ and }\bigcup_{\alpha \in A}V_\alpha \supset Y.
\]
We define the collection $\{W_\alpha\}_{\alpha \in A}$ of subsets of $Y$ by
\[
W_\alpha = V_\alpha \cap Y,\quad \forall \alpha \in A.
\]
According to~\Cref{prop:rel_top}, all these sets $W_\alpha$ are (relatively) open in $Y$. Moreover, they form a cover of $Y$ because
\[
Y\subset \bigcup_{\alpha \in A}V_\alpha \implies Y = Y\cap Y \subset \left(\bigcup_{\alpha \in A}V_\alpha\right)\cap Y = \bigcup_{\alpha \in A} (V_\alpha \cap Y) = \bigcup_{\alpha \in A}W_\alpha.
\]
Since $Y$ is compact, \textcolor{blue}{there exists a finite subset $F\subset A$} such that
\[
Y\subset \bigcup_{\alpha \in F}W_\alpha = \bigcup_{\alpha \in F}(V_\alpha \cap Y) = \left(\bigcup_{\alpha \in F}V_\alpha\right)\cap Y.
\]
This implies (exercise) that \textcolor{blue}{$Y \subset \bigcup_{\alpha \in F}V_\alpha$.} The text in \textcolor{blue}{blue} is what we wanted to prove.

\ref{cpct_int_2} $\implies$ \ref{cpct_int_1}: \textcolor{blue}{Fix any collection $\{W_\alpha\}_{\alpha \in A}$ of (relatively) open subsets of $(Y,d|_{Y\times Y})$ which cover $Y$.} In other words,
\[
W_\alpha\text{ is relatively open in $Y$ for every $\alpha \in A$ and }Y = \bigcup_{\alpha \in A}W_\alpha = Y.
\]
By~\Cref{prop:rel_top}, we can find $V_\alpha \in X$ such that
\[
V_\alpha \text{ is open in $X$ for every $\alpha \in A$ and }W_\alpha = V_\alpha \cap Y.
\]
Since $\{W_\alpha\}_{\alpha \in A}$ is a cover for $Y$, we also see that $\{V_\alpha\}_{\alpha \in A}$ covers $Y$ (and is a collection of open subsets of $X$) because (exercise)
\[
Y = \bigcup_{\alpha \in A}W_\alpha = \left(\bigcup_{\alpha \in A}V_\alpha\right)\cap Y \implies Y \subset \bigcup_{\alpha \in A}V_\alpha.
\]
Using the assumption~\ref{cpct_int_2}, \textcolor{blue}{there must exist a finite subset $F\subset A$} such that $\{V_\alpha\}_{\alpha \in F}$ covers $Y$ i.e.
\[
Y \subset \bigcup_{\alpha \in F}V_\alpha.
\]
This implies that \textcolor{blue}{$\{W_\alpha\}_{\alpha \in F}$ also covers $Y$} since
\[
Y = Y\cap Y \subset \left(\bigcup_{\alpha \in A}V_\alpha\right)\cap Y = \bigcup_{\alpha \in A}W_\alpha.
\]
The text in \textcolor{blue}{blue} is what we wanted to prove.
\end{proof}
From now on, if we want to check that $Y$ is compact as a subset of an ambient metric space $(X,d)$, we will check~\ref{cpct_int_2} in light of~\Cref{lem:cpct_int}.
\begin{eg}
Suppose $(X,d)$ is a non-empty metric space with $x\in X$. Then, the singleton $\{x\}$ is a compact set.
\begin{proof}
	\textcolor{blue}{Fix any open cover $\{V_\alpha\}_{\alpha \in A}$ of $\{x\}$.} By definition, we must have
	\[
	\{x\}\subset \bigcup_{\alpha \in A}V_\alpha \implies x\in \bigcup_{\alpha \in A}V_\alpha.
	\]
	Hence, there must exist an index $\beta\in A$ such that $x \in V_\beta$. \textcolor{blue}{Let us define the finite set $F = \{\beta\}$}. We then see that
	\[
	x \in V_\beta \implies x \in \bigcup_{\alpha \in F}V_\alpha \implies \textcolor{blue}{\{x\}\subset \bigcup_{\alpha \in F}V_\alpha}.
	\]
	The text in \textcolor{blue}{blue} gives~\ref{cpct_int_2} which is equivalent to saying $\{x\}\subset X$ is compact.
\end{proof}
\end{eg}
\begin{eg}
$\R$ is \textbf{not} compact.
\begin{proof}
We can exhibit an open cover of $\R$ which does not have a finite subcover. Namely, consider
\[
V_\alpha = (-\alpha,\alpha),\quad \alpha \in A = \mathbb{N}.
\]
It is easy to see (exercise) that $\R = \bigcup_{\alpha \in A}V_\alpha$. However, suppose $F\subset A$ is any finite (non-empty) subset of $\mathbb{N}$. Let us refer to $f := \max F < +\infty$. We see that
\[
\bigcup_{\alpha \in F}V_\alpha \subset (-f,f) \neq \R.
\]
In other words, no finite subset of $\{V_\alpha\}_{\alpha \in A}$ is a cover of $\R$.
\end{proof}
\end{eg}
Although~\Cref{defn:cpct} might seem weird and abstract, it turns out that it is equivalent to~\Cref{defn:seq_cpct} \textit{for metric spaces.}
\subsection{Equivalence in metric spaces}
\label{sec:cpct_seq}
This section is dedicated to proving~\Cref{thm:cpct_equiv} below.
\begin{thm}
\label{thm:cpct_equiv}
Let $(X,d)$ be a metric space. TFAE:
\begin{enumerate}
	\item \label{cpct_equiv_1} $(X,d)$ is compact.
	\item \label{cpct_equiv_2} $(X,d)$ is sequentially compact.
\end{enumerate}
\end{thm}
\begin{rem}
	\Cref{thm:cpct_equiv} is not true in the general setting of \textit{topological spaces.}
\end{rem}
We will need the following lemma.
\begin{lemma}[Lebesgue number]
	\label{lem:Lebesgue_number}
	Let $(X,d)$ be sequentially compact and suppose $\{U_\alpha\}_{\alpha \in A}$ is an open cover of $X$. Then, there exists $r>0$ such that
	\[
	\forall x \in X,\exists \alpha \in A\text{ s.t. }B_r(x)\subset U_\alpha.
	\]
\end{lemma}
We call $r>0$ from~\Cref{lem:Lebesgue_number} a \textit{Lebesgue number} of $\{U_\alpha\}_{\alpha \in A}$ for $X$.
\begin{proof}[Proof of~\Cref{lem:Lebesgue_number}]
	Let us assume, for a contradiction, that no such Lebesgue number exists. Specifically, for every $r>0$, we can find $x\in X$ (this $x$ depends on $r$) such that, for any $\alpha \in A$, we have
	\[
	B_r(x)\not\subset U_\alpha.
	\]
	In particular, for $n\in\N$, we choose $r = \frac{1}{n}$ to obtain a sequence $(x_n)_{n\in\N}$ such that
	\begin{equation}
		\label{eq:Leb_1}
	B_\frac{1}{n}(x_n)\not\subset U_\alpha,\quad \forall \alpha \in A,\quad \forall n\in\N.
	\end{equation}
	Since $X$ is sequentially compact, there exists a subsequence $(x_{n_k})_{k\in\N}\subset (x_n)_{n\in\N}$ and $x\in X$ such that $x_{n_k}\overset{k\to \infty}{\to}x$. Since $\{U_\alpha\}_{\alpha \in A}$ is an open cover of $X$, we must have $x\in U_{\alpha^*}$ for some $\alpha^*\in A$. Since $U_{\alpha^*}$ is open, according to~\Cref{defn:open_closed}, there must exist $\varepsilon>0$ such that
	\[
	B_\varepsilon(x)\subset U_{\alpha^*}.
	\]
	Since $(x_{n_k})_{k\in\N}$ converges to $x$, for the $\varepsilon>0$ above, there must exist some $J_1\in\N$ such that
	\begin{equation}
		\label{eq:Leb_2}
	k\ge J_1 \implies d(x_{n_k},x)< \frac{\varepsilon}{2}.
	\end{equation}
	Now, choose $J_2\in\N$ large enough such that
	\[
	k \ge J_2 \implies \frac{1}{n_k}< \frac{\varepsilon}{2}.
	\]
	We define $K:= \max(J_1,J_2)$. Fix any $y \in B_\frac{1}{n_K}(x_{n_K})$. The triangle inequality together with~\eqref{eq:Leb_2} yields
	\[
	d(x,y) \le d(x,x_{n_K}) + d(x_{n_K},y) < \frac{\varepsilon}{2} +\frac{\varepsilon}{2} = \varepsilon.
	\]
	This implies that $y \in B_\varepsilon(x)\subset U_{\alpha^*}$. Moreover, this shows $B_\frac{1}{n_K}(x_{n_K})\subset U_{\alpha^*}$ which contradicts~\eqref{eq:Leb_1}.
\end{proof}
\begin{prop}
\label{prop:seq_ball_inf}
Let $(X,d)$ be a non-empty metric space with $(x_n)_{n\in\N}\subset X$ and $x\in X$. Suppose that, for every $\varepsilon>0$, the ball $B_\varepsilon(x)$ contains $x_n$ for infinite many $n\in\N$. Then, $(x_n)_{n\in\N}$ has a subsequence converging to $x$.
\end{prop}
\begin{proof}
We construct the subsequence by choosing $\varepsilon = \frac{1}{j}$ for $j\in\N$.
\begin{enumerate}
\item Set $\varepsilon=1$: Since $B_1(x)$ contains $x_n$ for infinitely many $n\in\N$, it must contain $x_{n_1}$ for some $n_1\in\N$.
\item Set $\varepsilon = \frac{1}{2}$: Since $B_\frac{1}{2}(x)$ contains $x_n$ for an infinite number of $n\in\N$, it must contain $x_{n_2}$ for some $n_2 > n_1$.
\item Repeating in this way, for every $j\in\N$, we have that $x_{n_j}\in B_\frac{1}{j}(x)$ and $n_j>n_{j-1}$. 
\end{enumerate}
By construction, we deduce that the subsequence $(x_{n_j})_{j\in\N}\subset (x_n)_{n\in\N}$ converges to $x$.
\end{proof}
Using a combination of the previous results, we are ready to prove~\Cref{thm:cpct_equiv}.
\begin{proof}[Proof of~\Cref{thm:cpct_equiv}]
	\ul{$\implies$:} We assume $X$ is compact and assume, for a contradiction, that $X$ is not sequentially compact. This means that there exists $(x_n)_{n\in\N}\subset X$ which admits no convergent subsequence. The contrapositive of~\Cref{prop:seq_ball_inf} implies that, for every $x\in X$, there exists a corresponding $\varepsilon_x>0$ such that $B_{\varepsilon_x}(x)$ contains $x_n$ for only finitely many $n\in\N$. Of course, we see that $\{B_{\varepsilon_x}(x)\}_{x\in X}$ is an open cover for $X$. By assumption that $X$ is compact, there must exist a finite subcover which we call
	\[
	\{B_{\varepsilon_{x^i}}(x^i)\}_{i=1}^N\subset \{B_{\varepsilon_x}(x)\}_{x\in X},\quad X = \bigcup_{i=1}^NB_{\varepsilon_{x^i}}(x^i).
	\]
However, for every $i=1,\dots, N$, each set $B_{\varepsilon_{x^i}}(x^i)$ must only contain a  $x_n$ for finitely many $n$. This implies that $X$ contains $x_n$ for only finitely many $n\in\N$. This is a contradiction.
%In particular, no term in $(x_n)_{n\in\N}$ may appear infinitely often, otherwise the subsequence consisting of that term would be constant and, hence, convergent. We remove all  Therefore, there exists some $N\in\N$ such that
%	\[
%	\forall i,j\ge N, \, i\neq j \implies x_i \neq x_j.
%	\]
%	If this wasn't the case, then we would be able to construct a constant subsequence (why? Students: fill in details). Let us now consider the sequence $(x_n)_{n\ge N}$ where we have thrown away the first $N$ terms. In this sequence where every term is distinct, for every $n\in \N$, there exists $\varepsilon_n>0$ such that $B_{\varepsilon_n}(x_n)$ contains $x_n$ and no other terms in the sequence. If this wasn't the case, one could find a convergent subsequence (why? Students: fill in details). Let us rephrase this last point in terms of general open sets: for every $\alpha\in \N$ such that $\alpha \ge N$, we have an open set $U_\alpha := B_{\varepsilon_\alpha}(x_\alpha)$ such that $x_j\notin U_\alpha$ for every $j\neq \alpha$.
%	
%	Let us define $U_0:= X\setminus \left(\bigcup_{n \ge N}\{x_n\}\right)$. Now, since $\bigcup_{n\ge N}\{x_n\}$ is closed (why? Be careful since this is an infinite union of singletons.), we see that $U_0$ is open. From the previous discussions, we have thus constructed a collection $\{U_0,U_\alpha\}_{\alpha \ge N}$ of open sets. In particular, we see that
%	\[
%	X = U_0 \cup \bigcup_{\alpha \ge N}U_\alpha.
%	\]
%	We see that $\{U_0,U_\alpha\}_{\alpha \ge N}$ is an open cover of $X$. However, any finite subcollection of this cover will fail to include infinitely many terms in $(x_n)_{n\ge N}$ (why? Students fill in the details). Thus, no finite subcover can be extracted and this is a contradiction.
	
	\ul{$\impliedby$:} Let us now assume that $X$ is sequentially compact and \textcolor{blue}{fix any open cover $\{U_\alpha\}_{\alpha \in A}$ for $X$.} According to~\Cref{lem:Lebesgue_number}, there exists $r>0$ such that
	\[
	\forall x\in X,\exists \alpha \in A\text{ s.t. }B_r(x)\subset U_\alpha.
	\]
	By~\Cref{thm:hb_gen}, we know that $X$ is totally bounded. Thus, there exists a finite number of points $y_1,\dots, y_N\in X$ (depending on $r$) such that
	\[
	X = \bigcup_{i=1}^NB_r(y_i).
	\]
	In particular, for each $i=1,\dots, N$, there exists some associated index $\alpha(i)\in A$ such that $B_r(y_i)\subset U_{\alpha(i)}$. \textcolor{blue}{We extract $\{U_{\alpha(i)}\}_{i=1}^N$ and the previous discussions show
	\[
	X = \bigcup_{i=1}^N B_r(y_i)\subset \bigcup_{i=1}^N U_{\alpha(i)}.
	\]
	}
	The text in \textcolor{blue}{blue} verifies shows that $X$ is compact.
\end{proof}
	\section{Connectedness}
\label{sec:connected}
This section introduces the topological notion of connectedness. Intuitively, a connected space is one which `should not break into two (or more) disjoint parts.'
\begin{defn}[Connected]
	\label{defn:connected}
	A (non-empty) metric space $(X,d)$ is said to be \textit{connected} provided the following holds. If $A,B\subset X$ are open subsets such that $\emptyset \neq A \neq X$ and $\emptyset \neq B \neq X$, then at least one of the following statements:
	\[
	A\cup B = X\quad\text{or }A\cap B = \emptyset
	\]
	must be false. By convention, the empty set is connected.
\end{defn}
In practice, it may be difficult to check the criteria in~\Cref{defn:connected} to determine connectedness. We will come back to this later.
\begin{lemma}
	\label{lem:closed_connect}
	A metric space $(X,d)$ is connected iff the word `open' in~\Cref{defn:connected} is replaced by `closed.'
\end{lemma}
\begin{proof}
	\ul{$\implies$:} Let us assume that $(X,d)$ is connected and \textcolor{blue}{fix any closed sets $C,D\subset X$ such that $C,D\neq \emptyset$ and $C,D \neq X$}. According to~\Cref{defn:open_closed}, this means that the complements $A:= X\setminus C$ and $B:= X\setminus D$ are open and $A,B\neq \emptyset$ and $A,B\neq X$. Now, we use De Morgan's laws (exercise) to obtain
	\[
	A \cup B = (X\setminus C)\cup (X \setminus D) = X \setminus (C\cap D).
	\]
	And similarly, $A\cap B = X\setminus (C \cup D)$. Since $A$ and $B$ are open, we deduce that
	\[
	X = A\cup B \iff C\cap D = \emptyset\quad\text{and}\quad\emptyset = A \cap B \iff C \cup D = X.
	\]
	Thus, we see that \textcolor{blue}{at least one of $C \cup D = X$ or $C\cap D = \emptyset$ must be false.}
	
	\ul{$\impliedby$:} One can repeat the above argument in reverse.
\end{proof}
Recall from the last few points of~\Cref{eg:R_open_closed} that the notion of open and closed sets need not be negations of each other. Sets can be both open and closed (sometimes called `clopen'), exactly one of open or closed, or neither. In a \textit{connected} metric spaces, however, one can remove most of these possibilities.
\begin{prop}
	\label{prop:clopen_connected}
	Let $(X,d)$ be a metric space. TFAE:
	\begin{itemize}
		\item $X$ is connected.
		\item The only subsets of $X$ which are both open and closed are $X$ and $\emptyset$.
	\end{itemize}
\end{prop}
\begin{proof}
	\ul{$\implies$:} Suppose, for a contradiction, that $A\subset X$ is a non-empty set such that $A\neq X$ and $A$ is both open and closed. We can define $B:= X\setminus A$ which is open (and closed) and non-empty. We would then see that both
	\[
	A \cup B = X\quad\text{and}\quad A\cap B = \emptyset,
	\]
	which is a contradiction.
	
	\ul{$\impliedby$:} Let $A,B\subset X$ be non-empty open subsets of $X$. Assume, for a contradiction, that we have both
	\[
	A\cup B = X\quad\text{and}\quad A\cap B = \emptyset.
	\]
	This implies that $B = A^\complement$ and $A = B^\complement$. Thus, both $A$ and $B$ are also closed. This is a contradiction.
\end{proof}
We will have more to say about connected spaces especially in relation to continuous functions. This will be done later. For now, we end this with a close study of connectedness for $\R$ to confirm our intuition. 
\begin{defn}[Connectedness of subspaces]
	\label{defn:connect_sub}
	Let $(X,d)$ be a non-empty metric space and let $A\subset X$ be non-empty. We say that $A$ is \textit{connected} provided that $(A,d|_A)$ is a connected metric space (c.f.~\Cref{eg:sub_met}).
\end{defn}
\subsection{Connectedness on $\R$}
\label{sec:connect_R}
Our intuition on $\R$ about connectedness is very much linked with intervals. Given $a<b$ (we shall allow $a = -\infty$ and / or $b = +\infty$), intervals take on one of the following four forms:
\begin{equation}
	\label{eq:intervals}
	(a,b),\,(a,b],\,[a,b),\,\text{or }[a,b].
\end{equation}
Of course, the sets in~\eqref{eq:intervals} are connected (in the intuitive sense) and~\Cref{thm:connected_interval} below shows that is consistent with connectedness in the sense of~\Cref{defn:connected}. For technical convenience, we shall identify $\emptyset = (a,a)$.
\begin{thm}[On $\R$, connected $\iff$ interval]
	\label{thm:connected_interval}
	$I \subset \R$ is connected iff $I\subset \R$ is an interval in the form of~\eqref{eq:intervals}.
\end{thm}
We will first need to prove some other results.
\begin{prop}
	\label{prop:char_int}
	Let $I\subset \R$ be non-empty. TFAE:
	\begin{itemize}
		\item $I$ is an interval in the form of~\eqref{eq:intervals}.
		\item $I$ satisfies the property:
		\begin{equation}
			\label{eq:char_int}
			\text{Let }x,y\in I.\text{ If }z\in \R\text{ satisfies }x < z < y,\text{ then }z \in I.
		\end{equation}
	\end{itemize}
\end{prop}
\begin{proof}
	\ul{$\implies$:} This is immediate.
	
	\ul{$\impliedby$:} We define
	\[
	\begin{array}{ccc}
		a = \inf I 	&\text{or }-\infty, 	&\text{if $I$ is not bounded below,} 	\\
		b = \sup I &\text{or }+\infty, 	&\text{if $I$ is not bounded above.}
	\end{array}
	\]
	We will show that $(a,b)\subset I \subset [a,b]$ with the convention that $[$ or $]$ adjacent to $\pm \infty$ means the same as $($ or $)$ adjacent to $\pm \infty$. Again, we identify $\emptyset = (a,a)$. We would then be done (why?). By definition of $a,b$ we already have $I\subset [a,b]$ so it remains to show $(a,b) \subset I$.
	
	\textcolor{blue}{Fix $z \in (a,b)$.} By definition, we have $z> a = \inf I$ and / or $z >-\infty$ (if $I$ is not bounded below). In either case, \textcolor{blue}{there exists $x \in I$ such that $x<z$.} Similarly, since $z < b$, \textcolor{blue}{there exists $y\in I$ such that $z < y$.} By assumption that $I$ satisfies~\eqref{eq:char_int}, we conclude that \textcolor{blue}{$z \in I$.}
\end{proof}
The characterisation of intervals in~\Cref{prop:char_int} sets us up to prove~\Cref{thm:connected_interval}.
\begin{proof}[Proof of~\Cref{thm:connected_interval}]
	\ul{$\implies$:} Suppose, for a contradiction, that $I\subset \R$ is connected, but not an interval of the form~\eqref{eq:intervals}. According to~\Cref{prop:char_int}, there must exist $x,y\in I$ and $z \notin I$ such that $x < z < y$. We now define $A:= (-\infty,z) \cap I$ and $B:= (z,\infty) \cap I$. These are \textcolor{blue}{both (relatively) open and non-empty subsets of $I$ with $A,B\neq I$. We also see that both the following are true:
	\[
	A \cup B = I\quad\text{and}\quad A\cap B = \emptyset.
	\]
}
	This is a contradiction to~\Cref{defn:connected}.
	
	\ul{$\impliedby$:} Let $I\subset \R$ be an interval of the form~\eqref{eq:intervals} and suppose, for a contradiction, that $I$ is not connected. More specifically, we assume that there exist non-empty $A,B\subsetneq I$ which are open and closed (relative to $I$) such that
	\[
	A\cup B = I\quad\text{and}\quad A\cap B = \emptyset.
	\]
	Let us take $a\in A, b\in B$ and assume, wlog, that $a<b$ (otherwise we just exchange the roles of $A,B$). Since $a,b\in I$ and $I$ is an interval of the form~\eqref{eq:intervals}, we know by~\Cref{prop:char_int} that $[a,b] \subset I$.
	
	We now define $A' = A\cap [a,b]$ and $B' = B \cap [a,b]$. Note that $a\in A'$ and $b\in B'$. We see that both $A,B$ are (relatively) closed in $I$ and $[a,b]$ is relatively closed in $I$. Since $A',B'$ are intersections of (relatively) closed subsets of $I$, we deduce that $A',B'$ are (relatively) closed in $I$. Since $A',B'\subset [a,b]$, we can use~\Cref{prop:triple_rel} to deduce that $A',B'$ are (relatively) closed in $[a,b]$. Finally, since $[a,b]$ is closed in $\R$, we see (again, by~\Cref{prop:triple_rel}) that $A',B'$ are also closed in $\R$. 
	
The Completeness Axiom gives $c:=\sup A' \in A'$ (why?). Now, $c \le b$ because $c \in A' \subset [a,b]$ and, moreover, $c \neq b$ because $A' \cap B' = \emptyset$ and $b \in B'$. Hence, we deduce $c < b$. Since $A$ is open in $I$, we have that $A'$ is relatively open in $[a,b]$. But this implies that there exists some $\delta>0$ such that $(c-\delta,c+\delta)\cap [a,b] \subset A'$. Since $c < b$, there must exist some points in $(c,c+\delta)\cap [a,b]\subset A'$ which contradicts $c = \sup A'$.
\end{proof}
%	\chapter{Series on $\R$}
%	\label{ch:series_R}
%	\textcolor{ForestGreen}{Some / all of this chapter may have already been taught in Math 2080.} It is intended to cover much more than is necessary in order to deal with series of functions later on.
%	We use the tools we have learned from sequences to study convergence properties of infinite series. In your mathematical careers, you may have come across expressions like
\[
\sum_{n=1}^N a_n = a_1 + a_2 + \dots + a_N.
\]
We study such objects, in particular we investigate
\[
\lim_{N\to\infty} \sum_{n=1}^N a_n.
\]
\section{Summation notation}
Just to fix some notation, we recall that $\sum_{k=m}^na_k$ is just short hand for
\[
\sum_{k=m}^na_k = a_m + a_{m+1}+ \dots + a_n.
\]
The symbol $\sum_{k=m}^na_k$ tells us to add the terms $a_k$ starting from $k=m$ and increasing the value of $k$ by 1 until we reach the index $k=n$. We emphasise that the subscript $k$ in this summation is just a dummy variable. All of the following expressions equal each other
\[
\sum_{\heartsuit=m}^n a_\heartsuit =\sum_{j=m}^n a_j = \sum_{k=m}^na_k = a_m + a_{m+1}+ \dots + a_n.
\]
\begin{eg}
	For $n\in\N$, the symbol $\sum_{k=0}^n2^{-k}$ is just shorthand for
	\[
	\sum_{k=0}^n2^{-k} = 1 + \frac{1}{2} + \frac{1}{4} + \dots + \frac{1}{2^n}.
	\]
\end{eg}
\section{Infinite series}
We want to study \textbf{infinite series} which take the form $\sum_{j=1}^\infty a_j$. To be precise about what this means, we consider the sequence of \textbf{partial sums} given by
\[
s_n = \sum_{j=1}^n a_j = a_0 + a_1 + a_2 + \dots + a_n.
\]
\begin{defn}
	We say that the infinite series $\sum_{j=1}^\infty a_j$ \textbf{converges} if the sequence of partial sums $s_n = \sum_{j=1}^n a_j$ converges to a real number $S$. In this case, we write $\sum_{j=1}^\infty a_j = S$.
\end{defn}
In other words, the infinite series $\sum_{j=1}^\infty a_j$ is just a glorified limit of the sequence of partial sums
\[
\sum_{j=1}^\infty a_j = S \iff \lim_{n\to \infty}\left(\sum_{j=1}^n a_j\right)= S.
\]
A series that does not converge is said to \textbf{diverge}. Often, we will be concerned with properties of infinite series besides their exact values or what the starting index is. In other words, sometimes I will write $\sum a_j$ in place of $\sum_{j=m}^\infty a_j$.

If $(a_n)$ is a sequence of non-negative numbers, then the partial sums $s_n = \sum_{j=1}^n a_j$ form an increasing sequence. We know from Math 2080 that, if $(s_n)_{n\in\N}$ is bounded above, then $s_n \overset{n\to \infty}{\to}\sup_{n\in\N}s_n$.
\begin{defn}
	\label{defn:abscvgseries}
	We say that the infinite series $\sum a_j$ \textbf{converges absolutely} or is \textbf{absolutely convergent} if the infinite series $\sum |a_j|$ converges.
\end{defn}
We will see later on that absolute convergence implies convergence. Moreover, absolutely convergent series are very special and are much nicer than merely convergent series.
\begin{eg}
	[Geometric series]
	\label{eg:geoseries}
	One particular class of series of significance is \textbf{geometric series}. These are expressions of the form $\sum_{n=0}^\infty ar^n$ where $a$ and $r$ are real numbers. When $r\neq 1$ an explicit formula for the partial sums and infinite series is easy to obtain. The partial sums look like
	\begin{equation}
		\label{eq:geopartial}
		\sum_{j=0}^n ar^j = a\frac{1-r^{n+1}}{1-r}.
	\end{equation}
	An easy proof of~\eqref{eq:geopartial} goes as follows.
	\begin{align*}
		(1-r)\sum_{j=0}^n ar^j &= \sum_{j=0}^n ar^j - \sum_{j=0}^n ar^{j+1} 	\\
		&= a + ar + ar^2 + \dots + ar^n \\
		&\quad - (ar + ar^2 + ar^3 + \dots + ar^{n+1}) 	\\
		&= a - ar^{n+1}.
	\end{align*}
	When $|r|<1$, then we have (exercise) $\lim_{n\to\infty} r^{n+1} =0$. Under this assumption, we can pass to the limit $n\to \infty$ in~\eqref{eq:geopartial} to get
	\begin{equation}
		\label{eq:geoinf}
		|r|<1 \implies \sum_{n=0}^\infty ar^n = \frac{a}{1-r}.
	\end{equation}
\end{eg}
\begin{eg}
	We will prove later that
	\begin{equation}
		\label{eq:powerrecip}
		\sum_{n=1}^\infty \frac{1}{n^p} \quad \text{converges if and only if }p>1.
	\end{equation}
	In particular, when $0 \le p \le 1$, we have $\sum n^{-p}=+\infty$. There are exact values for this series when $p$ is even. For example,
	\[
	\sum_{n=1}^\infty \frac{1}{n^2}=\frac{\pi^2}{6}, \quad \sum_{n=1}^\infty \frac{1}{n^4} = \frac{\pi^4}{90}.
	\]
	We will not prove these formulas, instead we will refer to Fourier analysis for one possible technique to verify these formulas.
\end{eg}
We now introduce an important (and hopefully familiar) concept.
\begin{defn}
	[Cauchy criterion]
	\label{defn:cauchysum}
	Let $s_n = \sum_{j=1}^n a_j$ be the sequence of partial sums of $\sum a_j$. We say that the series $\sum a_j$ satisfies the \textbf{Cauchy criterion} if $(s_n)_{n\in\N}$ is a Cauchy sequence.
\end{defn}
In full detail, the Cauchy criterion means
\[
\forall \epsilon>0, \exists N\in\N \text{ such that }n,m\ge N \implies |s_n - s_m | <\epsilon.
\]
The roles of $n,m$ above are completely arbitrary so we may as well assume $n> m$. Moreover, whenever $m>N$ it is certainly true that $m-1\ge N$ so the statement above is equivalent to
\[
\forall \epsilon>0, \exists N\in\N \text{ such that }n\ge m>N \implies |s_n - s_{m-1} | <\epsilon.
\]
Finally, since $s_n- s_{m-1} = \sum_{j=1}^n a_j - \sum_{j=1}^{m-1}a_j = \sum_{j=m}^na_j$, we can rewrite the Cauchy criterion as
\begin{equation}
	\label{eq:cauchysum}
	\forall \epsilon>0, \exists N\in\N \text{ such that }n\ge m>N \implies \left|
	\sum_{j=m}^n a_j
	\right| <\epsilon.
\end{equation}
We will usually use~\eqref{eq:cauchysum} when referring to the \textbf{Cauchy criterion}.
\begin{thm}
	\label{thm:cauchysum}
	A series converges if and only if it satisfies the Cauchy criterion.
\end{thm}
\begin{proof}
	After unpacking the definitions, this follows from the fact that, on $\R$, Cauchy sequences are the same as convergent sequences.
\end{proof}
\begin{cor}
	\label{cor:summandvanish}
	Let $(a_n)_{n\in\N}$ be a sequence of real numbers. The following implication is true.
	\begin{equation}
		\label{eq:summandvanish}
		\sum_{n=1}^\infty a_n \quad \text{converges }\implies \lim_{n\to\infty}a_n = 0.
	\end{equation}
\end{cor}
Before we prove~\Cref{cor:summandvanish} it is important to discuss the direction of the implication in~\eqref{eq:summandvanish}. In particular, I will now point out a common misinterpretation/misuse. \textbf{\Cref{cor:summandvanish} does NOT say `$\lim_{n\to \infty}a_n = 0 \implies \sum_{n=1}^\infty a_n$ converges' because this is, in general, false.} We will see in~\Cref{eg:harmseries} on that $\sum_{n=1}^\infty \frac{1}{n} = +\infty$ even though $\frac{1}{n}\to 0$.
\begin{proof}
	[Proof of~\Cref{cor:summandvanish}]
	Assume $\sum_{n=1}^\infty a_n$ converges to some $S\in\R$. By definition, we have
	\[
	\lim_{n\to \infty} \sum_{j=1}^n a_j = S = \lim_{n\to \infty} \sum_{j=1}^{n-1}a_j.
	\]
	We can write $a_n$ as a difference of these partial sums and pass to the limit $n\to \infty$ to obtain
	\begin{align*}
		a_n 	& = \sum_{j=1}^n a_j - \sum_{j=1}^{n-1}a_j 	\\
		\implies \lim_{n\to \infty} a_n 	&= \lim_{n\to \infty}\left(\sum_{j=1}^n a_j\right) - \lim_{n\to \infty}\left(\sum_{j=1}^{n-1} a_j\right) = S - S = 0.
	\end{align*}
\end{proof}
\begin{cor}
	[Absolute convergence $\implies$ convergence]
	\label{cor:abscvgseries}
	Let $\sum a_n$ be an absolutely convergent series. Then, $\sum a_n$ is also a convergent series.
\end{cor}
\begin{proof}
	\textcolor{blue}{Fix any $\varepsilon>0$.} Now, $\sum a_n$ being absolutely convergent is equivalent to $\sum |a_n|$ being convergent. According to~\Cref{thm:cauchysum}, this is equivalent to $\sum |a_n|$ satisfying the Cauchy criterion. Hence, \textcolor{blue}{there exists some $N\in\N$} such that
	\[
	\forall n\ge m>N,\text{ we have }\left|\sum_{j=m}^n|a_j|\right| <\varepsilon.
	\]
	In particular, \textcolor{blue}{for any $n\ge m>N$}, the triangle inequality gives
	\[
	\textcolor{blue}{\left|\sum_{j=m}^na_j\right|} \le \sum_{j=m}^n |a_j| = \left|\sum_{j=m}^n|a_j|\right| \textcolor{blue}{< \varepsilon.}
	\]
\end{proof}
\section{Tests for convergence}
Next we list some tests for determining whether a series converges or not.
\begin{thm}
	[Comparison test]
	\label{thm:comptest}
	Let $(a_n)$ be a sequence where $a_n \ge 0$ for all $n\in\N$. Let $(b_n)$ be a sequence.
	\begin{enumerate}[label=\textbf{C\arabic*)}]
		\item \label{comp1} Assume $\sum a_n$ converges and $|b_n| \le a_n$ for all $n\in\N$. Then, $\sum b_n$ converges absolutely.
		\item \label{comp2} Assume $\sum a_n = +\infty$ and $b_n \ge a_n$ for every $n\in\N$. Then, $\sum b_n = +\infty$.
	\end{enumerate}
\end{thm}
\begin{rem}
	\ref{comp2} has an analogue when a series diverges to $-\infty$. The precise statement is that if $b_n\le a_n$ for every $n\in\N$ and $\sum a_n = -\infty$, then we have
	\[
	\sum b_n = -\infty.
	\]
\end{rem}
\begin{proof}
	\begin{enumerate}[label=\textbf{C\arabic*)}]
		\item We will show that $\sum |b_n|$ satisfies the Cauchy criterion~\eqref{eq:cauchysum} which, by~\Cref{thm:cauchysum}, is equivalent to proving $\sum |b_n|$ converges. This, of course, means that $\sum b_n$ converges absolutely. Fix $\epsilon>0$ and since $\sum a_n$ converges, it also satisfies the Cauchy criterion. Therefore, there exists $N\in\N$ such that
		\[
		n\ge m > N \implies \left|
		\sum_{j=m}^n a_j
		\right|<\epsilon.
		\]
		However, by definition of the absolute value and the assumption that $|b_j| \le a_j$, we also have
		\[
		\left|
		\sum_{j=m}^n |b_j|
		\right| = \sum_{j=m}^n |b_j| \le \sum_{j=m}^n a_j = \left|
		\sum_{j=m}^n a_j
		\right|.
		\]
		Therefore, the series $\sum |b_j|$ also fulfills the Cauchy criterion. Hence, by~\Cref{thm:cauchysum}, the series $\sum |b_j|$ is convergent which is equivalent to saying that $\sum b_j$ is absolutely convergent.
		\item Denote by $s_n$ and $t_n$ the partial sums of $\sum a_n$ and $\sum b_n$ i.e.
		\[
		s_n = \sum_{j=1}^n a_j, \quad t_n = \sum_{j=1}^n b_j.
		\]
		Since $b_n \ge a_n$, we also get $t_n \ge s_n$ for every $n\in\N$. Since $\lim_{n\to \infty} s_n = \sum_{j=1}^\infty a_j = +\infty$, we also see $\lim_{n\to\infty} t_n = +\infty$.
	\end{enumerate}
\end{proof}
\begin{eg}
	[Divergence of harmonic series]
	\label{eg:harmseries}
	We will use the Comparison Test to prove that $\sum \frac{1}{n}$ diverges. We set $a_n = \frac{1}{n}$ and we compare it to another sequence $b_n$ defined as follows
	\[
	\begin{array}{c|c|c|c|c|c|c|c|c|c|c|c}
		n &1	&2	&3	&4	&5	&6	&7	&8	&9	&10 &\dots	\\\hline
		a_n 	&1	&\frac{1}{2}	&\frac{1}{3}	&\frac{1}{4}	&\frac{1}{5}	&\frac{1}{6}	&\frac{1}{7}	&\frac{1}{8}	&\frac{1}{9}	&\frac{1}{10}  	&\dots	\\\hline
		b_n 	&1	&\frac{1}{2} 	&\frac{1}{4}	&\frac{1}{4}	&\frac{1}{8}	&\frac{1}{8} 	&\frac{1}{8} 	&\frac{1}{8} 	&\frac{1}{16} 	&\frac{1}{16} 	&\dots
	\end{array}
	\]
	So $b_n$ is identically equal to $2^{-n}$ for $2^{n-1}$ consecutive terms when $n\ge 2$. A more precise definition of $b_n$ for $n\ge 2$ goes as follows. Whenever $n\ge 2$, there exists a unique $k\in\N$ such that
	\[
	2^{k-1} < n \le 2^k.
	\]
	Therefore, for fixed $n\ge 2$, take the unique $k\in \N$ satisfying the inequalities above and define 
	\[
	b_1 := 1, \quad b_n := 2^{-k}.
	\]
	By construction (see also the table above), we easily see that $a_n \ge b_n >0$ for every $n\in\N$. The series $\sum b_n$ diverges because, when $N\in\N$ is a large natural number, the partial sums look like
	\begin{align*}
		\sum_{n=1}^N b_n &= 1 + \frac{1}{2} + \left(\frac{1}{4} + \frac{1}{4}\right) + \left(\frac{1}{8}+ \frac{1}{8}+ \frac{1}{8}+ \frac{1}{8}\right) + \dots + b_N 	\\
		&= 1 + \frac{1}{2} + 2*\frac{1}{4} + 4*\frac{1}{8} + \dots + b_N 	\\
		&= 1 + \frac{1}{2} + \frac{1}{2} + \frac{1}{2} + \dots + b_N.
	\end{align*}
	Explicitly, this expression of the partial sums shows that $\sum b_n$ adds the number $\frac{1}{2}$ \textit{infinitely many times.} In particular, $\sum b_n$ diverges to $+\infty$ and by the Comparison Test~\Cref{thm:comptest}, we see that $\sum \frac{1}{n}$ also diverges to $+\infty$.
\end{eg}
The comparison test in~\Cref{thm:comptest} requires knowing a priori that a related series converges. The next few tests we list allow us to calculate directly from $\sum a_n$ whether this series converges.
\begin{thm}
	[Root test]
	\label{thm:rootttest}
	Let $(a_n)_{n\in\N}$ be a sequence and set $\alpha = \limsup_{n\to\infty} |a_n|^\frac{1}{n}$.
	\begin{enumerate}[label=\textbf{Root \arabic*)}]
		\item $\alpha <1 \implies \sum a_n$ converges absolutely.
		\item $\alpha >1 \implies \sum a_n$ diverges.
		\item $\alpha =1$ gives no information.
	\end{enumerate}
\end{thm}
\begin{proof}
	Throughout, we denote $\alpha = \limsup_{n\to\infty} |a_n|^\frac{1}{n}$.
	\begin{enumerate}[label=\textbf{Root \arabic*)}]
		\item Assuming $\alpha<1$, the Archimedean property guarantees the existence of $\epsilon>0$ such that $\alpha + \epsilon<1$ (check). Then, by the definition of limit superior, there exists $N\in\N$ such that
		\[
		\sup\{ |a_n|^\frac{1}{n} \, : \, n \ge N\} < \alpha + \epsilon.
		\]
		In particular, whenever $n\ge N$ we have $|a_n|^\frac{1}{n} < \alpha + \epsilon$ and raising this inequality to the $n$th exponent yields
		\[
		n\ge N \implies |a_n| < (\alpha + \epsilon)^n.
		\]
		Since $0 < \alpha + \epsilon<1$, the geometric series $\sum_{n=N}^\infty (\alpha + \epsilon)^n$ converges absolutely and the comparison test~\Cref{thm:comptest} ensures $\sum_{n=N}^\infty a_n$ converges. We conclude (exercise) that the full series $\sum_{n=1}^\infty a_n = \sum_{n=1}^{N-1}a_n + \sum_{n=N}^\infty a_n$ converges.
		\item We assume that $\alpha>1$. By definition of $\alpha$, this means that
		\[
		1 < \inf_{n\in\N} \sup_{j\ge n} |a_j|^\frac{1}{j}.
		\]
		Therefore, we get
		\[
		1 < \sup_{j\ge n} |a_j|^\frac{1}{j}, \quad \forall n\in\N.
		\]
		This inequality shows (exercise) that there are infinitely many natural numbers $n\in\N$ such that $|a_n|^\frac{1}{n}>1$. Taking $n$th powers, there are an infinite number of natural numbers $n\in\N$ such that $|a_n|>1$. In particular, the sequence $(a_n)$ cannot converge to 0 (exercise), so by (the contrapositive of)~\Cref{cor:summandvanish}, the series $\sum a_n$ diverges.
		\item For both of the series $\sum \frac{1}{n}$ and $\sum \frac{1}{n^2}$, the value of $\alpha$ is identically equal to 1. However, the former series diverges according to~\Cref{eg:harmseries} while the latter converges (we will show this later). Hence, nothing can be concluded in this case.
	\end{enumerate}
\end{proof}
\begin{thm}
	[Ratio test]
	\label{thm:ratiotest}
	Let $(a_n)_{n\in\N}$ be a sequence such that $a_n \neq 0$ for every $n\in\N$.
	\begin{enumerate}[label=\textbf{Ratio \arabic*)}]
		\item $\limsup_{n\to \infty}\left|\frac{a_{n+1}}{a_n}\right| <1 \implies \sum a_n$ converges absolutely.
		\item $\liminf_{n\to \infty}\left|\frac{a_{n+1}}{a_n}\right| >1 \implies \sum a_n$ diverges.
		\item Otherwise, $\liminf_{n\to \infty}\left|\frac{a_{n+1}}{a_n}\right| \le 1 \le \limsup_{n\to \infty}\left|\frac{a_{n+1}}{a_n}\right|$ and no information is given.
	\end{enumerate}
\end{thm}
\begin{proof}
	Following the notation from the Root test~\Cref{thm:rootttest}, we set $\alpha = \limsup_{n\to\infty} |a_n|^\frac{1}{n}$. \textcolor{red}{MAKE REFERENCE TO~\eqref{eq:ratioroot} IN A PROBLEM SHEET OR INCLUDE PROOF}
	\begin{equation}
		\label{eq:ratioroot}
		\liminf_{n\to \infty}\left|\frac{a_{n+1}}{a_n}\right| \le \alpha \le \limsup_{n\to \infty}\left|\frac{a_{n+1}}{a_n}\right|.
	\end{equation}
	\begin{enumerate}[label=\textbf{Ratio \arabic*)}]
		\item If $\limsup_{n\to \infty}\left|\frac{a_{n+1}}{a_n}\right| <1$, then~\eqref{eq:ratioroot} ensures $\alpha <1$, so the Root test~\Cref{thm:rootttest} ensures that $\sum a_n$ converges absolutely.
		\item If $\liminf_{n\to \infty}\left|\frac{a_{n+1}}{a_n}\right|>1$, then~\eqref{eq:ratioroot} ensures $\alpha >1$, so the Root test~\Cref{thm:rootttest} ensures that $\sum a_n$ diverges.
		\item Consider again $\sum \frac{1}{n}$ and $\sum \frac{1}{n^2}$.
	\end{enumerate}
\end{proof}
Before we see how the previous Comparison, Root, and Ratio tests are used in many examples, we point out that the Root and Ratio tests are `blunt' tools. These are blunt tools in the sense that, a series can converge without the first conditions being true. For example, $\sum \frac{1}{n^2}$ is a convergent series however $\alpha = \limsup_{n\to \infty}|a_n|^\frac{1}{n} = 1$. In other words, a convergent series need not satisfy the condition $\alpha<1$.
\section{Examples}
We will see how the Comparison, Root, and Ratio tests can apply (or fail) to some examples. To save some space with notation, we will consider series of the form
\[
\sum a_n
\]
and we define $\alpha := \limsup_{n\to \infty}|a_n|^\frac{1}{n}$.
\begin{eg}
	Consider the series 
	\begin{equation}
		\label{eq:egseries1}
		\sum_{n=2}^\infty \left(-\frac{1}{3}\right)^n = \frac{1}{9} - \frac{1}{27} +\frac{1}{81} - \dots
	\end{equation}
	This is a geometric series so we can apply~\eqref{eq:geoinf} (since $r = -1/3$ verifies $|r|<1$) and so the limit is $1/12$.
	
	Alternatively, the series~\eqref{eq:egseries1} can be shown to converge by all of the previous tests. One could use the comparison test with the series $\sum 3^{-n}$. Moreover (exercise), for $a_n = (-1/3)^n$, the corresponding value of $\alpha$ is $\alpha = 1/3$. Since $\alpha<1$, the Root and Ratio test allow us to conclude that $\sum (-3)^{-n}$ converges.
\end{eg}
\begin{eg}
	Consider the series
	\begin{equation}
		\label{eq:egseries2}
		\sum \frac{n}{n^2+3}.
	\end{equation}
	This series diverges to $+\infty$ by the Comparison test, which we shall prove. Firstly, just to get some intuition, when $n$ takes on very large values, each summand `looks' like $\frac{1}{n}$. We know from~\Cref{eg:harmseries} that the harmonic series diverges so by comparison, the series in~\eqref{eq:egseries2} should also diverge.
	
	The intuitive discussion from the previous discussion is not rigorous so we translate the ideas there into a proper proof. Now, whenever $n\ge 1$, it is certainly true that
	\[
	n^2 + 3 \le n^2 + 3n^2 = 4n^2.
	\]
	Hence, by taking reciprocals, we see that
	\[
	n\ge 1 \implies \frac{n}{n^2 + 3}\ge \frac{n}{4n^2} = \frac{1}{4n}.
	\]
	Therefore, we have $\sum \frac{n}{n^2+3}\ge \frac{1}{4}\sum \frac{1}{n} = +\infty$ so the series in~\eqref{eq:egseries2} diverges to $+\infty$.
	
	Let us see what the Ratio and Root tests can tell us with $a_n = \frac{n}{n^2+3}$ for the series in~\eqref{eq:egseries2}. We have
	\[
	\frac{a_{n+1}}{a_n} = \frac{n+1}{(n+1)^2+3}\frac{n^2+3}{n} =\frac{n+1}{n}\frac{n^2+3}{n^2+2n+4} \to 1\quad \text{as }n\to\infty.
	\]
	Therefore, $\alpha=1$ and neither the Ratio nor Root test would have told us~\eqref{eq:egseries2} diverged.
\end{eg}
\begin{eg}
	Consider the series
	\begin{equation}
		\label{eq:egseries3}
		\sum \frac{n}{3^n}.
	\end{equation}
	Using the Ratio or Root test, we take $a_n = \frac{n}{3^n}$ and calculate
	\[
	\frac{a_{n+1}}{a_n}=\frac{n+1}{3^{n+1}}\frac{3^n}{n} = \frac{n+1}{3n} \to \frac{1}{3}\quad \text{as }n\to\infty.
	\]
	Therefore, by the Ratio test, we conclude~\eqref{eq:egseries3} converges.
\end{eg}
\begin{eg}
	Consider the series
	\begin{equation}
		\label{eq:egseries4}
		\sum\left(
		\frac{2}{(-1)^n-3}
		\right)^n.
	\end{equation}
	We see $a_n = \left(
	\frac{2}{(-1)^n-3}
	\right)^n$ and since it is raised to an $n$th power, it seems like the Root Test could help here. In particular, we see that
	\[
	|a_n|^\frac{1}{n} = \left\{
	\begin{array}{ll}
		1 	&n\text{ even}\\
		\frac{1}{2} 	&n\text{ odd}
	\end{array}
	\right.\implies \alpha = \limsup_{n\to\infty}|a_n|^\frac{1}{n}=1.
	\]
	Unfortunately, since $\alpha =1$, we cannot conclude anything by the Root Test (nor Ratio Test).
	
	However, if we were a bit more clever about the values of $a_n$ depending on the parity of $n$, we would notice
	\[
	a_n = \left\{
	\begin{array}{cl}
		1 	&n\text{ even} 	\\
		-\frac{1}{2^n} 	&n\text{ odd}
	\end{array}
	\right..
	\]
	This implies that the sequence $a_n$ cannot converge to zero, hence by (the contrapositive of)~\Cref{cor:summandvanish}, the series in~\eqref{eq:egseries4} must diverge.
\end{eg}
\begin{eg}
	[Difference between the ratio and root tests]
	Consider the series
	\begin{equation}
		\label{eq:egseries5}
		\sum_{n=0}^\infty 2^{(-1)^n - n}.
	\end{equation}
	\ul{Comparison test:} We set $a_n = 2^{(-1)^n - n}$ and it is easy to see that $a_n \le 2^{-n+1}$ so by comparing~\eqref{eq:egseries5} with the appropriate geometric series, we conclude that~\eqref{eq:egseries5} converges.
	
	\ul{Ratio test:} We see that
	\[
	\frac{a_{n+1}}{a_n} = \frac{2^{(-1)^{n+1}-(n+1)}}{2^{(-1)^n - n}} = 2^{-1-2(-1)^n} = \left\{
	\begin{array}{cc}
		\frac{1}{8} 	&n\text{ even} 	\\
		2 	&n\text{ odd}
	\end{array}
	\right..
	\]
	Therefore, we can only conclude
	\[
	\frac{1}{8} = \liminf_{n\to\infty}\left|
	\frac{a_{n+1}}{a_n}
	\right|< 1< \limsup_{n\to\infty}\left|
	\frac{a_{n+1}}{a_n}
	\right| = 2,
	\]
	and this gives us no information according to the Ratio test.
	
	\ul{Root test:} We see that
	\[
	|a_n|^\frac{1}{n} = 2^{\frac{(-1)^n}{n}-1} = \left\{
	\begin{array}{cc}
		2^{\frac{1}{n}-1} 	&n\text{ even} 	\\
		2^{-\frac{1}{n}-1}  	&n\text{ odd}
	\end{array}
	\right..
	\]
	Regardless of the parity of $n\in\N$, we nevertheless see $\lim_{n\to\infty} |a_n|^\frac{1}{n} = \alpha = \limsup_{n\to\infty}|a_n|^\frac{1}{n} = \frac{1}{2}<1$. Hence, by the Root test, we can conclude that~\eqref{eq:egseries5} converges.
\end{eg}
\section{Riemann's rearrangement theorem}
\label{sec:rearrange}
In this section, we prove a seemingly surprising result by Riemann.
\begin{thm}[Riemann's rearrangement theorem]
	\label{thm:rearrange}
	Let $\sum a_n$ be a convergent, but not absolutely convergent series.
\end{thm}
	\chapter{Continuous functions on metric spaces}
	\label{ch:cts_met}
	This chapter follows up from~\Cref{ch:metric_spaces} by generalising the notion of continuity for functions defined on and with values in metric spaces. Recall that $f:\R\to \R$ is said to be continuous at $x_0\in\R$ iff the following is true:
	\[
	\forall \varepsilon>0,\exists \delta>0\text{ s.t. }\forall x\text{ satisfying }|x-x_0|<\delta,\text{ we have }|f(x)-f(x_0)|<\varepsilon.
	\]
	In~\Cref{sec:def_met}, we saw how we can view $\R$ as a metric space with
	\[
	d(x,y) = |x-y|.
	\]
	Of course, we can now generalise the absolute value terms $|x-x_0|$ and $|f(x) - f(x_0)|$ with metrics. Moreover, we can consider functions $f:X \to Y$ where $X,Y$ are two (possibly different) metric spaces.
	%\section{Continuous functions}
%\label{sec:cts_fun}
\begin{defn}[Continuous functions]
\label{defn:cty}
Let $(X,d_X)$ and $(Y,d_Y)$ be metric spaces and $f:X\to Y$ be a function.
\begin{itemize}
	\item For $x_0\in X$, we say that $f$ is \textit{continuous at $x_0$} iff
	\[
	\forall \varepsilon>0,\exists \delta>0\text{ s.t. }\forall x\in X\text{ satisfying }d_X(x,x_0)<\delta,\text{ we have }d_Y(f(x),f(x_0))<\varepsilon.
	\]
	\item We saw that $f$ is \textit{continuous (on $X$)} iff $f$ is continuous at every $x_0\in X$.
\end{itemize}
\end{defn}
\begin{ex}
Suppose $f:X\to Y$ and $K\subset X$. Show that the restriction $f|_K:K\to Y$ of $f$ to $K$ is also continuous.
\end{ex}
\begin{defn}[Image and preimage]
	\label{defn:preimage}
	Given a function $f:X\to Y$ and subsets $A\subset X, B \subset Y$, we denote the \textit{image of $A$ through $f$} and the \textit{preimage of $B$ under $f$} by
	\[
	f(A) = \left\{
	f(x) \, : \, x \in A
	\right\}\quad\text{and }f^{-1}(B) := \left\{
	x\in X\, : \, f(x) \in B
	\right\}.
	\]
\end{defn}
\begin{ex}
	\label{ex:cts_image}
	Let $(X,d_X)$ and $(Y,d_Y)$ be metric spaces with $a\in X$ and a function $f:X\to Y$. Prove that TFAE:
	\begin{itemize}
		\item $f$ is continuous at $a$.
		\item For every $\varepsilon>0$, there exists $\delta>0$ such that
		\[
		f(B_\delta(a))\subset B_\varepsilon(f(a)).
		\]
		\item For every $\varepsilon>0$, there exists $\delta>0$ such that
		\[
		B_\delta(a) \subset f^{-1}(B_\varepsilon(f(a))).
		\]
	\end{itemize}
\end{ex}
The third point from~\Cref{ex:cts_image} can be generalised to open sets.
\begin{prop}[Preimage of continuous functions]
	\label{prop:cts_pre_open}
	Let $(X,d_X)$ and $(Y,d_Y)$ be metric spaces with $f:X\to Y$. TFAE:
	\begin{enumerate}
		\item $f$ is continuous.
		\item For every open set $A\subset Y$, the preimage $f^{-1}(A)$ is open (as a subset of $X$).
	\end{enumerate}
\end{prop}
\begin{proof}
	\ul{$\implies$:} \textcolor{blue}{Fix any open set $A\subset Y$ and element $a \in f^{-1}(A)$.} We wish to show that we can find $\delta>0$ such that $B_\delta(a) \subset f^{-1}(A)$. From the text in \textcolor{blue}{blue}, we know that there exists $\varepsilon>0$ such that
	\[
	B_\varepsilon(f(a))\subset A.
	\]
	By assumption that $f$ is continuous, \textcolor{blue}{there exists $\delta>0$} such that
	\[
	d_X(x,a)<\delta \implies d_Y(f(x),f(a))<\varepsilon.
	\]
	In terms of balls, this means
	\[
	x\in B_\delta(a)\implies f(x) \in B_\varepsilon(f(a)).
	\]
	\textcolor{blue}{Fix any $x\in B_\delta(a)$.} Since $B_\varepsilon(f(a))\subset A$, the above implies that $f(x) \in A$. In terms of balls, this means that \textcolor{blue}{$x \in f^{-1}(A)$.}
	
	\ul{$\impliedby$:} \textcolor{blue}{Fix any $\varepsilon>0$ and $a\in X$.} Motivated by the converse direction, we define the set $A:= B_\varepsilon(f(a))\subset Y$. Recall from~\Cref{eg:balls_are_open} that balls are open so $A$ is certainly open. By assumption, we obtain that $f^{-1}(A)$ is an open subset of $X$. Notice that $a\in f^{-1}(A)$. Hence, \textcolor{blue}{there exists some $\delta>0$} such that
	\[
	B_\delta(a) \subset f^{-1}(A).
	\]
	This means that \textcolor{blue}{for any $x\in X$ with $d_X(x,a)<\delta$}, we have $x \in f^{-1}(B_\varepsilon(f(a)))$. Specifically, this means that $f(x) \in B_\varepsilon(f(a))$ which is the same as \textcolor{blue}{$d_Y(f(x),f(a))<\varepsilon$.}
\end{proof}
\begin{rem}
	In the general setting of \textit{topological spaces}, the equivalence in~\Cref{prop:cts_pre_open} is actually used as the definition of continuous functions.
\end{rem}
\begin{ex}
	Prove~\Cref{prop:cts_pre_open} when the word `open' is replaced with `closed.'
\end{ex}
\begin{ex}
	\label{ex:cts_im}
	Let $f:X\to Y$ be continuous. Is it true that if $A\subset X$ is open, then $f(A)\subset Y$ is open? Answer the same question with the word `open' replaced by `closed.'
\end{ex}
	\section{Stronger notions of continuity}
\label{sec:strong_cts}
In this brief section, we introduce various different stronger notions of continuity. We leave as informative exercises the comparison between them.
\begin{defn}[Uniform continuity]
	\label{defn:unif_cty}
	A function $f:X\to Y$ between two metric spaces $(X,d_X)$ and $(Y,d_Y)$ is said to be \textit{uniformly continuous} provided the following holds true:
	\[
		\forall \varepsilon>0,\exists \delta>0\text{ s.t. }\forall a,b\in X, \, d_X(a,b) < \delta \implies d_Y(f(a),f(b))<\varepsilon.
	\]
\end{defn}
\begin{defn}[H\"older and Lipschitz continuity]
	\label{defn:holder}
	Let $(X,d_X)$ and $(Y,d_Y)$ be two metric spaces. Let $\alpha \in (0,1]$. A function $f: X\to Y$ is said to be \textit{H\"older continuous with exponent $\alpha$} provided there exists a constant $C>0$ such that
	\[
	d_Y(f(a),f(b))\le C\left(d_X(a,b)\right)^\alpha,\quad \forall a,b\in X.
	\]
	When $\alpha =1$, we refer to such functions as \textit{Lipschitz continuous}.
\end{defn}
\begin{lemma}[H\"older $\implies$ uniform continuity $\implies$ continuity]
	\label{lem:cty_hierarchy}
	Let $(X,d_X)$ and $(Y,d_Y)$ be two metric spaces. Let $\alpha \in (0,1]$ and $f:X\to Y$.
	\begin{enumerate}
		\item \label{unif_imp_vanilla} If $f$ is uniformly continuous, then $f$ is continuous.
		\item \label{hold_imp_unif} If $f$ is H\"older continuous with exponent $\alpha$, then $f$ is uniformly continuous.
	\end{enumerate}
\end{lemma}
\begin{proof}
	\ref{unif_imp_vanilla}: \textcolor{blue}{Fix any $\varepsilon>0$ and $a\in X$.} Since $f$ is uniformly continuous, \textcolor{blue}{there exists $\delta>0$} such that any $x_1,x_2\in X$ satisfying $d_X(x_1,x_2)<\delta$ implies $d_Y(f(x_1),f(x_2))<\varepsilon$. In particular, \textcolor{blue}{for any $x\in X$}, set $x_1 = x$ and $x_2 = a$ to see that
	{\color{blue}
	\[
	d_X(x,a)<\delta \implies d_Y(f(x),f(a))<\varepsilon.
	\]
	}
	This shows that $f$ is continuous at $a\in X$. Since $a\in X$ was arbitrary, we see that $f$ is continuous.
	
	\ref{hold_imp_unif}:\textcolor{blue}{Fix any $\varepsilon>0$.} Since $f$ is H\"older continuous, there exists $C>0$ such that
	\[
	d_Y(f(a),f(b))\le C\left(d_X(a,b)\right)^\alpha,\quad \forall a,b\in X.
	\]
	\textcolor{blue}{We define $\delta = \left(\frac{\varepsilon}{C}\right)^\frac{1}{\alpha}>0$. Then, for any $a,b\in X$ satisfying $d_X(a,b)<\delta$, we have
	}
	\[
	\textcolor{blue}{d_Y(f(a),f(b))}\le C\left(d_X(a,b)\right)^\alpha \textcolor{blue}{<}C\left(\left(\frac{\varepsilon}{C}\right)^\frac{1}{\alpha}\right)^\alpha = \textcolor{blue}{\varepsilon.}
	\]
\end{proof}
\begin{ex}
	\label{ex:cts_counter}
	Can either implications in~\Cref{lem:cty_hierarchy} be reversed? In other words, does continuity imply uniform continuity? Does uniform continuity imply H\"older continuity?
\end{ex}
	\section{Basic properties of continuous functions}
\label{sec:basic_prop_cts}
\begin{prop}[Composition of continuous functions is continuous]
	\label{prop:comp_cts}
	Let $(X,d_X),(Y,d_Y),(Z,d_Z)$ be metric spaces. Let $f:X\to Y$ and $g:Y \to Z$ be functions. Let $a\in X$ and assume $f$ is continuous at $a$ and $g$ is continuous at $f(a)$. Then, the composition $g\circ f$ is continuous at $a$.
\end{prop}
\begin{proof}
	\textcolor{blue}{Fix $\varepsilon>0$.} Since $g$ is continuous at $f(a)$ then according to~\Cref{defn:cty}, there must exist some $\delta_1>0$ such that
	\[
	\forall y \in Y \text{ satisfying }d_Y(y,f(a))<\delta_1,\text{ we have }d_Z(g(y),g(f(a)))<\varepsilon.
	\]
	For this $\delta_1>0$, according to~\Cref{defn:cty}, \textcolor{blue}{there must exist some $\delta>0$} such that
	\[
	\forall x \in X\text{ satisfying }d_X(x,a) < \delta,\text{ we have }d_Y(f(x),f(a))<\delta_1.
	\]
	Putting together the previous implications in particular yields
	\[
	\textcolor{blue}{\forall x \in X\text{ satisfying }d_X(x,a)<\delta,\text{ we have }d_Z(g(f(x)),g(f(a)))<\varepsilon.}
	\]
\end{proof}
The next few results will involve product metric spaces. Given metric spaces $(X,d_X)$ and $(Y,d_Y)$, we will endow the product $X\times Y$ with the following metric
\[
d_{X\times Y}((x_1,y_1),\,(x_2,y_2)):= d_X(x_1,x_2) + d_Y(y_1,y_2).
\]
\begin{prop}[Tensor product of continuous functions is continuous]
	\label{prop:tens_prod_cts}
	Let $(X,d_X),(X'd_{X'}),(Y,d_Y)$, and $(Y',d_{Y'})$ be metric spaces with $a\in X$ and $b\in Y$. Let $f:X\to X'$ and $g:Y \to Y'$ be continuous at $a$ and $b$, respectively. Then, the tensor product map $f\times g :X\times Y \to X'\times Y'$ defined by
	\[
	(f\times g)(x,y) := (f(x),\,g(y)),\quad x\in X,\,y \in Y,
	\]
	is continuous at $(a,b) \in X\times Y$.
\end{prop}
\begin{proof}
	\textcolor{blue}{Fix $\varepsilon>0$.} Since $f$ and $g$ are continuous at $a$ and $b$, respectively, we can find $\delta_1,\delta_2>0$ such that
	\[
	d_X(x,a)<\delta_1\implies d_{X'}(f(x),f(a))<\frac{\varepsilon}{2}\quad\text{and }d_Y(y,b)<\delta_2\implies d_{Y'}(g(y),g(b))<\frac{\varepsilon}{2}.
	\]
	\textcolor{blue}{Define $\delta=\min(\delta_1,\delta_2)$ and fix any $(x,y) \in X\times Y$ such that $d_{X\times Y}((x,y),\,(a,b))<\delta$.} We have
	\[
	d_{X\times Y}((x,y),\,(a,b)) = d_X(x,a) + d_Y(y,b)<\delta \implies d_X(x,a)<\delta_1\text{ and }d_Y(y,b) < \delta_2.
	\] 
	Hence, we get
	\begin{align*}
		&\quad \textcolor{blue}{d_{X'\times Y'}((f(x),g(y)),(f(a),g(b)))} 	\\
		&= d_{X'}(f(x),f(a)) + d_{Y'}(g(y),g(b)) 	\\
		&\textcolor{blue}{<} \frac{\varepsilon}{2} + \frac{\varepsilon}{2} = \textcolor{blue}{\varepsilon.}
	\end{align*}
\end{proof}
Given metric spaces $(X,d_X)$ and $(Y,d_Y)$, we define the projections
\[
p_X:(x,y)\in X\times Y \mapsto x \in X,\quad p_Y:(x,y)\in X\times Y \mapsto y\in Y.
\]
\begin{prop}[Projections are continuous]
	The maps $p_X$ and $p_Y$ are continuous.
\end{prop}
\begin{proof}
	\textcolor{blue}{Fix $\varepsilon>0$ and any $(a,b)\in X\times Y$. We define $\delta = \varepsilon$. Fix any $(x,y)\in X\times Y$ such that $d_{X\times Y}((x,y),(a,b))<\delta$.} We see that
	\[
	\textcolor{blue}{d_X(p_X(x,y),p_X(a,b))} = d_X(x,a) \le d_{X\times Y}((x,y),(a,b)) \textcolor{blue}{<\delta = \varepsilon.}
	\]
	A nearly identical argument shows that $p_Y$ is continuous.
\end{proof}
\begin{defn}
	\label{defn:diag_map}
	Given a metric space $(X,d_X)$, the \textit{diagonal map} is the function
	\[
	\Delta : x\in X \mapsto (x,x)\in X\times X.
	\]
\end{defn}
\begin{prop}[Diagonal map is continuous]
	For any metric space $(X,d_X)$, the diagonal map in~\Cref{defn:diag_map} is continuous.
\end{prop}
\begin{proof}
	\textcolor{blue}{Fix any $\varepsilon>0$ and any $a\in X$. We define $\delta = \frac{\varepsilon}{2}$. For every $x\in X$ such that $d_X(x,a) < \delta$, we have the following.}
	\[
	\textcolor{blue}{d_{X\times X}(\Delta(x),\Delta(a))} = d_{X\times X}((x,x),(a,a)) = d_X(x,a) + d_X(x,a) \textcolor{blue}{< 2\delta = \varepsilon.}
	\]
\end{proof}
We conclude this section with an introductory connection between continuity and some of the metric space concepts we've learned previously. In particular, the following result should be very familiar. I encourage you to compare the proof of the result below with what you've seen previously for functions $f:\R \to \R$.
\begin{prop}[Continuity and convergent sequences]
	\label{prop:cty_cvg_seq}
	Let $(X,d_X)$ and $(Y,d_Y)$ be metric spaces with $a \in X$ and a function $f:X\to Y$. TFAE:
	\begin{enumerate}
		\item \label{cty_cvg_seq_1} $f$ is continuous at $a$.
		\item \label{cty_cvg_seq_2} If $(x_n)_{n\in\N}\subset X$ is a sequence converging to $a$, then $(f(x_n))_{n\in\N}\subset Y$ converges to $f(a)$.
	\end{enumerate}
\end{prop}
\begin{rem}
	\Cref{prop:cty_cvg_seq} is not true in the more general setting of \textit{topological spaces}.
\end{rem}
\begin{proof}
	\ul{$\implies$:} \textcolor{blue}{Fix any sequence $(x_n)_{n\in\N}\subset X$ such that $x_n \to a$.} We wish to show that $f(x_n)\to f(a)$. \textcolor{blue}{Fix any $\varepsilon>0$.} By assumption that $f$ is continuous at $a$, we know that there exists some $\delta>0$ such that
	\[
	d_X(x,a) < \delta \implies d_Y(f(x),f(a))<\varepsilon.
	\]
	For this $\delta>0$, since $x_n \to a$, \textcolor{blue}{there exists $N\in\N$} such that
	\[
	n\ge N \implies d_X(x_n,a)<\delta.
	\]
	In particular, we see that
	\[
	\textcolor{blue}{n\ge N\implies} d_X(x_n,a)<\delta \implies \textcolor{blue}{d_Y(f(x_n),f(a))<\varepsilon}.
	\]
	\ul{$\impliedby$:} We will sketch the proof by contradiction. Suppose, for a contradiction, that $f$ is not continuous at $a$. This means that there exists some $\varepsilon>0$ such that
	\begin{equation}
		\label{eq:cty_cvg_seq_cont}
		\forall \delta>0,\exists x\in X\text{ s.t. }d_X(x,a)<\delta\text{ but }d_Y(f(x),f(a))\ge \varepsilon.
	\end{equation}
	In particular, for every $n\in\N$, we apply~\eqref{eq:cty_cvg_seq_cont} with $\delta = \frac{1}{n}$. This gives a sequence $(x_n)_{n\in\N}\subset X$ such that
	\[
	d_X(x_n,a)<\frac{1}{n}\text{ but }d_Y(f(x_n),f(a))\ge \varepsilon,\quad \forall n\in\N.
	\]
	By construction, we see that $x_n\to a$ since $d_X(x_n,a) \to 0$ (c.f.~\Cref{ex:cvg_metric}). By assumption, this implies that $f(x_n)\to f(a)$. There must exist some $N\in\N$ such that
	\[
	d_Y(f(x_N),f(a))<\varepsilon.
	\]
	This is the desired contradiction.
\end{proof}
	\section{Continuous functions on compact spaces}
\label{sec:cpct_cts}
In this (short) section, we prove some results for continuous functions on compact metric spaces.
\begin{prop}[Continuous image of compact sets is compact]
	\label{prop:cts_image_cpct}
	Let $(X,d_X)$ and $(Y,d_Y)$ be metric spaces with $f:X\to Y$ continuous. If $K\subset X$ is compact, then the image $f(K)$ is compact.
\end{prop}
\begin{rem}
	\begin{itemize}
		\item Compare~\Cref{prop:cts_image_cpct} with~\Cref{ex:cts_im}.
		\item The proof of~\Cref{prop:cts_image_cpct} below is deliberately written without explicit reference to any metrics.
	\end{itemize}
\end{rem}
\begin{proof}[Proof of~\Cref{prop:cts_image_cpct}]
	We will use the characterisation of continuous functions in~\Cref{prop:cts_pre_open} for this proof. \textcolor{blue}{Let $\{V_\alpha\}_{\alpha\in A}$ be an open cover of $f(K)$.} We have
	\[
	f(K) \subset \bigcup_{\alpha \in A}V_\alpha \iff K \subset f^{-1}\left(\bigcup_{\alpha \in A}V_\alpha\right) \iff K \subset \bigcup_{\alpha \in A}f^{-1}(V_\alpha).
	\]
	According to~\Cref{prop:cts_pre_open}, since each $V_\alpha$ is open, we have that every $f^{-1}(V_\alpha)$ is also open. We see that $\{f^{-1}(V_\alpha)\}_{\alpha \in A}$ is therefore an open cover of $K$. Since $K$ is compact, \textcolor{blue}{there must exist a finite subset $F\subset A$} such that $\{f^{-1}(V_\alpha)\}_{\alpha \in F}$ covers $K$. Following the same equivalences above, we obtain
	\[
	K \subset \bigcup_{\alpha \in F}f^{-1}(V_\alpha) \iff K \subset f^{-1}\left(\bigcup_{\alpha \in F}V_\alpha\right) \iff \textcolor{blue}{f(K) \subset \bigcup_{\alpha \in F}V_\alpha}.
	\]
	The third inclusion demonstrates that \textcolor{blue}{$\{V_\alpha\}_{\alpha \in F}$ is an open over of $f(K)$.}
\end{proof}
An application of~\Cref{prop:cts_image_cpct} is the Extreme Value Theorem.
\begin{thm}[Extreme Value Theorem]
	\label{thm:EVT}
	Let $(X,d)$ be a non-empty compact metric space and $f:X\to \R$ be continuous. Then, there exist (possibly the same) $x_m,x_M\in X$ such that
	\[
	f(x_m)= \min_{x\in X}f(x)\quad\text{and }f(x_M)=\max_{x\in X}f(x).
	\]
\end{thm}
\begin{proof}
	According to~\Cref{prop:cts_image_cpct}, we know that $f(X)$ is a non-empty compact subset of $\R$. By~\Cref{thm:hb_real} and~\Cref{thm:cpct_equiv}, $f(X)$ is also a closed and bounded subset of $\R$. Finally, \Cref{ex:r_closed_bdd} completes the proof.
\end{proof}
\begin{prop}
	\label{prop:cpct+cts_imp_unif}
	Let $(X,d_X)$ and $(Y,d_Y)$ be two metric spaces and $f:X\to Y$ a function. Assume that $X$ is non-empty and compact and $f$ is continuous. Then, $f$ is uniformly continuous.
\end{prop}
\begin{rem}
	Compare~\Cref{prop:cpct+cts_imp_unif} with~\Cref{ex:cts_counter} and~\Cref{lem:cty_hierarchy}.
\end{rem}
\begin{proof}[Proof of~\Cref{prop:cpct+cts_imp_unif}]
	\textcolor{blue}{Fix any $\varepsilon>0$.} Since $f$ is continuous, we know that, for every $a\in X$, there exists $\delta_a>0$ (emphasis on the dependence of $\delta$ with respect to $a$) such that 
	\[
	d_X(x,a) < \delta_a \implies d_Y(f(x),f(a))<\frac{\varepsilon}{2}.
	\]
	The collection of balls $\{B_{\frac{\delta_a}{2}}(a)\}_{a\in X}$ forms an open cover of $X$. By assumption that $X$ is compact, there exists $N\in\N$ and finitely many $\{a_i\}_{i=1}^N \subset X$ such that
	\[
	X = \bigcup_{i=1}^N B_{\frac{\delta_{a_i}}{2}}(a_i).
	\]
	\textcolor{blue}{Define $\delta := \frac{1}{2}\min_{i=1,\dots, N}\delta_{a_i}>0$. Fix any $x_1,x_2\in X$ such that $d_X(x_1,x_2)<\delta$.} There is some index $j=1,\dots, N$ such that $x_1 \in B_{\frac{\delta_{a_j}}{2}}(a_j)$. Moreover, the triangle inequality gives
	\[
	d_X(x_2,a_j)\le d_X(x_2,x_1) + d_X(x_1,a_j) < \delta + \frac{\delta_{a_j}}{2} \le \frac{\delta_{a_j}}{2} + \frac{\delta_{a_j}}{2} = \delta_{a_j}.
	\]
	Hence, we have $x_2 \in B_{\delta_{a_j}}(a_j)$. Using the triangle inequality again, we have
	\[
	\textcolor{blue}{d_Y(f(x_1),f(x_2))}\le d_Y(f(x_1),f(a_j)) + d_Y(f(a_j),f(x_2)) \textcolor{blue}{<}\frac{\varepsilon}{2} + \frac{\varepsilon}{2} = \textcolor{blue}{\varepsilon.}
	\]
\end{proof}
	\section{Continuous functions on complete spaces}
\label{sec:comp_cts}
A significant result in Analysis is~\Cref{thm:banach} below.
\begin{thm}[Banach's fixed point / Contraction Mapping theorem]
	\label{thm:banach}
	Let $(X,d)$ be a non-empty complete metric space. Let $f:X\to X$ and $\kappa \in (0,1)$ be such that
	\begin{equation}
		\label{eq:contract}
		d(f(a),f(b))\le \kappa d(a,b),\quad \forall a,b\in X.
	\end{equation}
	Then, there exists a unique element $p\in X$ such that $f(p) = p$.
\end{thm}
\begin{rem}
	\begin{itemize}
		\item A function satisfying~\eqref{eq:contract} is sometimes called a \textit{(strict) contraction.}
		\item A point satisfying $p = f(p)$ is called a \textit{fixed point of $f$}.
		\item The assumption of~\eqref{eq:contract} \textbf{cannot} be weakened (exercise) to merely
		\[
		d(f(a),f(b))<d(a,b),\quad \forall a,b\in X.
		\]
	\end{itemize}
\end{rem}
\begin{proof}[Proof of~\Cref{thm:banach}]
	\ul{Uniqueness:} Let us first prove that, if fixed points exist, they must be unique. Let $p_1,p_2\in X$ be such that $f(p_i)=p_i$ for $i=1,2$. By~\eqref{eq:contract}, we have
	\[
	d(p_1,p_2) = d(f(p_1),f(p_2)) \le \kappa d(p_1,p_2)\implies (1-\kappa)d(p_1,p_2)\le 0.
	\]
	Since $\kappa<1$, this can only happen when $d(p_1,p_2) = 0$ which implies, by~\ref{m:pos}, that $p_1=p_2$.
	
	\ul{Existence:} Since $X$ is non-empty, it contains an element which we call $x_0 \in X$. We define a sequence $(x_n)_{n\in\N}$ by iteratively applying $f$:
	\[
	x_n := f(x_{n-1}),\quad \forall n\in\N.
	\]
	Suppose $n,m\in \N$ and $n<m$. By the triangle inequality and inductively applying~\eqref{eq:contract}, we have
	\begin{align}
		\label{eq:contract_estimate}
		\begin{split}
		d(x_n,x_m) &\le \sum_{j=n}^{m-1}d(x_j,x_{j+1}) = \sum_{j=n}^{m-1}d(f(x_{j-1}),f(x_j)) 	\\
		&\le \sum_{j=n}^{m-1}\kappa^{j-1}d(x_0,x_1) < d(x_0,x_1)\sum_{j=n}^\infty \kappa^{j-1} = \frac{d(x_0,x_1)}{1-\kappa}\kappa^{n-1}.
		\end{split}
	\end{align}
	In~\eqref{eq:contract_estimate} above, we used the geometric series with $\kappa\in(0,1)$. This estimate shows that $(x_n)_{n\in\N}$ is Cauchy. Indeed, fix any $\varepsilon>0$. Since $\kappa^{n-1}\overset{n\to \infty}{\to}0$, we can find $N\in\N$ large enough such that
	\[
	n\ge N \implies \kappa^{n-1} < \frac{1-\kappa}{1 + d(x_0,x_1)}\varepsilon.
	\]
	Then, for any $m>n\ge N$, we see from the above and~\eqref{eq:contract_estimate} that $d(x_n,x_m)<\varepsilon$. Since $X$ is complete, we deduce that $(x_n)_{n\in\N}$ is convergent to \textcolor{blue}{some limit $p\in X$}. Going back to the construction of the sequence, we pass to the limit to obtain
	\[
	x_n = f(x_{n-1})\implies \lim_{n\to \infty} x_n = \lim_{n\to \infty}f(x_{n-1}) \implies p = f\left(\lim_{n\to \infty}x_{n-1}\right) \iff \textcolor{blue}{p = f(p)}.
	\]
	The second implication $\lim_{n\to \infty}f(x_{n-1}) = f\left(\lim_{n\to \infty}x_{n-1}\right)$ comes from the fact that $f$ is continuous (why? Hint: \Cref{prop:cty_cvg_seq}). 
\end{proof}
\Cref{thm:banach} is a very important result with many applications. We shall investigate one of those applications later on in the context of solving differential equations.
	\section{Continuous functions and connected spaces}
\label{sec:connected_cts}
This section discusses the relation between connected spaces and continuous functions. In this discussion, we will refer to $\{0,1\}$ as the \textit{two-point space} and we will endow it with the discrete metric. In this setting, the open sets of the two-point space are
\[
\{0,1\},\,\{1\},\,\{0\},\text{ and }\emptyset.
\]
Recall also that we say $f:X\to Y$ is a \textit{surjection} provided
\[
\forall y\in Y,\exists x\in X\text{ s.t. }f(x) = y.
\]
\begin{prop}
	\label{prop:connected_two-point}
	Let $(X,d)$ be a metric space. TFAE:
	\begin{itemize}
		\item $X$ is connected.
		\item There are no continuous surjections from $X$ to the two-point space.
	\end{itemize}
\end{prop}
This result formally says that a connected space is one that `cannot be continuously split into two distinct parts.'
\begin{proof}
	\ul{$\implies$:} Assume, for a contradiction, that $X$ is connected, but there exists a continuous surjection $f:X \to \{0,1\}$. \textcolor{blue}{Define $A:= f^{-1}(\{0\})$ and $B:= f^{-1}(\{1\})$.} According to~\Cref{prop:cts_pre_open}, we see that $A,B \subset X$ are \textcolor{blue}{non-empty and open and $A,B\neq X$}. However, since $f^{-1}(C)\cup f^{-1}(D) = f^{-1}(C\cup D)$ and similarly for intersections, we deduce \textcolor{blue}{both
	\[
	A \cup B = X\quad\text{and}\quad A\cap B = \emptyset.
	\]
	}
	\ul{$\impliedby$:} Assume, for a contradiction, that $X$ is not connected. This means that we can find non-empty $A,B\subsetneq X$ which are open such that
	\[
	A\cup B = X\quad\text{and}\quad A\cap B = \emptyset.
	\]
	We define $f:X\to \{0,1\}$ by
	\[
	f:=\left\{
	\begin{array}{cc}
		0,	&\text{if }x\in A 	\\
		1, 	&\text{if }x \in B
	\end{array}
	\right..
	\]
	Clearly, we see that $f$ is a surjection. As we mentioned before this result, all the open sets of the two-point space are
	\[
	\{0,1\},\,\{1\},\,\{0\},\text{ and }\emptyset.
	\]
	Their respective pre-images through $f$ are
	\[
	X,\,A,\,B,\text{ and }\emptyset,
	\]
	which are all open sets. Hence $f$ is continuous, by~\Cref{prop:cts_pre_open}, which is a contradiction.
\end{proof}
\begin{ex}
	\begin{itemize}
		\item Let $X$ be a non-empty finite set with at least two elements endowed with the discrete metric. Prove that $X$ cannot be connected.
		\item Let $X$ be a countable set endowed with the discrete metric. Prove that $X$ cannot be connected.
	\end{itemize}
\end{ex}
\begin{prop}[Continuous image of a connected set is connected]
	\label{prop:cts_im_connect}
	Let $(X,d_X)$ and $(Y,d_Y)$ be metric spaces and $f:X\to Y$ is a continuous function. If $X$ is connected, then $f(X)$ is connected.
\end{prop}
\begin{proof}
	We show that we can reduce to the case that $f$ is surjective, wlog. Suppose we know~\Cref{prop:cts_im_connect} is true for continuous surjections. In the case when $f$ is not surjective, we can consider $f_1: X \to f(X)$ defined by $f_1(x) = f(x)$ for all $x\in X$. The map $f_1$ is a continuous surjection by construction and so~\Cref{prop:cts_im_connect} for continuous surjections says that $f_1(X)$ is connected. But we also know that $f_1(X) = f(X)$ which shows that $f(X)$ is connected.
	
	Now we assume that $f$ is a surjection so that $Y = f(X)$. Suppose, for a contradiction, that $Y$ is not connected. This means that there exists non-empty and open $U,V\subset Y$ such that
	\[
	U\cup V = Y\quad\text{and}\quad U\cap V = \emptyset.
	\]
	\textcolor{blue}{We define $A:= f^{-1}(U)$ and $B := f^{-1}(V)$.} Since $U,V$ are non-empty and open, the continuity of $f$ (c.f.~\Cref{prop:cts_pre_open}) gives that \textcolor{blue}{$A,B$ are also non-empty and open and $A,B\neq X$.} In particular, we have
	\[
	\textcolor{blue}{A \cup B = }f^{-1}(U)\cup f^{-1}(V) = f^{-1}(U\cup V) = f^{-1}(Y) = \textcolor{blue}{X}
	\]
	and
	\[
	\textcolor{blue}{A\cap B = }f^{-1}(U)\cap f^{-1}(V) = f^{-1}(U\cap V) = f^{-1}(\emptyset) = \textcolor{blue}{\emptyset}.
	\]
	The text in \textcolor{blue}{blue} contradicts the assumption that $X$ is connected.
\end{proof}
In the particular case when $Y = \R$, we can give a proof of the intermediate value theorem for general metric space domains.
\begin{thm}[Intermediate Value Theorem]
	\label{thm:IVT}
	Let $(X,d)$ be a connected metric space and $f:X\to \R$ be a continuous function. Let $y_1,y_2\in f(X)$. Then, for every $y\in \R$ between $y_1,y_2$ (meaning $y_1\le y\le y_2$ if $y_1 \le y_2$ or $y_1 \ge y \ge y_2$ if $y_1 \ge y_2$), there exists $\xi \in X$ such that $f(\xi) = y$.
\end{thm}
\begin{proof}
	According to~\Cref{prop:cts_im_connect}, we know that $f(X)$ is a connected subset of $\R$. Hence, by~\Cref{thm:connected_interval} we see that $f(X)$ must be an interval. Moreover, according to~\Cref{prop:char_int} since $y_1,y_2 \in f(X)$ and $y$ is between $y_1$ and $y_2$, we must also have that $y \in f(X)$. By definition, this implies \textcolor{blue}{there exists some $\xi \in X$ such that $f(\xi) = y$.}
\end{proof}
The next result discusses how adding limit points does not change the connectedness of a set.
\begin{prop}
	Let $(X,d)$ be a metric space and suppose $A\subset X$ is a connected set. Suppose $B\subset X$ satisfies $A\subset B \subset \bar{A}$. Then, $B$ is connected.
\end{prop}
\begin{proof}
	Assume, for a contradiction, that $B$ is not connected. According to~\Cref{prop:connected_two-point}, there exists some continuous surjection $f:B\to \{0,1\}$. Since $A$ is connected, we must have that $f|_A$ is a constant. Suppose, wlog, that
	\[
	f(a) = 0,\quad\forall a\in A,
	\]
	otherwise, exchange the roles of $0$ and $1$. Now, since we assume that $f:B\to \{0,1\}$ is a surjection, \textcolor{blue}{there must exist some $b\in B\setminus A$ such that $f(b) = 1$.} Since $B\subset \overline{A}$, we see that $b$ is a limit point of $A$. According to~\Cref{lem:lim_pts}, there exists a sequence $(a_n)_{n\in\N}\subset A$ such that $a_n \overset{n\to\infty}{\to} b$. By continuity of $f$, we obtain the following contradiction:
	\[
	0 = f(a_n) \overset{n\to\infty}{\to} f(b) = 1.
	\]
\end{proof}
We next investigate a notion which can lead to the beginnings of algebraic topology. For our purposes, it is a specific case of connectedness.
\subsection{Path-Connectedness}
\label{sec:path_connect}
\begin{defn}[Path-connectedness]
	Let $(X,d)$ be a non-empty metric space. We say that $X$ is \textit{path-connected} provided the following holds true: For any $x,y\in X$, there exists a continuous map $f:[0,1]\to X$ such that $f(0) = x$ and $f(1) = y$. We refer to such an $f$ as a \textit{path (joining $x$ to $y$)}.
\end{defn}
\begin{eg}
	$\R^d$ is path-connected for any $d\in\N$. Indeed, given $x,y\in\R^d$, one can define
	\[
	f(t) := (1-t)x + ty.
	\]
\end{eg}
\begin{prop}[Path-connected implies connected]
	\label{prop:path_connect_implies_connect}
	Let $(X,d)$ be a path-connected metric space. Then, $X$ is also connected.
\end{prop}
\begin{proof}
	We will use~\Cref{prop:connected_two-point} for this. Suppose, for a contradiction that there exists a continuous surjection $g:X\to \{0,1\}$. There must exist $x,y\in X$ such that $g(x) = 0$ and $g(y) = 1$. By assumption that $X$ is path-connected, let $f$ be a path joining $x$ to $y$. Owing to~\Cref{prop:comp_cts}, the composition $g\circ f:[0,1]\to \{0,1\}$ must be continuous and surjective. Appealing to~\Cref{thm:connected_interval}, this contradicts the fact that $[0,1]$ is connected.
\end{proof}
The converse to~\Cref{prop:path_connect_implies_connect} is not, however, true in general.
\begin{eg}[Connected does not imply path-connected]
	\label{eg:top_sine}
	Define the following function
	\[
	g(x):=\left\{
	\begin{array}{cc}
		0, 	&\text{if }x=0 	\\
		\sin \frac{1}{x},	&\text{if }x >0
	\end{array}
	\right..
	\]
	This is known as the \textit{topologist's sine curve} and a guided analysis will be carried through on a problem sheet.
\end{eg}
Despite~\Cref{eg:top_sine}, for open subsets of $\R^d$, the result is true.
\begin{thm}
	\label{thm:open_connected_real}
	Let $S\subset \R^d$ be (non-empty) open and connected. Then $S$ is path-connected.
\end{thm}
\begin{proof}
	Fix $x\in S$. We partition $S$ into $A\subset S$ where
	\[
	A:= \left\{
	y \in S \, | \, \exists\text{ continuous path between $x$ and $y$ lying in $S$}
	\right\}
	\]
	and $B:= S\setminus A$. Clearly, we have
	\[
	S = A\cup B\quad\text{and}\quad A\cap B = \emptyset.
	\]
	We will demonstrate that $A = S$ which, since $x\in S$ was arbitrary, will prove that $S$ is path-connected. 
	
	To this aim, we prove that $A,B$ are open. Fix $a\in A$. By definition, there exists a continuous path between $x$ and $a$ in $S$. Since $a\in S$ and $S$ is open, there exists $r>0$ such that $B_r(a) \subset S$. For every $y \in B_r(a)$, one can create a continuous path between $x$ and $y$ by concatenating first the path between $x$ and $a$ and then the straight line between $a$ and $y$ (exercise). Hence, we obtain $y\in A$ and, moreover, $B_r(a) \subset A$.
	
	Fix any $b\in B$. Again, since $S$ is open, there exists $r>0$ such that $B_r(b) \subset S$. If there was $y\in B_r(b)$ such that a continuous path between $x$ and $y$ existed, then one can build from this (in a similar way as in the previous paragraph) a continuous path between $x$ and $b$. Hence, no continuous path between $x$ and any $y\in B_r(b)$ can exist and so $B_r(b)\subset B$.
	
	The previous two paragraphs showed that $A,B$ are open. Recall from construction that that $A\cup B = S$ and $A\cap B = \emptyset$. Since $S$ is connected, this can only happen when $\{A,B\} = \{S,\emptyset\}$. $A$ cannot be empty because $x\in A$ (the trivial path staying at $x$ is continuous), hence $B = \emptyset$. Therefore, we obtain $A = S$.
\end{proof}
	\chapter{Metric space of continuous functions}
	\label{ch:C^zero}
	This chapter focuses on the setting when $(X,d)$ is a metric space \textit{whose elements are (continuous) functions.} This chapter will combine what we have learned previously while also developing insights into function spaces. 
	
	We will discuss extremely useful results about function spaces which appear in modern research.
%	
%	To set the notation for this chapter, given metric spaces $(S,d_S),(Y,d_Y)$, we will denote the set
%	\[ 
%	C(S,Y) := \{\text{continuous functions }f:S\to Y\}.
%	\]
%	Specifically, $f\in C(S,Y)$ means that $f$ is a function from $S$ to $Y$ which is continuous in the sense of~\Cref{defn:cty}. In the event that $Y = \R$, we will abbreviate $C(S) := C(S,\R)$. We will often explicitly identify the metric spaces $S,Y$, but our standing assumptions this chapter will be:
%	\begin{itemize}
%		\item $S$ is a \textbf{compact} metric space and
%		\item $Y=\R$. In particular, $Y$ is \textbf{complete}.
%	\end{itemize}
%	With this in mind, we will focus on the \textit{sup metric} defined by
%	\[
%	d(f,g):= \sup_{x\in S}|f(x) - g(x)|,\quad \forall f,g \in C(S).
%	\]
%	\begin{ex}
%		Show that $d:C(S)\times C(S)\to [0,+\infty)$ satisfies being a metric as in~\Cref{defn:metric}, hence $(C(S),d)$ is a metric space. \textit{Note:} What prevents $d(f,g)$ from attaining $+\infty$?
%	\end{ex}
%	The role of $Y = \R$ of course makes the sup metric well-defined. The completeness property of it will be exploited later on.
%	
%	For notational convenience, we will also use the \textit{sup norm} which is any of
%	\[
%	\|f - g\|,\text{ or }\|f-g\|_\infty,\text{ or }\|f - g\|_{C(S)},
%	\]
%	all of the above referring to $d(f,g)$.
%	
%	Another space which we will not focus on is the following:
%	\[
%	B(S,Y):= \{
%	\text{bounded functions }f:S\to Y
%	\}.
%	\]
%	This set consists of functions $f:S\to Y$ such that the image set $f(S)\subset Y$ is bounded (c.f.~\Cref{defn:bdd}).
%	\begin{ex}
%		Let $(S,d_S)$ be a metric space. Show that $(B(S,\R),d)$ is a metric space. Here, $B(S,\R)$ is also equipped with the sup metric.
%	\end{ex}
%	\begin{prop}[$C(S,Y)$ is complete]
%		\label{prop:C^0_complete}
%		Let $(S,d_S)$ be a compact metric space and $Y = \R^d$ for any $d\ge 1$. Then $C(S,Y)$ equipped with the sup metric is a complete metric space.
%	\end{prop}
%	\begin{proof}
%		\textcolor{blue}{Let $(f_n)_{n\in\N}\subset C(S,Y)$ be a Cauchy sequence.} This means that, for any $\varepsilon>0$, there exists $N\in\N$ such that
%		\[
%		n,m\ge N \implies \sup_{x\in S}|f_n(x) - f_m(x)|<\varepsilon.
%		\]
%		For any fixed $x\in S$, in particular, we see that the sequence $(f_n(x))_{n\in\N}$ of real vectors is a Cauchy sequence. By completeness of $\R^d$, this sequence admits a limit which we shall denote $f(x)$, emphasising the dependence on $x\in S$. Specifically, we can view this as a function \textcolor{blue}{$f:X \to Y$.} We now seek to show that this object $f$ belongs to $C(S,Y)$.
%	\end{proof}
	\section{Notions of convergence}
\label{sec:cvg_cts}
Let $S$ be a set and, for $n\in\N$, let $f_n:S \to \R$ be a function so that we have a sequence of functions $(f_n)_{n\in\N}$. For every $x\in S$, we obtain a sequence of real numbers $(f_n(x))_{n\in\N}$. By looking at this sequence for each $x\in S$, we can immediately lift~\Cref{defn:cvg_real} to our first notion of convergence of functions.
\begin{defn}[Pointwise convergence]
	\label{defn:ptwise}
	Let $S$ be a set. Let $f:S\to \R$ and $(f_n)_{n\in\N}$ be a sequence of functions $f_n:S\to \R$. We say that $f_n$ \textit{converges pointwise to $f$ on $S$} provided the following holds:
	\[
	\forall \varepsilon>0,\,\forall x\in S,\, \exists N\in\N\text{ s.t. }n\ge N \implies |f_n(x) - f(x)|<\varepsilon.
	\]
	We also write
	\[
	f(x) = \lim_{n\to \infty} f_n(x).
	\]
\end{defn}
The notion of pointwise convergence may seem very natural. In practice, it is extremely poor for our purposes. This is somewhat remedied when one considers the context of measure theory, but this is beyond the scope of these notes. The following examples illustrate the general flaw of~\Cref{defn:ptwise}: The sequence $(f_n)_{n\in\N}$ might satisfy a certain property, but the pointwise limit $f$ can fail to satisfy that property.
\begin{ex}[Pointwise limit of bounded functions]
	Let $S = [0,1]$ and consider the sequence of functions
	\[
	f_n(x):=\left\{
	\begin{array}{cl}
		\min \left(n,\frac{1}{x}\right), 	&\text{if }0<x\le 1 	\\
		0, 	&\text{if }x=0
	\end{array}
	\right.,\quad n\in\N.
	\]
	Denote
	\[
	f(x) := \left\{
	\begin{array}{cl}
		\frac{1}{x},	&\text{if }0<x\le 1 	\\
		0, 	&\text{if }x=0
	\end{array}
	\right..
	\]
	Prove the following:
	\begin{itemize}
		\item For every $n\in\N$, the function $f_n$ is bounded. i.e. show that $f_n([0,1])$ is a bounded subset of $\R$.
		\item $f_n$ converges pointwise to $f$ on $S$.
		\item $f$ is an unbounded function.
	\end{itemize}
\end{ex}
\begin{ex}[Pointwise limit of continuous functions]
	Let $S = [0,1]$ and consider the sequence of functions
	\[
	f_n(x) := x^n,\quad \forall x\in [0,1],\,\forall n\in\N,
	\]
	together with the function
	\[
	f(x) :=\left\{
	\begin{array}{cl}
		0, 	&\text{if }x \neq 1 	\\
		1, 	&\text{if }x = 1
	\end{array}
	\right..
	\]
	Notice that, for every $n\in\N$, $f_n$ is continuous on $[0,1]$, but $f$ is not. Prove that $f_n$ converges pointwise to $f$ on $[0,1]$.
\end{ex}
\begin{ex}[Pointwise limit of integrable functions]
	\label{ex:pt_int}
	Let $S = [0,1]$ and consider the sequence of functions
	\[
	f_n(x):= \left\{
	\begin{array}{cl}
		n, 	&\text{if }0 < x < \frac{1}{n} 	\\
		0, 	&\text{otherwise}
	\end{array}
	\right.,\quad x\in [0,1],\,n\in\N.
	\]
	Prove the following:
	\begin{itemize}
		\item Show that $f_n$ converges pointwise to $0$ in $[0,1]$. Here, the limit is the function which is identically zero on $[0,1]$.
		\item Calculate and compare
		\[
		\lim_{n\to \infty}\left(\int_0^1f_n(x)\,\mathrm{d}x\right)\quad\text{versus}\quad \int_0^1\left(\lim_{n\to\infty}f_n(x)\right)\,\mathrm{d}x.
		\]
		Are their values the same or different?
	\end{itemize}
\end{ex}
All the previous examples are meant to illustrate the deficiency of pointwise convergence in a variety of situations. The stronger and more useful notion is the following.
\begin{defn}[Uniform convergence]
	\label{defn:unif_cvg}
	Let $S$ be a set. Let $f:S\to \R$ and $(f_n)_{n\in\N}$ be a sequence of functions $f_n:S\to \R$. We say that $f_n$ \textit{converges uniformly to $f$ on $S$} provided the following holds:
	\[
	\forall \varepsilon>0,\,\exists N\in\N\text{ s.t. }\forall x \in S, n\ge N \implies |f_n(x) - f(x)|<\varepsilon.
	\]
	In this case, we also write `$f_n\to f$ uniformly on $S$ (as $n\to \infty$).'
\end{defn}
Notice the placement of `$\forall x \in S$' in~\Cref{defn:ptwise} and~\Cref{defn:unif_cvg}! In~\Cref{defn:ptwise}, the chosen $N\in \N$ may change depending on $x\in S$. By contrast in~\Cref{defn:unif_cvg}, the chosen $N\in\N$ does not depend on $x\in S$. This is the main crucial difference and it is what makes uniform convergence more advantageous than pointwise convergence.
\begin{ex}
	Revisit the previous examples demonstrating the failure of pointwise convergence. Show, in each of those examples, that the sequence of functions does not converge uniformly to the stated limit functions.
\end{ex}
\begin{rem}[Equivalence of~\Cref{defn:unif_cvg} with sup]
	\label{rem:sup_unif_cvg}
	Notice (exercise) the following equivalence
	\[
	|f_n(x) - f(x)|<\varepsilon,\,\forall x \in S \iff \sup_{x\in S}|f_n(x) - f(x)|.
	\]
	Hence, we could have formulated~\Cref{defn:unif_cvg} to say
	\[
	\forall \varepsilon>0,\exists N\in \N\text{ s.t. }n\ge N\implies \sup_{x\in S}|f_n(x) - f(x)|<\varepsilon.
	\]
	This is a purely cosmetic distinction which does not change the mathematics whatsoever. We will use the supremum convention almost exclusively much later on.
\end{rem}
Let us clarify that uniform convergence is actually stronger than pointwise convergence.
\begin{lemma}[Uniform implies pointwise]
	Let $S$ be a set. Let $f:S\to \R$ and $(f_n)_{n\in\N}$ be a sequence of functions $f_n:S\to \R$. Suppose that $f_n$ converges uniformly to $f$ on $S$. Then, $f_n$ converges pointwise to $f$ on $S$ as well.
\end{lemma}
\begin{proof}
	Exercise.
\end{proof}
We will begin our specialisation to continuous functions.
\begin{prop}[Uniform limit of continuous functions]
	\label{prop:unif_cvg_cts}
	Let $(S,d)$ be a metric space and $f:S\to \R$. Let $(f_n)_{n\in\N}$ be a sequence of continuous functions $f_n:S\to \R$ for each $n\in\N$. Suppose that $f_n$ converges uniformly to $f$ on $S$. Then, $f$ is continuous.
\end{prop}
\begin{proof}
	\textcolor{blue}{Fix any $c \in S$ and $\varepsilon>0$.} Owing to the uniform convergence, we know that there exists $N\in \N$ such that
	\[
	n\ge N\implies |f_n(x) - f(x)|<\frac{\varepsilon}{3},\quad \forall x \in S.
	\]
	In particular, since $f_N$ is continuous, we know that \textcolor{blue}{there exists $\delta>0$} such that, for any $x \in S$, we have
	\[
	d(x,c)<\delta\implies |f_N(x) - f_N(c)|<\frac{\varepsilon}{3}.
	\]
	We combine the previous statements. \textcolor{blue}{Fix any $x\in S$ such that $d(x,c)<\delta$.} The triangle inequality yields
	\begin{align*}
		\textcolor{blue}{|f(x)-f(c)|}&\le |f(x) - f_N(x)| + |f_N(x) - f_N(c)| + |f_N(c) - f(c)| 	\\
		&\textcolor{blue}{<} \frac{\varepsilon}{3}+\frac{\varepsilon}{3}+\frac{\varepsilon}{3} = \textcolor{blue}{\varepsilon.}
	\end{align*}
\end{proof}
We end this section with a useful result.
\begin{thm}[Cauchy criterion for uniform convergence]
	\label{thm:unif_cauchy}
	Let $S$ be a set and let $(f_n)_{n\in\N}$ be a sequence of functions $f_n:S\to \R$. TFAE:
	\begin{itemize}
		\item There exists a function $f:S\to \R$ such that $f_n$ converges to $f$ uniformly on $S$.
		\item $\forall \varepsilon>0, \exists N\in\N$ such that
		\[
		m,n\ge N\implies |f_n(x) - f_m(x)|<\varepsilon,\quad \forall x\in S.
		\]
	\end{itemize}
\end{thm}
\begin{proof}
	\ul{$\implies$:} \textcolor{blue}{Fix any $\varepsilon>0$.} By the assumption of uniform convergence, \textcolor{blue}{there exists $N\in \N$} such that
	\[
	n\ge N \implies |f_n(x) - f(x)|<\frac{\varepsilon}{2},\quad \forall x \in S.
	\]
	In particular, the triangle inequality gives the following:
	\begin{align*}
		\textcolor{blue}{n,m\ge N\implies|f_n(x) - f_m(x)|} &\le |f_n(x) - f(x)| + |f_m(x) - f(x)| \\
		&\textcolor{blue}{<}\frac{\varepsilon}{2} + \frac{\varepsilon}{2} = \textcolor{blue}{\varepsilon}.
	\end{align*}
	\ul{$\impliedby$:} For any $x\in S$, consider the sequence of real numbers $(f_n(x))_{n\in\N}$. By assumption, we see that this is a Cauchy sequence. Hence, we can define a function $f:S\to \R$ by
	\[
	f(x) = \lim_{n\to \infty}f_n(x),\quad \forall x\in S.
	\]
	We aim now to prove that $f_n$ converges uniformly to $f$. \textcolor{blue}{Fix any $\varepsilon>0$.} By assumption, \textcolor{blue}{there exists $N\in \N$} such that
	\[
	m,n\ge N\implies |f_n(x) - f_m(x)|< \frac{\varepsilon}{2},\quad \forall x \in S.
	\]
	In particular, \textcolor{blue}{fix any $x\in S$} and let $M\ge N$ be such that
	\[
	m \ge M \implies |f_m(x) - f(x)| < \frac{\varepsilon}{2}.
	\]
	By the triangle inequality, \textcolor{blue}{for any $n\ge N$, we have}
	\[
	\textcolor{blue}{|f_n(x) - f(x)|} \le |f_n(x) - f_M(x)| + |f_M(x) - f(x)| \textcolor{blue}{<}\frac{\varepsilon}{2} + \frac{\varepsilon}{2} = \textcolor{blue}{\varepsilon}.
	\]
\end{proof}
The usefulness of~\Cref{thm:unif_cauchy} comes from the fact that that the second point makes no explicit mention of a limit function $f$. Hence, the Cauchy criterion can be used to prove uniform convergence of a sequence of functions without relying on knowing the limit.
	\section{Uniform convergence and integration}
\label{sec:unif_int}
We have seen in~\Cref{ex:pt_int} how interchanging limits with integration can be false for pointwise convergent sequences of functions. In this short section, we prove a result which guarantees when this is valid.

We first recall some concepts about integration and integrability.
\begin{defn}[Riemann sums, integrability]
	\label{defn:riemann}
	For a closed and bounded interval $[a,b]\subset\R$, a \textit{partition} of $[a,b]$ is any finite set of the form
	\[
	P = \{x_0,x_1,\dots, x_n\},\quad a = x_0 < x_1 < \dots < x_n = b.
	\]
	Given a bounded function $f:[a,b]\to \R$ and a partition $P = \{x_i\}_{i=0}^n$ of $[a,b]$, its \textit{upper Riemann sum} is
	\[
	U(f,P) = \sum_{i=1}^n (x_i-x_{i-1})\left(\sup_{x_{i-1}\le x \le x_i}f(x)\right).
	\]
	The corresponding \textit{lower Riemann sum} is
	\[
	L(f,P) = \sum_{i=1}^n (x_i-x_{i-1})\left(\inf_{x_{i-1}\le x \le x_i}f(x)\right).
	\]
	We say that a (bounded) function $f:[a,b]\to \R$ is \textit{(Riemann) integrable on $[a,b]$} provided the following is true
	\[
	\sup_{P\subset [a,b]}L(f,P) = \inf_{P\subset [a,b]}U(f,P).
	\]
	Above, the supremum / infimum runs over all partitions $P$ of $[a,b]$.
\end{defn}
An important result to determine Riemann integrability is the following. \textcolor{blue}{A proof of this can be found in an earlier course.}
\begin{thm}[Cauchy criterion for Riemann integrability]
	\label{thm:cauchy_int}
	Let $f:[a,b]\to\R$ be bounded. TFAE
	\begin{itemize}
		\item $f$ is integrable on $[a,b]$.
		\item For every $\varepsilon>0$, there exists a partition $P$ of $[a,b]$ such that
		\[
		U(f,P) - L(f,P)<\varepsilon.
		\]
	\end{itemize}
\end{thm}
Before we proceed to the main theorem of this section, we present a preliminary result in the same theme as~\Cref{prop:unif_cvg_cts} -- uniform convergence preserves properties of the function sequence.
\begin{prop}
	Let $-\infty < a < b < +\infty$. For every $n\in\N$, let $f_n:[a,b]\to\R$ be a bounded function. Moreover, assume that there exists $f:[a,b]\to\R$ such that $f_n$ converges uniformly to $f$ on $[a,b]$. Then, the function $f$ is also bounded.
\end{prop}
\begin{proof}
	With $\varepsilon=1$, the uniform convergence assumption guarantees the existence of $N\in\N$ such that
	\[
	n\ge N \implies \sup_{x\in[a,b]}|f_n(x) - f(x)|<1.
	\]
	The boundedness assumption applied particularly to $f_N$ guarantees the existence of $\tilde{M}>0$ such that
	\[
	\sup_{x\in[a,b]}|f_N(x)|\le \tilde{M}.
	\]
	\textcolor{blue}{We define $M:= \tilde{M}+1$. For any $x\in[a,b]$,} we deduce from the triangle inequality and the previous deductions that
	\[
	\textcolor{blue}{|f(x)| \le} |f(x) - f_N(x)| + |f_N(x)| \le 1 + \tilde{M} = \textcolor{blue}{M}.
	\]
\end{proof}
This result says that we can calculate the Riemann sums of a uniform limit of bounded functions. We want to say more with respect to integration.
\begin{thm}
	\label{thm:unif_int}
	Let $-\infty<a<b<+\infty$ and $f_n:[a,b]\to\R$ be a sequence of bounded functions for $n\in\N$. Assume, moreover that $f_n$ is integrable on $[a,b]$ for every $n\in\N$ and that there exists $f:[a,b]\to\R$ such that $f_n$ converges uniformly to $f$ on $[a,b]$. Then, we have the following:
	\begin{enumerate}
		\item $f$ is integrable on $[a,b]$.
		\item The integral is given by
		\[
		\int_a^bf(x)\,\mathrm{d}x = \lim_{n\to\infty}\int_a^b f_n(x)\,\mathrm{d}x.
		\]
	\end{enumerate}
\end{thm}
\begin{proof}
	We will use~\Cref{thm:cauchy_int}.
	\begin{enumerate}
		\item \textcolor{blue}{Fix any $\varepsilon>0$.} Since $f_n$ converges uniformly to $f$ on $[a,b]$, we know that there exists some $N\in\N$ such that
		\[
		n\ge N\implies |f_n(x) - f(x)|<\frac{\varepsilon}{3(b-a)},\quad \forall x\in[a,b].
		\]
		In particular, for any partition $Q$ of $[a,b]$, we can deduce from this statement that (exercise)
		\[
		\left|U(f-f_N,Q)\right| < \frac{\varepsilon}{3}\text{ and }\left|L(f-f_N,Q)\right| < \frac{\varepsilon}{3}.
		\]
		Since $f_N$ is integrable, we know from~\Cref{thm:cauchy_int} that \textcolor{blue}{there exists a partition $P$ of $[a,b]$} such that
		\[
		U(f_N,P) - L(f_N,P)< \frac{\varepsilon}{3}.
		\]
		The previous deductions give
		\begin{align*}
			&\quad \textcolor{blue}{U(f,P)-L(f,P)} \le U(f-f_N,P) + U(f_N,P) - L(f-f_N,P) - L(f_N,P) \\
			&\le \left|U(f-f_N,Q)\right| +\left|L(f-f_N,Q)\right| + U(f_N,P) - L(f_N,P) \textcolor{blue}{<} \frac{\varepsilon}{3} + \frac{\varepsilon}{3} + \frac{\varepsilon}{3} = \textcolor{blue}{\varepsilon.} 
		\end{align*}
		Hence, $f$ satisfies the Cauchy criterion for integrability and so, by~\Cref{thm:cauchy_int}, we have established that $f$ is integrable on $[a,b]$.
		\item We will show that the sequence of numbers $\left(\int_a^bf_n(x)\,\mathrm{d}x\right)_{n\in\N}$ converges to $\int_a^b f(x)\,\mathrm{d}x$. \textcolor{blue}{Fix any $\varepsilon>0$.} Since $f_n$ converges uniformly to $f$ on $[a,b]$, we know that \textcolor{blue}{there exists some $N\in\N$} such that
		\[
		n\ge N\implies |f_n(x) - f(x)|<\frac{\varepsilon}{b-a},\quad\forall x\in[a,b].
		\]
		\textcolor{blue}{For any $n\ge N$, }we apply the triangle inequality for integrals to see that
		\begin{align*}
			&\quad \textcolor{blue}{\left|\int_a^b f_n(x)\,\mathrm{d}x - \int_a^b f(x)\,\mathrm{d}x\right|} = \left|\int_a^bf_n(x) - f(x)\,\mathrm{d}x\right| \le \int_a^b\left|f_n(x) - f(x)\right| \,\mathrm{d}x \\
			&\textcolor{blue}{<}\int_a^b\frac{\varepsilon}{b-a}\,\mathrm{d}x = \frac{\varepsilon}{b-a}\int_a^b\,\mathrm{d}x = \textcolor{blue}{\varepsilon}.
		\end{align*}
	\end{enumerate}
\end{proof}
	\section{Uniform convergence with derivatives}
\label{sec:unif_diff}
\Cref{prop:unif_cvg_cts} illustrated how uniform convergence preserves continuity. In this section, we will explore uniform convergence with derivatives of functions, function series, and their combination.
\begin{thm}[Uniform convergence of differentiable functions]
	\label{thm:unif_cvg_diff}
	Let $(a,b)\subset \R$ be an open and bounded interval. Let $(f_n)_{n\in\N}$ be a sequence of functions $f_n:(a,b)\to \R$ satisfying the following assumptions:
	\begin{enumerate}[label = \textbf{A\roman*)}]
		\item \label{unif_cvg_diff_1} For every $n\in\N$, the function $f_n$ is differentiable on $(a,b)$.
		\item \label{unif_cvg_diff_2} There exists a point $x_0 \in (a,b)$ such that the sequence $(f_n(x_0))_{n\in\N}$ is convergent.
		\item \label{unif_cvg_diff_3} There exists a function $g:(a,b)\to \R$ such that $f_n'$ converges uniformly to $g$ on $(a,b)$.
	\end{enumerate}
	Then, we obtain the following:
	\begin{enumerate}[label = \textbf{C\roman*)}]
		\item \label{unif_cvg_diff_C1} There exists a function $f:(a,b)\to \R$ such that $f_n$ converges uniformly to $f$ on $(a,b)$.
		\item \label{unif_cvg_diff_C2} The function $f$ is differentiable on $(a,b)$ and $f' = g$.
	\end{enumerate}
\end{thm}
\begin{proof}
	For every $c\in (a,b)$ and $n\in N$, we define
	\[
	g_n^c(x):= \left\{
	\begin{array}{cl}
		\frac{f_n(x)-f_n(c)}{x-c}, 	&\text{if }x\neq c 	\\
		f_n'(c), 	&\text{if }x=c
	\end{array}
	\right.,\quad x\in (a,b).
	\]
	Notice that $g_n^c(x)$ is continuous at $x=c$ and $\lim_{n\to\infty} g_n^c(c) = \lim_{n\to\infty}f_n'(c) = g(c)$ by~\ref{unif_cvg_diff_3}. 
	
	\ul{$(g_n^c)_{n\in\N}$ is uniformly convergent:} \textcolor{blue}{Fix any $\varepsilon>0$.} Owing to~\ref{unif_cvg_diff_3}, \textcolor{blue}{there exists $N\in \N$} such that
	\[
	m,n\ge N\implies |f_n'(x) - f_m'(x)|<\varepsilon,\quad \forall x\in (a,b).
	\]
	Now, for any $m,n\in\N$ and $x\in (a,b)$ with $x\neq c$, the Mean-Value Theorem (applied to the function $f_n(x) - f_m(x)$) gives the existence of some $\xi$ between $x$ and $c$ such that
	\begin{align*}
	g_n^c(x) - g_m^c(x) &= \frac{f_n(x)-f_n(c)}{x-c} - \frac{f_m(x)-f_m(c)}{x-c} = \frac{[f_n(x) - f_m(x)] - [f_n(c) - f_m(c)]}{x-c} 	\\
	&= f_n'(\xi) - f_m'(\xi).
	\end{align*}
	In particular, \textcolor{blue}{for any $m,n\ge N$,} the previous points yield
	\[
	\textcolor{blue}{|g_n^c(x) - g_m^c(x)|} = |f_n'(\xi) - f_m'(\xi)| \textcolor{blue}{< \varepsilon.\quad \forall x\in (a,b)\text{ s.t. }x\neq c.}
	\]
	Recall that $g_n^c(c) \overset{n\to \infty}{\to} g(c)$. The text in \textcolor{blue}{blue} allows us to apply the Cauchy criterion from~\Cref{thm:unif_cauchy}. Specifically, for every $c\in (a,b)$, we know that there exists some function $G^c:(a,b)\to \R$ such that $g_n^c$ converges uniformly to $G^c$ on $(a,b)$.
	
	\ul{Construction of $f$:} From above, we set $c = x_0$ and we write
	\[
	f_n(x) = f_n(x_0) + g_n^{x_0}(x)(x-x_0),\quad \forall n\in\N,\,\forall x\in(a,b).
	\]
	\textcolor{blue}{Fix any $\varepsilon>0$.} According to~\ref{unif_cvg_diff_2}, there exists $N_1\in\N$ such that
	\[
	m,n\ge N_1\implies |f_n(x_0) - f_m(x_0)|<\frac{\varepsilon}{2}.
	\]
	By what we have just done with $g_n$, we also know that there exists some $N_2\in\N$ such that
	\[
	m,n\ge N_2\implies |g_n^{x_0}(x) - g_m^{x_0}(x)|<\frac{\varepsilon}{2(1 + b-a)},\quad \forall x \in (a,b).
	\]
	\textcolor{blue}{Define $N:= \max(N_1,N_2)$} We combine the previous deductions to see that, \textcolor{blue}{for any $m,n\ge N$ and $x\in (a,b)$, we have}
	\begin{align*}
		\textcolor{blue}{|f_n(x) - f_m(x)|} &\le |f_n(x_0) - f_m(x_0)| + |x-x_0||g_n^{x_0}(x) - g_m^{x_0}(x)| 	\\
		&\le |f_n(x_0) - f_m(x_0)| + (b-a)|g_n^{x_0}(x) - g_m^{x_0}(x)| 	\\
		&< \frac{\varepsilon}{2} + \left(\frac{b-a}{1 + b-a}\right)\frac{\varepsilon}{2} \textcolor{blue}{< \varepsilon.}
	\end{align*}
	Again, the text in \textcolor{blue}{blue} allows us to apply the Cauchy criterion from~\Cref{thm:unif_cauchy} to deduce the existence of a function $f:(a,b)\to \R$ such that $f_n$ converges uniformly to $f$ on $(a,b)$. This achieves~\ref{unif_cvg_diff_C1}.
	
	\ul{Derivative of $f$:} Recall that~\ref{unif_cvg_diff_1} yields continuity of $g_n^c$, particularly at the point $x=c$, so that
	\[
	\lim_{x\to c}g_n^c(x) = g_n^c(c) = f_n'(c).
	\]
	Applying~\Cref{prop:unif_cvg_cts} gives that $G^c$ is continuous, also particularly at the point $x=c$. Together with uniform convergence, we get
	\[
	\lim_{x\to c}G^c(x) = G^c(c) = \lim_{n\to \infty}g_n^c(c) = \lim_{n\to \infty}f_n'(c) = g(c).
	\]
	On the other hand, for every $x\neq c$, we have
	\[
	G^c(x) = \lim_{n\to \infty}g_n^c(x) = \lim_{n\to \infty}\frac{f_n(x)-f_n(c)}{x-c} = \frac{f(x) - f(c)}{x-c}.
	\]
	Putting the previous two points together implies that
	\[
	g(c) = \lim_{x\to c}\frac{f(x)-f(c)}{x-c},
	\]
	which establishes~\ref{unif_cvg_diff_C2}.
\end{proof}
We mentioned in~\Cref{sec:intro} that, for a series of functions given by
\[
f(x) = \sum_{n=1}^\infty f_n(x),
\]
an attractive identity would be
\[
f'(x) \overset{!}{=} \sum_{n=1}^\infty f_n'(x).
\]
Again, the exclamation point is meant to emphasise that this is not true in general! The derivative of an infinite series is not necessarily the same as the infinite series of derivatives. Nevertheless, \Cref{thm:unif_cvg_diff} gives a crucial ingredient into proving the case where this is true. To get there, we will first need to develop some preliminaries for series of functions
	\section{Uniform convergence with series of functions}
\label{sec:unif_fun_series}
We turn to objects of the form
\[
\sum_{n=1}^\infty f_n(x)
\]
and we will make sense of this analogously to~\Cref{defn:ptwise} and~\Cref{defn:unif_cvg}.
\begin{defn}[Pointwise and uniform convergence of series of functions]
	\label{defn:series_cvg}
	Let $S$ be a set and $(f_n)_{n\in\N}$ be a sequence of functions $f_n:S\to \R$ for each $n\in\N$. Let $f:S\to \R$ be a function. We define the \textit{partial sum of the series $\sum f_n$} to be
	\[
	s_n(x):= \sum_{j=1}^nf_j(x),\quad x \in S.
	\]
	We say that:
	\begin{itemize}
		\item \textit{$\sum f_n$ converges pointwise to $f$ on $S$} provided $s_n$ converges pointwise to $f$ on $S$. In this case, we write
		\[
		f(x) = \sum_{n=1}^\infty f_n(x),\quad\text{pointwise on }S.
		\]
		\item \textit{$\sum f_n$ converges uniformly to $f$ on $S$} provided $s_n$ converges uniformly to $f$ on $S$. In this case, we write
		\[
		f(x) = \sum_{n=1}^\infty f_n(x),\quad \text{uniformly on }S.
		\]
	\end{itemize}
\end{defn}
\begin{thm}[Cauchy criterion for uniform convergence of function series]
	\label{thm:unif_cauchy_series}
	Let $S$ be a set and $(f_n)_{n\in\N}$ be a sequence of functions $f_n:S\to \R$ for each $n\in \N$. TFAE:
	\begin{itemize}
		\item There exists $f:S\to \R$ such that $\sum f_n$ converges uniformly to $f$ on $S$.
		\item $\forall \varepsilon>0,\exists N\in\N$ such that
		\[
		m\ge n\ge N \implies \left|
		\sum_{j=n}^mf_j(x)
		\right|<\varepsilon,\quad \forall x \in S.
		\]
	\end{itemize}
\end{thm}
\begin{proof}
	$\implies$: Assuming that $\sum f_n$ converges uniformly to $f$ on $S$, this means that the partial sums $s_n(x) = \sum_{j=1}^nf_j(x)$ converges uniformly to $f$ on $S$. \textcolor{blue}{Fix any $\varepsilon>0$.} Since $s_n\to f$ uniformly, we know that there exists $N_*\in\N$ such that
	\[
	n \ge N_* \implies |s_n(x) - f(x)|<\frac{\varepsilon}{2},\quad \forall x \in S.
	\]
	\textcolor{blue}{Define $N = N_*+1$ and fix $m\ge n\ge N$.} We certainly have that $m,n-1 \ge N_*$ and so the previous part together with the triangle inequality gives
	\[
	\textcolor{blue}{\left|\sum_{j=n}^mf_j(x)\right|} = \left|\sum_{j=1}^mf_j(x) - \sum_{j=1}^{n-1}f_j(x)\right| = |s_m(x) - s_{n-1}(x)| \le |s_m(x) - f(x)| + |s_{n-1}(x) - f(x)| \textcolor{blue}{<}\frac{\varepsilon}{2} + \frac{\varepsilon}{2} = \textcolor{blue}{\varepsilon.}
	\]
	We used the triangle inequality above.
	
	$\impliedby$: We will show that $(s_n)_{n\in\N}$ satisfies the Cauchy criterion for uniform convergence of sequences of functions. \textcolor{blue}{Fix any $\varepsilon>0$}, by assumption, we know that \textcolor{blue}{there exists $N\in \N$} such that
	\[
	m\ge n\ge N\implies |s_m(x) - s_{n-1}(x)|<\varepsilon,\quad \forall x \in S.
	\]
	In particular, \textcolor{blue}{for any $m,n\ge N$}, let us assume without loss of generality that $m\ge n$ (otherwise swap their labels) and we certainly have $n+1 \ge n \ge N$. Hence, we see that
	\[
	\textcolor{blue}{|s_m(x)-s_n(x)|<\varepsilon,\quad \forall x \in S.}
	\]
	The text in \textcolor{blue}{blue} shows that $(s_n)_{n\in\N}$ satisfies the Cauchy criterion in~\Cref{thm:unif_cauchy}. Hence, there exists a function $f:S\to \R$ such that $s_n\to f$ uniformly on $S$. According to~\Cref{defn:series_cvg}, this is exactly what it means for $\sum f_n$ to converge to $f$ on $S$.
\end{proof}
We will use the following variation of~\Cref{thm:unif_cauchy_series}.
\begin{ex}
	\label{ex:unif_abs_series}
	In the context of~\Cref{thm:unif_cauchy_series}, TFAE:
	\begin{itemize}
		\item There exists a function $f:S\to \R$ such that $\sum f_n$ is absolutely convergent and $\sum f_n$ converges uniformly to $f$ on $S$.
		\item For every $\varepsilon>0$, there exists $N\in\N$ such that
		\[
		m\ge n \ge N\implies \sum_{j=n}^m |f_j(x)| < \varepsilon,\quad \forall x\in S.
		\]
	\end{itemize}
\end{ex}
\begin{thm}[Weierstrass $M$-test]
	\label{thm:m_test}
	Let $S$ be a set and $(f_n)_{n\in\N}$ a sequence of functions $f_n:S\to \R$ for each $n\in\N$. Assume that $(M_n)_{n\in\N}\subset [0,+\infty)$ is a sequence satisfying
	\[
	\sup_{x\in S}|f_n(x)| \le M_n,\quad \forall n\in\N.
	\]
	Moreover, assume that $\sum M_n$ is a convergent series. Then, $\sum f_n$ converges absolutely and uniformly to some function $f:S\to \R$.
\end{thm}
\begin{proof}
	We will use~\Cref{ex:unif_abs_series}. \textcolor{blue}{Fix any $\varepsilon>0$.} Since $\sum M_n$ is a convergent series, the Cauchy criterion for it gives \textcolor{blue}{the existence of some $N\in\N$} such that
	\[
	m\ge n \ge N \implies \left|
	\sum_{j=n}^m M_j
	\right| = \sum_{j=n}^m M_j <\varepsilon.
	\]
	In particular, \textcolor{blue}{for any $m\ge n\ge N$, we have}
	\[
	\textcolor{blue}{ \sum_{j=n}^m |f_j(x)|} \le \sum_{j=n}^m M_j \textcolor{blue}{< \varepsilon, \text{ for every }x\in S.}
	\]
	The text in \textcolor{blue}{blue} allows us to apply~\Cref{ex:unif_abs_series} and conclude the result.
\end{proof}
We return to one of the motivations which was the interchange of derivatives with series. 
\begin{thm}[Uniformly convergent series of differentiable functions]
	\label{thm:unif_cvg_series_diff}
	Let $(a,b)\subset \R$ be an open and bounded interval. Let $(f_n)_{n\in\N}$ be a sequence of functions $f_n:(a,b)\to \R$ satisfying the following assumptions:
	\begin{enumerate}
		\item For every $n\in\N$, the function $f_n$ is differentiable on $(a,b)$.
		\item There exists a point $x_0\in(a,b)$ such that the series $\sum f_n(x_0)$ converges.
		\item There exists a function $g:(a,b)\to \R$ such that $\sum f_n'$ converges uniformly to $g$ on $(a,b)$.
	\end{enumerate}
	Then, we obtain the following:
	\begin{enumerate}
		\item There exists a function $f:(a,b)\to \R$ such that $\sum f_n$ converges uniformly to $f$ on $(a,b)$.
		\item The function $f$ is differentiable on $(a,b)$ and $f' = g$.
	\end{enumerate}
\end{thm}
\begin{proof}
	This is a simple application of~\Cref{thm:unif_cvg_diff} after unpacking the definitions with the partial sums.
\end{proof}
	\section{Power series}
\label{sec:power_series}
\begin{defn}[Power series]
	\label{defn:power_series}
	An infinite series of the form
	\[
	\sum_{n=0}^\infty a_n (x-x_0)^n
	\]
	is called a \textit{power series} centred at $x_0$.
\end{defn}
For a power series, we typically think of $(a_n)_{n\in\N}$ and $x_0$ as given numbers i.e. constants. Conversely, we treat $x$ as an independent variable so that the power series is actually a function of $x$. It is clear that a power series converges at least at $x=x_0$ (not that interesting). What is more interesting is whether there are other points $x\neq x_0$ for which it converges.
\begin{rem}
	Suppose that a power series converges absolutely at some point $x=x_1$. Then, it must also converge absolutely at all other points $x$ satisfying $|x-x_0|\le |x_1-x_0|$ i.e. all points which are closer to $x_0$.
	
	Indeed, for any fixed $n\in \N\cup\{0\}$, we have
	\[
	|a_n(x-x_0)^n| \le |a_n||x_1-x_0|^n.
	\]
	By the comparison test, we obtain convergence.
\end{rem}
\begin{thm}[Convergence of a power series]
	\label{thm:cvg_power}
	Consider a power series $\sum_{n=0}^\infty a_n(x-x_0)^n$ and set
	\[
	\lambda:= \limsup_{n\to \infty}|a_n|^\frac{1}{n}.
	\]
	The following statements holds.
	\begin{enumerate}
		\item If $\lambda = +\infty$, then the series converges only for $x=x_0$.
		\item If $0<\lambda < +\infty$, then:
		\begin{enumerate}
			\item the power series converges absolutely for any $x$ satisfying $|x-x_0|< \frac{1}{\lambda}$.
			\item the power series diverges for any $x$ such that $|x-x_0|>\frac{1}{\lambda}$.
		\end{enumerate}
		Moreover, the convergence is uniform on any closed interval contained in $\left(x_0 - \frac{1}{\lambda},x_0 + \frac{1}{\lambda}\right)$.
		\item If $\lambda = 0$, then the series converges absolutely for all $x\in\R$. Moreover, the convergence is uniform on any closed bounded interval.
	\end{enumerate}
\end{thm}
The quantity $r = \frac{1}{\lambda}$ (with $r = 0$ when $\lambda = +\infty$ and $r=+\infty$ when $\lambda=0$) is called the \textit{radius of convergence}. The interval $(x_0-r,x_0+r)$ is called the \textit{disk / interval of convergence}.
\begin{proof}
	We will apply the root test to the power series. Recall that, for a series $\sum b_n$ with $\rho = \limsup_{n\to\infty}|b_n|^\frac{1}{n}$, the root test says
	\begin{enumerate}
		\item If $\rho <1$, then $\sum b_n$ converges absolutely.
		\item If $\rho>1$, then $\sum b_n$ diverges.
	\end{enumerate}
	We apply the root test with $b_n = a_n(x-x_0)^n$ so that
	\[
	\limsup_{n\to\infty}|a_n(x-x_0)^n|^\frac{1}{n} = |x-x_0|\limsup_{n\to\infty}|a_n|^\frac{1}{n} = \lambda|x-x_0|.
	\]
	Therefore, we see that the power series converges absolutely if $|x-x_0|<\frac{1}{\lambda}$ and diverges if $|x-x_0|> \frac{1}{\lambda}$. We now turn to the uniform convergence statements.
	
	Fix any closed and bounded interval $[a,b]$ within the disk of convergence i.e.
	\[
	[a,b] \subset (x_0 - r,x_0+r).
	\]
	There exists (why?) $0<\delta < r$ such that
	\[
	[a,b] \subset [x_0 - r + \delta, x_0 + r - \delta].
	\]
	Define now $y:= x_0 - r + \delta$. For every $x\in[a,b]$, we have
	\[
	|a_n(x-x_0)^n| \le |a_n| |y - x_0|^n =:M_n.
	\]
	Observe that $|y - x_0| = r - \delta < r = \frac{1}{\lambda}$. Applying the root test to $\sum M_n$ gives
	\[
	\limsup_{n\to\infty}|M_n|^\frac{1}{n} = \limsup_{n\to\infty}|a_n|^\frac{1}{n}|y-x_0| = |y-x_0|\limsup_{n\to\infty}|a_n|^\frac{1}{n} < 1.
	\]
	We see that $\sum M_n$ converges (absolutely) and, owing to~\Cref{thm:m_test}, we conclude uniform convergence on $[a,b]$.
\end{proof}
Anything can happen at the boundary of the disk of convergence.
\begin{ex}
	Consider the series (with centres $x_0 = 0$)
	\[
	\sum_{n=0}^\infty x^n,\, \sum_{n=0}^\infty \frac{x^n}{n},\text{ and }\sum_{n=0}^\infty \frac{x^n}{n^2}.
	\]
	Verify that the radius of convergence of all of these series is 1. What happens at $x=\pm 1$ for these series?
\end{ex}
\begin{thm}
	\label{thm:cty_power}
	Suppose that a power series $\sum_{n=0}^\infty a_n(x-x_0)^n$ has a strictly positive radius of convergence i.e.
	\[
	r = \frac{1}{\lambda}>0\iff \lambda = \limsup_{n\to \infty}|a_n|^\frac{1}{n} < +\infty.
	\]
	Then, the function $f:(x_0 - r,x_0+r)\to \R$ defined by
	\[
	f(x) = \sum_{n=0}^\infty a_n(x-x_0)^n
	\]
	is continuous on its domain.
\end{thm}
\begin{proof}
	\textcolor{blue}{Fix any $y\in (x_0-r,x_0+r)$}. We consider the following closed and bounded interval $I$ defined by
	\[
	(x_0-r,x_0+r)\supset I = \left\{
	\begin{array}{cl}
		[x_0-|y-x_0|,\,x_0+|y-x_0|], 	&\text{if }y\neq x_0 	\\
		\left[x_0 - \frac{\min(1,r)}{2},\,x_0 + \frac{\min(1,r)}{2}\right], 	&\text{if }y = x_0
	\end{array}
	\right..
	\]
	Hence, by~\Cref{thm:cvg_power}, the power series converges uniformly to $f$ on $I$. Since the partial sums $\sum_{j=0}^n a_j (x-x_0)^j$ are polynomials (which are continuous), we have by~\Cref{prop:unif_cvg_cts} that $f$ is continuous on $I$. In particular, $y\in I$ and so \textcolor{blue}{$f$ is continuous at $y$.}
\end{proof}
While working within the interval of convergence in the previous result, we extrapolated from the continuity of (finite degree) polynomials to deduce continuity of the infinite power series. Since polynomials are also differentiable, we would also like to deduce differentiability of the infinite power series.
\begin{cor}[Differentiability of convergent power series]
	\label{cor:diff_power}
	Under the same setting as~\Cref{thm:cty_power}, the function $f$ is differentiable on $(x_0 - r,x_0+r)$. Moreover, we have
	\[
	f'(x) = \sum_{n=0}^\infty (n+1)a_{n+1}(x-x_0)^n.
	\]
	The radius of convergence of this power series is also $r$.
\end{cor}
\begin{proof}
	\ul{Uniform convergence of $\sum_{n=0}^\infty (n+1)a_{n+1}(x-x_0)^n$:} Recall that $\lim_{n\to\infty} n^\frac{1}{n} =1$. Moreover, this allows to deduce (exercise) that
	\[
	\limsup_{n\to \infty}|na_{n+1}|^\frac{1}{n} = \limsup_{n\to\infty}|a_{n+1}|^\frac{1}{n}.
	\]
	Hence, the radius of convergence of $\sum_{n=1}^\infty na_n(x-x_0)^n$ is also $r$. Dividing by $(x-x_0)$, we see that $\sum_{n=0}^\infty (n+1)a_{n+1}(x-x_0)^n$ also has radius of convergence equal to $r$. Owing to~\Cref{thm:cvg_power}, we see that
	\begin{equation}
		\label{eq:diff_power}
		\sum_{n=0}^\infty (n+1)a_{n+1}(x-x_0)^n\text{ converges uniformly on any closed interval contained in }(x_0-r,x_0+r).
	\end{equation}
	\ul{Derivative of $f$ through partial sums:} Let us denote the following
	\[
	f_n(x):= \sum_{j=0}^n a_j(x-x_0)^j,\quad g(x) = \sum_{n=0}^\infty (n+1)a_{n+1}(x-x_0)^n.
	\]
	We will apply~\Cref{thm:unif_cvg_diff} (or~\Cref{thm:unif_cvg_series_diff} if you like). \textcolor{blue}{Fix any $y \in (x_0-r,x_0+r)$.} We consider the closed bounded interval
	\[
	(x_0-r,x_0+r)\supset I = \left\{
	\begin{array}{cl}
		[x_0-|y-x_0|,\,x_0+|y-x_0|], 	&\text{if }y\neq x_0 	\\
		\left[x_0 - \frac{\min(1,r)}{2},\,x_0 + \frac{\min(1,r)}{2}\right], 	&\text{if }y = x_0
	\end{array}
	\right..
	\]
	For every $n\in\N$, the function $f_n$ is a polynomial, hence~\ref{unif_cvg_diff_1} is satisfied. Certainly at $x=x_0$, the series $f_n(x_0) = a_0$ is convergent so~\ref{unif_cvg_diff_2} is true. Moreover, the previous paragraph (in particular~\eqref{eq:diff_power}) shows that $f_n'$ converges uniformly to $g$ on $I$ so~\ref{unif_cvg_diff_3} is true. The hypothesis of~\Cref{thm:unif_cvg_diff} is true and we conclude that $f$ is differentiable on $I$ and its derivative is given by
	\[
	f'(x) = g(x) = \sum_{n=0}^\infty (n+1)a_{n+1}(x-x_0)^n.
	\]
	In particular, \textcolor{blue}{$f$ is differentiable at $y\in I$. Since $y\in(x_0-r,x_0+r)$ was arbitrary, we see that $f$ is differentiable in the interval of convergence.}
\end{proof}
By successively differentiating, we can repeat~\Cref{cor:diff_power} over and over to deduce infinitely many derivatives.
\begin{cor}[Infinite differentiability of power series]
	\label{cor:analytic}
	Under the same setting as~\Cref{thm:cty_power}, the function $f$ has infinitely many derivatives. In other words, for every $k\in\N$, the $k$th order derivative $f^{(k)}$ exists and is given by
	\[
	f^{(k)}(x) = \sum_{n=k}^\infty \frac{n!}{(n-k)!}a_n(x-x_0)^{n-k},\quad \forall x\in (x_0-r,x_0+r).
	\]
	In particular, we have $f^{(k)}(x_0) = k!a_k$ for every $k\in\N\cup \{0\}$.
\end{cor}
In anticipation of~\Cref{sec:taylor}, another way of self-expressing the power series in~\Cref{cor:analytic} is
\[
f(x) = \sum_{n=0}^\infty \frac{f^{(n)}(x_0)}{n!}(x-x_0)^n,\quad \forall x\in(x_0-r,x_0+r).
\]
One consequence of~\Cref{cor:analytic} is that convergent power series are uniquely determined by the coefficients $(a_n)_{n\in\N\cup\{0\}}$.
\begin{cor}[Uniqueness of power series]
	\label{cor:uniq_analytic}
	Suppose that $\sum_{n=0}^\infty b_n (x-x_0)^n$ is a convergent power series. Let $\sum_{n=0}^\infty a_n(x-x_0)^n$ and $r>0$ be as in~\Cref{thm:cty_power} and assume that there exists some $\delta>0$ such that
	\[
	\sum_{n=0}^\infty a_n (x-x_0)^n = \sum_{n=0}^\infty b_n (x-x_0)^n,\quad \forall x \in (x_0-r,x_0+r)\cap (x_0-\delta,x_0+\delta).
	\]
	Then, $a_n = b_n$ for every $n\in\N\cup \{0\}$.
\end{cor}
\begin{proof}
	Let $f(x) = \sum_{n=0}^\infty a_n(x-x_0)^n$. Then, \Cref{cor:analytic} says $f^{(k)}(x_0) = k!a_k$ for every $k\in\N\cup\{0\}$. But by assumption, we also see that $f^{(k)}(x_0) = k!b_k$ for every $k\in\N\cup\{0\}$. This implies the result.
\end{proof}
\begin{rem}
	\Cref{cor:uniq_analytic} is the core of what is sometimes referred to as `analytic continuation' or `analytic extension.' In practice, it allows one to possibly expand the disk of convergence of $\sum_{n=0}^\infty b_n (x-x_0)^n$ when more information about $\sum_{n=0}^\infty a_n (x-x_0)^n$ is present. This~\href{https://www.youtube.com/watch?v=sD0NjbwqlYw&pp=ygUhYW5hbHl0aWMgY29udGludWF0aW9uIDNibHVlMWJyb3du}{link} is a famous pop science video summarising analytic continuation. If you're interested in learning more, you should take a course on Complex Analysis. One peculiarity of analytic continuation is the following weird looking equation
	\[
	\sum_{n=1}^\infty n \overset{!}{=} -\frac{1}{12}.
	\]
	Strictly speaking, the above is false so do not use it. The rough explanation is that one considers the \textit{Riemann zeta function} defined by
	\[
	R(\zeta) = \sum_{n=1}^\infty \frac{1}{n^\zeta}.
	\]
	We have already seen that $R(1)$ is the harmonic series which diverges. In the theory of complex analysis, one considers $R(\zeta)$ for $\zeta$ belonging to a subset of $\mathbb{C}$, the complex numbers. In principle, then, $R(\zeta)$ is only well-defined provided that the Re$(\zeta)>1$. However, a result from complex analysis is the following:
	\begin{enumerate}
		\item There exists a function $f(\zeta)$ which is infinitely differentiable on a subset of $\mathbb{C}$ containing $\zeta = -1$.
		\item Moreover, $f = R$ in the intersection of their domains.
	\end{enumerate}
	Analytic continuation then provides a meaning for $R(-1)$ by saying
	\[
	R(-1) := f(-1).
	\]
	The \textit{series representation} of $R(\zeta)$ does not make sense at $\zeta = -1$. However, we can still define the function at $\zeta = -1$ through $f$. Again, if you're interested in learning more, you should take a course on Complex Analysis (followed up by Analytic Number Theory).
\end{rem}
	\section{Taylor Series}
\label{sec:taylor}
We directly follow up from the previous section by `reversing' the starting point in the following way:
\begin{enumerate}
	\item \Cref{sec:power_series} started with convergent power series $\sum_{n=0}^\infty a_n(x-x_0)^n$. We considered these as functions $f(x) = \sum_{n=0}^\infty a_n(x-x_0)^n$ and then investigated the regularity properties (c.f.~\Cref{thm:cty_power},\Cref{cor:diff_power}, and~\Cref{cor:analytic}). In particular, we deduced that $f(x)$ has infinitely many derivatives.
	\item This section starts with a given function $f(x)$ which we assume has infinitely many derivatives. We then study the formal power series generated by $f(x)$ (c.f. below~\Cref{cor:analytic}) which should be given by
	\begin{equation}
		\label{eq:taylor_series}
		\sum_{n=0}^\infty \frac{f^{(n)}(c)}{n!}(x-c)^n.
	\end{equation}
\end{enumerate}
We already mentioned infinite differentiability, but let us be precise below.
\begin{defn}
	\label{defn:inf_diff_taylor}
	Let $I\subset \R$ and $f:I\to \R$. We say that $f$ is \textit{infinitely differentiable on $I$} provided
	\[
	\forall x\in I,\,\forall k\in \N,\, f^{(k)}(x)\text{ exists.}
	\]
	The set of all functions which are infinitely differentiable on $I$ is $C^\infty(I)$.
	
	Let $c\in\R$ and $\delta>0$ and suppose $f\in C^\infty(c-\delta,c+\delta)$. We define the \textit{Taylor series generated by $f$ around $c$} as the power series in~\eqref{eq:taylor_series}. When $c=0$, the Taylor series~\eqref{eq:taylor_series} is called the Maclaurin series.
\end{defn}
Clearly, the Taylor series~\eqref{eq:taylor_series} converges at $x=c$ and both $f(x)$ and its Taylor series agree at this point. However, a major distinction between the present discussion and previously in~\Cref{sec:power_series} is that~\eqref{eq:taylor_series} might not converge at other points $x\neq c$. Moreover, even if~\eqref{eq:taylor_series} converges at a point $x\neq c$, the Taylor series need not equal $f(x)$.
\begin{eg}
	Let $c=0$ and consider the function
	\[
	f(x) = \left\{
	\begin{array}{cl}
		e^{-\frac{1}{x}}, 	&\text{if }x>0 	\\
		0, 	&\text{if }x\le 0
	\end{array}
	\right..
	\]
	We have the following.
	\begin{enumerate}
		\item (Exercise) $f$ is infinitely differentiable on $\R$. In particular,
		\[
		f^{(n)}(0) = 0,\quad \forall n\in\N.
		\]
		\item However, the Maclaurin series (Taylor series with $c=0$) is identically
		\[
		\sum_{n=0}^\infty \frac{f^{(n)}(0)}{n!}x^n \equiv 0.
		\]
		Since $f(x) \neq 0$ for $x>0$, we see a discrepancy between the function and its Taylor series.
	\end{enumerate}
\end{eg}
In order to study how much a function and its Taylor series differ from each other, let us recall (a version of) Taylor's theorem.
\begin{thm}[Taylor's theorem]
	\label{thm:taylor}
	Let $a<c<b$ be three real numbers and $n\in\N$. Let $f$ be n times differentiable on $(a,b)$. Then, for every $x\in(a,b)\setminus\{c\}$, there exists some $y$ between $c$ and $x$ such that
	\[
	f(x) = \sum_{k=0}^{n-1}\frac{f^{(k)}(c)}{k!}(x-c)^k + \frac{f^{(n)}(y)}{n!}(x-c)^n.
	\]
\end{thm}
\begin{proof}
	\textcolor{blue}{Blank unless necessary.}
\end{proof}
Clearly, the sum on the right-hand side in~\Cref{thm:taylor} is the partial sum of the Taylor series. Therefore, a condition to ensure that a function is equal to its Taylor series is that the so-called \textit{remainder} term
\[
R_n(x) = \frac{f^{(n)}(y)}{n!}(x-c)^n
\]
vanishes as $n\to \infty$.
\begin{thm}
	\label{thm:taylor_eq}
	Let $f\in C^\infty((a,b))$ and fix $c\in (a,b)$. Assume that there exists $\delta>0$ and a constant $M\ge 0$ such that
	\[
	|f^{(n)}(x)|\le M^n,\quad \text{for every }x\in (c-\delta,c+\delta)\cap (a,b)\text{ and every }n\in\N.
	\]
	Then, we have
	\[
	f(x) = \sum_{n=0}^\infty \frac{f^{(n)}(c)}{n!}(x-c)^n,\quad \forall x\in (c-\delta,c+\delta)\cap (a,b).
	\]
\end{thm}
\begin{proof}
	Let $x\in (c-\delta,c+\delta)\cap (a,b)$. According to~\Cref{thm:taylor}, for any $n\in\N$, there exists some $y$ between $x$ and $c$ such that
	\[
	f(x) = \sum_{k=0}^{n-1}\frac{f^{(k)}(c)}{k!}(x-c)^k + R_n(x),\quad R_n(x) = \frac{f^{(n)}(y)}{n!}(x-c)^n.
	\]
	By assumption, we can estimate the remainder by
	\[
	\left|R_n(x)\right|\le \frac{(M\delta)^n}{n!},\quad\forall x\in (c-\delta,c+\delta)\cap (a,b).
	\]
	In particular, the remainder vanishes as $n\to \infty$ because (exercise)
	\[
	\lim_{n\to\infty}\frac{(M\delta)^n}{n!} = 0.
	\]
\end{proof}
	\section{The space of continuous functions}
\label{sec:C^0}
From now on, we will focus our attention to the space of continuous functions. To be specific, we will analyse a metric space whose elements are functions. To fix notation, given a metric space $(X,d)$, we denote the set
\[
C(X):=\{\text{continuous functions }f:X\to \R\}.
\]
We will endow $C(X)$ with a metric, but our standing assumption throughout the rest of these notes is to take $X$ to be \textbf{compact}. In applications, we will usually be explicit about $X$ as a compact subset of $\R$ or $\R^d$, but we keep it general for now. In this setting, we define the \textit{sup metric}
\[
d_\infty(f,g):= \sup_{x\in X}|f(x) - g(x)|,\quad f,g\in C(X).
\]
Based on notational preference, we may also refer to this with the \textit{sup norm} which is any of
\[
\|f-g\|,\,\|f-g\|_{\infty},\,\text{or }\|f-g\|_{C(X)},
\]
all of the above referring to $d_\infty(f,g)$. Let us first confirm that we are working with a valid metric space.
\begin{prop}[$C(X)$ is a metric space]
	Let $X$ be a compact metric space. The pair $(C(X),\,d_\infty)$ as defined above is a metric space.
\end{prop}
\begin{proof}
	We need to check that $d_\infty$ satisfies the criteria in~\Cref{defn:metric} of being a metric. First, let us check that it is well-defined as a map from $C(X)\times C(X)$ with values in $[0,+\infty)$. For any $f,g \in C(x)$, the Extreme Value Theorem from~\Cref{thm:EVT} yields that $\sup_{x\in X}|f(x) - g(x)|$ is a real and non-negative number. Hence, $d_\infty :C(X)\times C(X)\to [0,+\infty)$ is well-defined.
	
	By definition, we see that $d_\infty(f,g)\ge 0$ for every $f,g\in C(X)$. We have
	\[
	d_\infty(f,g) = 0 \iff \sup_{x\in X}|f(x) - g(x)| = 0 \iff f(x) = g(x),\, \forall x \in X \iff f = g.
	\]
	This verifies~\ref{m:pos}. The symmetry condition in~\ref{m:sym} is also immediate because
	\[
	d_\infty(f,g) = \sup_{x\in X}|f(x) - g(x)| = \sup_{x\in X}|g(x) - f(x)| = d_\infty(g,f).
	\]
	As for the triangle inequality in~\ref{m:tri}, for any fixed $f,g,h\in C(X)$ and any $x\in X$, the usual triangle inequality for real numbers gives
	\[
	|f(x) - g(x)| \le |f(x) - h(x)| + |h(x) - g(x)|.
	\]
	On the right-hand side, we have $|f(x) - h(x)| \le \sup_{x\in X}|f(x) - h(x)|$ and $|h(x) - g(x)| \le \sup_{x\in X}|h(x) - g(x)|$ so we get
	\[
	|f(x) - g(x)| \le d_\infty(f,h) + d_\infty (h,g).
	\]
	Since this inequality is true for all $x\in X$, in particular, we obtain
	\[
	d_\infty(f,g) \le d_\infty(f,h) + d_\infty(h,g).
	\]
	We will prove that $C(X)$ enjoys more metric properties. For now, let us comment on the crucial connection with uniform convergence. Suppose $(f_n)_{n\in\N}\subset C(X)$ is a sequence converging to $f \in C(X)$ with respect to the metric $d_\infty$. Let's unpack this definition. According to~\Cref{defn:cvg_metric}, this means that
	\[
	\forall \varepsilon>0,\exists N\in\N\text{ s.t. }n\ge N \implies d_\infty(f_n,f)<\varepsilon.
	\]
	Based on how we've defined $d_\infty$, we have
	\begin{align*}
		d_\infty(f_n,f)<\varepsilon &\iff \sup_{x\in X}|f_n(x) - f(x)|<\varepsilon 	\\
		&\iff |f_n(x) - f(x)|<\varepsilon,\,\forall x\in X.
	\end{align*}
	So convergence in $C(X)$ is the same as uniform convergence in~\Cref{defn:unif_cvg} (c.f.~\Cref{rem:sup_unif_cvg} as well)!
\end{proof}
We will conclude this introductory section on $C(X)$ by showing that it is complete.
\begin{prop}[$C(X)$ is complete]
	\label{prop:C^0_complete}
	Let $X$ be a compact metric space. The metric space $(C(X),d_\infty)$ as defined above is complete.
\end{prop}
\begin{proof}
	\textcolor{blue}{Fix $(f_n)_{n\in\N}\subset C(X)$ to be a Cauchy sequence.} Our goal is to construct a limit in $C(X)$ for this sequence. For every $x \in X$, we see that the sequence of real numbers $(f_n(x))_{n\in\N}$ is Cauchy. Since $\R$ is complete, $(f_n(x))_{n\in\N}$ is a convergent sequence for every $x\in X$. Since the limit depends on $x\in X$, \textcolor{blue}{we construct $f:X\to \R$} as the pointwise limit by
	\[
	f(x) := \lim_{n\to \infty}f_n(x),\quad \forall x \in X.
	\]
	\ul{Uniform convergence:} Fix any arbitrary $x\in X$. Owing to the construction above, we see that
	\begin{align*}
	|f_n(x) - f(x)| &= \lim_{m\to \infty}|f_n(x) - f_m(x)| \le \limsup_{m\to \infty} |f_n(x) - f_m(x)| \\
	&\le\limsup_{m\to \infty} \left(\sup_{x\in X}|f_n(x) - f_m(x)|\right) = \limsup_{m\to \infty}\|f_n-f_m\|_\infty.
	\end{align*}
	Since $x\in X$ was arbitrary, we obtain
	\[
	\|f_n-f\|_\infty \le \limsup_{m\to \infty}\|f_n-f_m\|_\infty.
	\]
	By assumption, $(f_n)_{n\in\N}$ is a Cauchy sequence, hence for any $\varepsilon>0$, there exists $N\in \N$ such that $n\ge N$ implies (why?)
	\[
	\limsup_{m\to\infty}\|f_n-f_m\|_\infty < \varepsilon.
	\]
	This shows (why? Students fill in details) that $f_n$ converges uniformly to $f$ on $X$. According to~\Cref{prop:unif_cvg_cts}, we obtain \textcolor{blue}{$f \in C(X)$.} In particular, from the discussion before this result, we conclude \textcolor{blue}{$f_n$ converges to $f$ in $C(X)$ with respect to $d_\infty$.}
\end{proof}
We also demonstrate another proof by looking at
\[
\mathcal{B}(X) = \left\{f:X\to \R \, | \, \sup_{x\in X}|f(x)| < +\infty\right\}.
\]
\begin{ex}
The pair $(\mathcal{B}(X),d_\infty)$ where
\[
d_\infty(f,g) := \sup_{x\in X}|f(x)-g(x)|,\quad f,g\in \mathcal{B}(X)
\]
is a metric space.
\end{ex}
More interesting is the following result.
\begin{prop}
\label{prop:B_comp}
$(\mathcal{B}(X),d_\infty)$ is a complete metric space.
\end{prop}
\begin{proof}
Fix $(f_n)_{n\in\N}\subset \mathcal{B}(X)$, a Cauchy sequence of bounded functions. We aim to demonstrate that this sequence converges and we proceed in a few steps.
\begin{enumerate}
\item \ul{Candidate for the limit:} For every $x\in X$, the sequence (of real numbers) $(f_n(x))_{n\in\N}$ is Cauchy just like in~\Cref{prop:C^0_complete}. Since $\R$ is complete, we see that each sequence $(f_n(x))_{n\in\N}$ converges to some real number which we label $f(x)$. Equivalently, this defines the pointwise limit $f:X\to\R$ where
\[
f(x) = \lim_{n\to\infty} f_n(x).
\]
This is our candidate for the limit and the remaining steps of this proof will be dedicated to showing that it is also a bounded function and that it is the uniform limit of $(f_n)_{n\in\N}$.
\item \ul{$f \in \mathcal{B}(X)$:} Since $(f_n)_{n\in\N}$, for the choice of $\varepsilon = 1>0$, there exists some $N\in\N$ such that $\sup_{x\in X}|f_n(x) - f_m(x)| < 1$ for every $n,m\ge N$. In particular, we see that
\[
|f_N(x) - f_m(x)|\le 1,\quad \forall x\in X,\, \forall m \ge N.
\]
Since $f(x) = \lim_{m\to \infty}f_m(x)$, we pass to the limit $m\to\infty$ in the above inequality to obtain
\[
|f_N(x) - f(x)|\le 1,\quad \forall x\in X.
\]
Since $f_N\in \mathcal{B}(X)$, there exists some $M>0$ such that
\[
|f_N(x)|\le M,\quad \forall x\in X.
\]
Putting the previous two points together with the triangle inequality gives
\[
|f(x)| \le |f(x) - f_N(x)| + |f_N(x)| = M+1,\quad \forall x\in X.
\]
This shows that $f\in \mathcal{B}(X)$.
\item \ul{$f_n\to f$ uniformly:} Fix $\varepsilon>0$. Since $(f_n)_{n\in\N}$ is Cauchy, there exists $N\in \N$ such that
\[
n\ge N \implies \sup_{x\in X}|f_n(x) - f_N(x)|<\frac{\varepsilon}{2}.
\]
Being that $f$ is the pointwise limit of $(f_n)_{n\in\N}$, for any $x\in X$, we see that
\[
|f_N(x) - f(x)| = \lim_{m\to\infty}|f_N(x) - f_m(x)| < \frac{\varepsilon}{2},\quad \forall x \in X.
\]
Putting the previous two parts together with the triangle inequality, we obtain
\[
|f_n(x) - f(x)| \le |f_n(x) - f_N(x)| + |f_N(x) - f(x)| < \varepsilon,\quad \forall x\in X.
\]
This demonstrates the uniform convergence $f_n \to f$.
\end{enumerate}
\end{proof}
With~\Cref{prop:B_comp} in hand, we can present a different proof of~\Cref{prop:C^0_complete}.
\begin{proof}[Alternative proof of~\Cref{prop:C^0_complete}]
We aim to show that $C(X)$ is a closed subspace of $\mathcal{B}(X)$. Then, \Cref{prop:comp_clos} will allow us to conclude that $C(X)$ is complete. 

Let $(f_n)_{n\in\N} \subset C(X)$ be any sequence such that it converges to $f\in \mathcal{B}(X)$ (valid when $X$ is compact by the Extreme Value Theorem). Since convergence with respect to $d_\infty$ is the same as uniform convergence, we have that $(f_n)_{n\in\N}$ uniformly converges to $f$. Since each $f_n$ is continuous, we can apply~\Cref{prop:unif_cvg_cts} to deduce that $f$ is continuous i.e. $f \in C(X)$. Thus, by~\Cref{lem:closed_lim}, we see that $C(X)$ is closed.
\end{proof}
%	\subfile{C^0}
	\section{ODE application: Picard's theorem}
\label{sec:Picard}
Using what we have learned about complete metric spaces, we can now present an important basic application for solving differential equations. Fix $T>0, x_0\in\R$, and let $f:[0,T]\times\R\to\R$ be a function. We would like to ascertain when we can solve the following differential equation
\begin{equation}
	\label{eq:ODE}
	\frac{d}{dt}x(t) = f(t,x(t)),\quad t\in (0,T),\quad x(0) = x_0.
\end{equation}
We have deliberately written~\eqref{eq:ODE} with the unknown being $x$ as a function of $t$ to evoke the following understanding from physics: Typically, $x(t)$ represents a particle's position (in space) at time $t$. Physics, in particular classical mechanics, is able to provide many rules relating a particle's velocity (or even acceleration) with forces exerted on it. To be precise, the time derivative $\frac{dx}{dt}$ represents the velocity of the particle and classical mechanics can describe how this relates to various effects experienced by the particle. The question that physics seeks to answer is the following: given full knowledge on the forces that a particle experiences (i.e. given a differential equation relating velocity with forces), can we deduce what the position of the particle is? In abstract mathematical formality, this motivates trying to solve~\eqref{eq:ODE}

There are many points to clarify. When we say `solve~\eqref{eq:ODE}' we mean to answer several questions, the first of which is given below.
\begin{q}[Existence]
	Does there exist a differentiable function $x:[0,T]\to \R$ such that
	\[
	\frac{d}{dt}x(t) = f(t,x(t)),\quad \forall t\in(0,T)
	\]
	and $x(0) = x_0$?
\end{q}
The equation $\frac{d}{dt}x(t) = f(t,x(t))$ is referred to as the `differential equation' whereas the pointwise evaluation $x(0) = x_0$ is the `initial condition.'
\begin{eg}
	Let $x_0 = 0$ and $f(t,x) = t$. The function $x(t) = \frac{1}{2}t^2$ satisfies
	\[
	\frac{d}{dt}x(t) = t\quad\text{and}\quad x(0) = 0.
	\]
\end{eg}
\begin{eg}
	Let $x_0 = 1$ and $f(t,x) = x$. Then the function $x(t) = e^t$ satisfies
	\[
	\frac{d}{dt}x(t) = x(t)\quad\text{and}\quad x(0) = 1.
	\]
\end{eg}
\begin{eg}
	Let $x_0 = 1$ and $f(t,x) = -xt$. Then the function $x(t) = e^{-\frac{1}{2}t^2}$ satisfies
	\[
	\frac{d}{dt}x(t) = -tx(t)\quad\text{and}\quad x(0) = 1.
	\]
\end{eg}
In your previous courses in differential equations, you learned many techniques to explicitly write down a formula for the solution of~\eqref{eq:ODE}. The purpose of this section is to illustrate that solving~\eqref{eq:ODE} does not necessarily mean finding equations.
\begin{eg}
	Does a differentiable function $x(t)$ exist such that its derivative is given by
	\[
	\frac{d}{dt}x(t) = \sin(\cos(\cos t))
	\]
	and it satisfies the initial condition $x(0) = 5$? An application of~\Cref{thm:Picard} will answer this example. For now, can you summon up the tools and tricks you've learned to solve this example?
\end{eg}
Another aspect which may have been ignored in your differential equations classes is whether the techniques you learned capture \textit{all} possible solutions or if they miss out on \textit{some} solutions. For example, how do you know that the method of separation of varriables / variation of parameters / undetermined coefficients / etc\dots applied to differential equations will yield the `correct' solution? This is summarised in the following question.
\begin{q}[Uniqueness]
	Suppose $x_1,x_2:[0,T]\to \R$ are differentiable functions which both satisfy
	\[
	\frac{d}{dt}x(t) = f(t,x(t)),\quad \forall t\in(0,T)\quad\text{and}\quad x(0) = x_0,
	\]
	with $x(t) = x_1(t)$ or $x(t) = x_2(t)$. Can we conclude $x_1(t) = x_2(t)$ for every $t\in[0,T]$? Is it possible that $x_1\neq x_2$?
\end{q}
We will be able to answer both existence and uniqueness for~\eqref{eq:ODE} by reformulating the problem and making precise assumptions on $f$ below in~\Cref{thm:Picard}. The first assumption (which we will eventually strengthen) is that $f$ is a \textbf{continuous} function. 
\begin{prop}[Integral form of~\eqref{eq:ODE}]
	\label{prop:ODE_int}
	Let $f:[0,T]\times\R\to\R$ be continuous. TFAE:
	\begin{itemize}
		\item $x:[0,T]\to \R$ is a continuously differentiable solution to~\eqref{eq:ODE}.
		\item \[
		x(t) = x_0 + \int_0^t f(s,x(s))\,\mathrm{d}s.
		\]
	\end{itemize}
\end{prop}
\begin{proof}
	This equivalence is immediate from integrating~\eqref{eq:ODE} and the Fundamental Theorem of Calculus (exercise).
\end{proof}
The simple equivalence in~\Cref{prop:ODE_int} gives us an iterated formula for solutions to~\eqref{eq:ODE}. With this in mind, let us think more abstractly in terms of metric spaces, specifically with $C([0,T])$. We define the map
\[
I:C([0,T])\to C([0,T])
\]
by the following way: Given any $x\in C([0,T])$, $I[x]$ belongs to $C([0,T])$ and it is defined as
\[
I[x](t):= x_0 + \int_0^tf(s,x(s))\,\mathrm{d}s.
\]
Here, we are using square brackets $I[x]$ to mean that $I$ takes the function $x$ as an input and the output is the function $I[x](t)$.
\begin{ex}
	Let $f:[0,T]\times \R\to\R$ be continuous. Verify that $I:C([0,T])\to C([0,T])$ described above is well-defined.
\end{ex}
With this setting, another way of describing solutions to~\eqref{eq:ODE} according to~\Cref{prop:ODE_int} is by searching for \textbf{fixed points} of the map $I$. Indeed, suppose $x\in C([0,T])$ is a fixed point of $I$. This gives
\[
I[x](t) = x(t),\quad\forall t\in[0,T]\overset{\Cref{prop:ODE_int}}{\iff} x\text{ is a solution to~\eqref{eq:ODE}.}
\]
Our strategy for analysing existence and uniqueness of solutions to~\eqref{eq:ODE} consists now of finding fixed points to the map $I$. Here, we will rely on~\Cref{thm:banach}.
\begin{thm}[Picard-Lindel\"of / Cauchy-Lipschitz]
	\label{thm:Picard}
	Fix $T>0, x_0\in\R$ and a continuous function $f:[0,T]\times\R\to \R$. Assume that $f$ satisfies the following uniform Lipschitz property:
	\[
	\exists L> 0\text{ s.t. }\forall t\in[0,T],\,\forall x,y\in \R,\text{ we have }|f(t,x)-f(t,y)|\le L|x-y|.
	\]
	Then, there exists a unique continuously differentiable function $x:[0,T]\to \R$ such that~\eqref{eq:ODE} is satisfied.
\end{thm}
\begin{proof}
	\ul{Short-time solution:} From~\Cref{prop:C^0_complete}, we see that \textcolor{blue}{$C([0,\tau])$ is a complete metric space for any $\tau>0$.} We will eventually explicitly choose $0 < \tau \ll T$. Specifically, $\tau$ will be chosen to make $I:C([0,\tau])\to C([0,\tau])$ a contraction. Fix any $x_1,x_2\in C([0,\tau])$ and we calculate
	\[
		I[x_1](t) - I[x_2](t) = \int_0^t f(s,x_1(s)) - f(s,x_2(s))\,\mathrm{d}s,\quad \forall t\in[0,\tau].
	\]
	By the triangle inequality and the uniform Lipschitz assumption on $f$, for any $t\in [0,\tau]$, we have that
	\begin{align*}
		|I[x_1](t) - I[x_2](t)| &\le \int_0^t |f(s,x_1(s)) - f(s,x_2(s))|\,\mathrm{d}s 	\\
		&\le L\int_0^t |x_1(s) - x_2(s)|\,\mathrm{d}s 	\\
		&\le L\int_0^t \left(\sup_{s\in[0,\tau]}|x_1(s)-x_2(s)|\right)\,\mathrm{d}s 	\\
		&\le Lt \|x_1-x_2\|_{C([0,\tau])}.
	\end{align*}
	Since this inequality is true for any arbitrary $t\in[0,\tau]$, we obtain
	\[
	\|I[x_1]-I[x_2]\|_{C([0,\tau])}\le L \tau \|x_1-x_2\|_{C([0,\tau])}.
	\]
	In particular, whenever $\tau < \frac{1}{L}$ (for example, set $\tau = \frac{1}{2L+1}$), this estimate shows \textcolor{blue}{$I:C([0,\tau])\to C([0,\tau])$ is a contraction.} The text in \textcolor{blue}{blue} shows that we are in a position to apply~\Cref{thm:banach}: There exists a unique $x\in C([0,\tau])$ such that
	\[
	I[x](t) = x(t),\quad \forall t\in[0,\tau].
	\]
	Owing to the discussion before this result, we see that this function $x$ is a solution to the ODE \textit{but only for times} $0 < t < \tau$. We next seek to extend the time horizon to $T$.
	
	\ul{Extension to global time:} We repeat the previous argument for $[\tau,2\tau]$. Let
	\[
	S:C([\tau,2\tau])\to C([\tau,2\tau])
	\]
	be defined by
	\[
	S[y](t):= x(\tau) + \int_\tau^{2\tau}f(s,y(s))\,\mathrm{d}s,\quad t\in[\tau,2\tau].
	\]
	By exactly the same arguments as before, we see that there exists a unique $y\in C([\tau,2\tau])$ such that
	\[
	\frac{d}{dt}y(t) = f(t,y(t)),\quad \forall t\in[\tau,2\tau],\quad y(\tau) = x(\tau).
	\]
	We extend $x\in C([0,\tau])$ from the previous paragraph by defining
	\[
	x(t) := y(t),\quad t\in[\tau,2\tau].
	\]
	This can be repeated on intervals of the form $[n\tau,(n+1)\tau]$ for $n\in\N$ until $x$ is extended to $[0,T]$. Finally, we discuss uniqueness.
	
	\ul{Uniqueness:} Suppose $x_1,x_2:[0,T]\to\R$ are continuously differentiable functions solving~\eqref{eq:ODE}. Using the uniform Lipschitz property of $f$, we have
	\[
	\frac{d}{dt}\frac{1}{2}|x_1(t)-x_2(t)|^2 = (x_1(t)-x_2(t))(f(t,x_1(t))-f(t,x_2(t)))\le L|x_1(t)-x_2(t)|^2,\quad \forall t\in[0,T].
	\]
	For simplicity, let us set $y(t) = \frac{1}{2}|x_1(t)-x_2(t)|^2$ so that the differential inequality above reads
	\[
	\frac{d}{dt}y(t) \le 2Ly(t),\quad \forall t\in[0,T].
	\]
	By introducing the integrating factor of $e^{-2Lt}$, we see that
	\[
	\frac{d}{dt}\left( e^{-2Lt}y(t) \right)\le 0 \implies e^{-2Lt}y(t) \le y(0)\implies y(t) \le y(0) e^{2Lt},\quad \forall t\in[0,T].
	\]
	However, since $x_1,x_2$ both satisfy the same initial condition, we have $y(0) = \frac{1}{2}|x_1(0)-x_2(0)|^2 = 0$ and so we deduce (by definition, $y(t) \ge 0$ for all $t$)
	\[
	y(t) = 0,\quad \forall t\in[0,T].
	\]
	In terms of $x_1$ and $x_2$, this shows that $x_1 = x_2$.
\end{proof}
\begin{rem}
	The technique for proving uniqueness above is a specific demonstration of what is generally known as applying \textit{Gr\"onwall's inequality / lemma.}
\end{rem}
There are many generalizations of~\Cref{thm:Picard}, some of which are described in a problem sheet. The main crucial point is the Lipschitz property of $f$ with respect to its second argument.
\begin{eg}[An application]
	Let $x_0 = 1$ and $f(t,x) = t\cos(\cos x)$. Can you write down the formula for $x(t)$ where $x(t)$ satisfies
	\[
	\frac{d}{dt}x(t) = f(t,x(t)),\quad \forall t\in(0,T),\quad x(0) = 1?
	\]
	I certainly don't know enough tricks to explicitly solve this differential equation. Nevertheless, I know because of~\Cref{thm:Picard} that this differential equation admits a unique solution.
	
	\ul{Lipschitz:} Fix any $x,y\in \R$ and $t\in[0,T]$ and consider
	\[
	f(t,x) - f(t,y) = t \left(\cos(\cos x) - \cos(\cos y)\right).
	\]
	Owing to the Mean Value Theorem, there is some $\xi$ between $x$ and $y$ such that
	\[
	f(t,x) - f(t,y) = t\left[\sin(\cos \xi)\right]\left[\sin \xi\right] (x-y).
	\]
	We know that $|\sin z| \le 1$ for every $z\in\R$ so the above equation implies
	\[
	|f(t,x) - f(t,y)| \le T|x-y|,\quad \forall x,y\in \R,\,\forall t\in[0,T].
	\]
	Hence, we see that $f$ satisfies the desired criteria and so~\Cref{thm:Picard} applies.
\end{eg}
\begin{eg}[Lipschitz is essential]
	Let us consider $x_0 = 0$ and $f(t,x) = x^\frac{1}{2}$. This is slightly improper compared to~\Cref{thm:Picard} because $f(t,\cdot)$ is only defined on $[0,+\infty)$. However, there is a version of~\Cref{thm:Picard} which allows to consider the domain of $f(t,\cdot)$ as a strict subset of $\R$. Consider the following differential equation
	\begin{equation}
		\label{eq:ODE_root}
		\frac{d}{dt}x(t) = x(t)^\frac{1}{2},\quad x(0) = 0.
	\end{equation}
	We will demonstrate \textit{infinitely} many solutions to~\eqref{eq:ODE_root}. First and foremost is the trivial solution $x(t) \equiv 0$. The next obvious one (using familiar techniques) is $x(t) = \frac{1}{4}t^2$. These two solutions can be combined to create infinitely many more in the following way: Let $\tau>0$ be any positive number, define
	\[
	x_\tau(t):=\left\{
	\begin{array}{cc}
		0, 	&\text{if }t\le \tau 	\\
		\frac{1}{4}(t - \tau)^2, 	&\text{if }t > \tau
	\end{array}
	\right..
	\]
	The function $x_\tau$ satisfies~\eqref{eq:ODE_root} for every $\tau>0$.
\end{eg}
\end{document}